% =====================================================================
% Detailed Notes on a Minimal Symmetry-Allowed Model for TmFeO_3
% in the Gamma_2 Phase --- Rewritten edition (May 2026)
% =====================================================================
% Rewrite history:
%   v1 (original): see tmfeo3_notes_v1_original.tex (253 kB, 3705 lines)
%   v2 (this file): full rewrite addressing all items in
%        audit_tmfeo3_notes.tex.  Notable fixes vs. v1:
%        - corrected qFM and qAFM Herrmann formulas with explicit DM gap;
%        - corrected Bertaut effective-field normalization (factor of 4);
%        - explicit T-symmetric gauge for non-Kramers projected dipoles;
%        - explicit m_J->-m_J parity assignment for the three CF singlets;
%        - explicit 4x4 linearized magnon matrix M^{Gamma_2} written out;
%        - Lindblad/CPTP form of Tm dissipator;
%        - SI emission formula with explicit Z_0/c prefactor on M-dipole;
%        - explicit Liouville-pathway derivation of every 2DCS peak;
%        - bond-by-bond derivation of the DM Bertaut prefactor;
%        - separation of microscopic per-site H from macroscopic Bertaut H;
%        - corrected spin-reorientation phenomenology (one continuous SR
%          bracketed by two second-order transitions T_SR1, T_SR2).
% =====================================================================

\documentclass[reprint,amsmath,amssymb,aps,prb,longbibliography]{revtex4-2}

\usepackage{graphicx}
\usepackage{bm}
\usepackage{hyperref}
\usepackage{xcolor}
\usepackage{enumitem}
\usepackage{booktabs}
\usepackage{mathtools}
\usepackage{dsfont}

\hypersetup{colorlinks=true,linkcolor=blue,citecolor=blue,urlcolor=blue}

\setlength{\emergencystretch}{3em}

\newcommand{\note}[1]{#1}
\newcommand{\fix}[1]{#1}

% Operator shortcuts
\newcommand{\Jeff}{\mathbf{J}^{\rm eff}}
\newcommand{\deff}{\mathbf{d}^{\rm eff}}
\newcommand{\bfS}{\mathbf{S}}
\newcommand{\bfH}{\mathbf{H}}
\newcommand{\bfF}{\mathbf{F}}
\newcommand{\bfG}{\mathbf{G}}
\newcommand{\bfA}{\mathbf{A}}
\newcommand{\bfC}{\mathbf{C}}
\newcommand{\bfX}{\mathbf{X}}
\newcommand{\bfY}{\mathbf{Y}}
\newcommand{\bfZ}{\mathbf{Z}}
\newcommand{\bfP}{\mathbf{P}}
\newcommand{\bfM}{\mathbf{M}}
\newcommand{\bfE}{\mathbf{E}}
\newcommand{\bfD}{\mathbf{D}}
\newcommand{\bfr}{\mathbf{r}}
\newcommand{\bfa}{\mathbf{a}}
\newcommand{\bfb}{\mathbf{b}}
\newcommand{\bfc}{\mathbf{c}}
\newcommand{\bfw}{\mathbf{w}}
\newcommand{\Tr}{\mathrm{Tr}}
\newcommand{\half}{\tfrac{1}{2}}
\newcommand{\quart}{\tfrac{1}{4}}
\DeclareMathOperator{\diag}{diag}

\begin{document}

\title{A Symmetry-Allowed Minimal Model for Two-Dimensional Coherent Terahertz Spectroscopy of TmFeO\texorpdfstring{$_3$}{3} in the \texorpdfstring{$\Gamma_2$}{Gamma2} Phase}

\author{Author list (placeholder)}
\affiliation{Affiliations (placeholder)}
\date{\today}

\begin{abstract}
We construct a symmetry-allowed minimal Hamiltonian for the rare-earth
orthoferrite TmFeO$_3$ in its low-temperature $\Gamma_2$ phase and use it to
derive the third-order nonlinear response that produces the observed
two-dimensional coherent terahertz spectroscopy (2DCS) peak structure,
for the
three pump--detect geometries
$(H\!\parallel\! a, E\!\parallel\! c,\mathrm{VV})$,
$(H\!\parallel\! a, E\!\parallel\! c,\mathrm{VH})$, and
$(H\!\parallel\! c, E\!\parallel\! a,\mathrm{VH})$. The model couples a
four-sublattice Fe Landau--Lifshitz--Gilbert sector (with Heisenberg,
Dzyaloshinskii--Moriya, and biaxial single-ion anisotropy) to a Tm$^{3+}$
non-Kramers crystal-field qutrit
$\{|E_1(A_1)\rangle,|E_2(A_1)\rangle,|E_3(A_2)\rangle\}$ projected onto the
SU(3) Gell--Mann basis. We derive Bertaut linear and quadratic
spin-reorientation (SR) modes carefully, including the
Dzyaloshinskii--Moriya gap on the qAFM mode (which a naive
exchange-only treatment misses by orders of magnitude), and we cast the
Tm dissipation in CPTP-compliant Lindblad form. Each 2DCS peak is
identified with a specific Liouville pathway in $\chi^{(3)}$, with explicit
double-sided Feynman labelling, rephasing/non-rephasing assignment, and a
formula for the relative amplitude in terms of the projected matrix
elements. The derivation closes with an explicit $4\times 4$ linearized
magnon matrix $M^{\Gamma_2}$ and a sympy/numpy verification script that
reproduces every algebraic claim.
\end{abstract}

\maketitle

\tableofcontents

% =====================================================================
\section{Introduction}
\label{sec:intro}

Rare-earth orthoferrites $R$FeO$_3$ are a paradigmatic platform for
coherent ultrafast control of antiferromagnetic order. Resonant terahertz
excitation of the canted G-type Fe$^{3+}$ sublattice~\cite{Kimel2004,Kimel2005}
launches the high-frequency quasi-antiferromagnetic (qAFM) and the
low-frequency quasi-ferromagnetic (qFM) magnon modes, while resonant
excitation of the rare-earth crystal-field (CF) levels modifies the
effective Fe anisotropy and drives a coherent spin
reorientation~\cite{Baierl2016,Schlauderer2019,Mashkovich2021}. TmFeO$_3$
is a benchmark case because Tm$^{3+}$ is a non-Kramers ion ($J=6$, $^3H_6$),
its low-energy spectrum consists of three well-separated singlets within
the bottom 8\,meV
($\epsilon_2-\epsilon_1=1.92$\,meV, $\epsilon_3-\epsilon_1=7.84$\,meV)
that are addressable by single-cycle THz fields, and the compound exhibits
a continuous Fe spin reorientation in the $ac$ plane bounded by
\emph{two} second-order phase transitions, $T_{\rm SR1}\!\approx\!81$\,K and
$T_{\rm SR2}\!\approx\!91$\,K, between the high-temperature $\Gamma_4$ and
the low-temperature $\Gamma_2$ phases, with the mixed
$\Gamma_{24}$ rotation between them~\cite{LeendersHerrmann1971,
ShapiroAxeBirgeneau1974,Belov1976,Yamaguchi1974,Schlauderer2019}.

Two-dimensional coherent terahertz spectroscopy (2DCS,
2DTS)~\cite{LuJonas2017,WanArmitage2019,MahmoodWan2021,LiNorrisHuber2022,
HuberPRX2024} resolves the third-order nonlinear susceptibility
$\chi^{(3)}(t_1,t_2,t_3)$ of magnetic materials into individual Liouville
pathways. In TmFeO$_3$ this is the natural language for asking which
combinations of the four Fe Bertaut modes ($\bfF,\bfG,\bfC,\bfA$) and the
three Tm CF coherences ($\omega_{12},\omega_{13},\omega_{23}$) couple,
and through which microscopic vertices.

\subsection{What this paper provides}

\begin{enumerate}[leftmargin=*]
\item A symmetry-derived minimal Hamiltonian, written in a single,
internally-consistent normalization that we audit end to end. The
Hamiltonian derivation is collected first, through the uniform
$\mathbf q=0$ working model of Sec.~\ref{sec:Htotal}.
\item A non-Kramers Tm projected-dipole sector with explicit
$T$-symmetric gauge for the three singlets (Sec.~\ref{sec:CEF}), so that
all projected operators $\Jeff_a$ are purely imaginary and the
$m_J\!\to\!-m_J$ relative parities of $|E_{1,2}\rangle$ and $|E_3\rangle$
are stated.
\item A dedicated 2DCS framework section (Sec.~\ref{sec:2dcsframework})
that introduces the nonlinear field, response-function language, and the
three experimental geometries before any dynamical calculation.
\item A later dynamical analysis in which the reduced Fe magnon sector,
the Tm Bloch equations, the joint EOM, and the detailed Liouville
pathways are developed only after the Hamiltonian and 2DCS setup are in
place (Secs.~\ref{sec:bloch}--\ref{sec:EOMexpansion}).
\item A Tm--Fe coupling sector (Sec.~\ref{sec:TmFe}) with bond-resolved
linear exchange and on-site Van Vleck back-action, classified by the
$Pbnm$ point group as $A_1^\pm$ and $A_2^\pm$, with the on-site
inversion-pair cancellation analyzed explicitly.
\item A CPTP-compliant Tm dissipator (Sec.~\ref{sec:bloch}) and a
Liouville-pathway derivation of every peak in the six-peak target list
(Sec.~\ref{sec:2dcs}).
\item A coherent-emission expression with the SI impedance prefactor
$Z_0/c$ on the magnetic-dipole channel (Sec.~\ref{sec:emission}).
\item A bond enumeration for the Fe--Fe and Fe--Tm shells
(App.~\ref{app:bonds}) and a glossary of normalization conventions
(App.~\ref{app:conventions}) so that every prefactor in every formula can
be traced back to a single source.
\end{enumerate}

The organizational logic of the manuscript is therefore: first derive the
Hamiltonian only (Secs.~\ref{sec:crystal}--\ref{sec:Htotal}); then
introduce the 2DCS setup and experimental geometries
(Sec.~\ref{sec:2dcsframework}); only after that do we develop the reduced
dynamical equations, coherent emission, Liouville pathways, and
perturbative EOM treatment (Secs.~\ref{sec:bloch}--\ref{sec:EOMexpansion}).

\subsection{Notation}
We use $\hbar = 1$ in all dynamical equations (restored in
final numerical evaluations). $\mu_B$ is the Bohr magneton, $g_J=7/6$ the
Tm$^{3+}$ Land\'e factor in $^3H_6$, and $\gamma_e = g_e\mu_B/\hbar$ the
Fe gyromagnetic ratio (with $g_e\approx 2$). The eight Gell-Mann matrices
satisfy $\Tr(\lambda^a\lambda^b)=2\delta^{ab}$, $[\lambda^a,\lambda^b]=
2if^{abc}\lambda^c$, $\{\lambda^a,\lambda^b\}=\frac{4}{3}\delta^{ab}+
2d^{abc}\lambda^c$. Cartesian indices are $a,b,c\in\{x,y,z\}$ for spin
and $\mu,\nu\in\{1,2,3,4\}$ for the four Fe (or Tm) sublattices. The
Bertaut basis variables are $\bfF$ (ferromagnetic), $\bfG$ (G-type AFM),
$\bfC$ (C-type AFM), $\bfA$ (A-type AFM).

% =====================================================================
\section{Crystal structure and sublattice frames}
\label{sec:crystal}

\subsection{Pbnm setting and Wyckoff positions}

TmFeO$_3$ adopts the orthorhombic $Pbnm$ setting (space group No.~62,
historical orthoferrite convention; the IT standard \emph{Pnma} differs
by an axis permutation). Lattice constants at 4\,K are $a=5.2534$\,\AA,
$b=5.5707$\,\AA, $c=7.6076$\,\AA~\cite{Marezio1970}. The four Fe
sublattices occupy Wyckoff position $4b$ ($\bar 1$ site symmetry):
\begin{equation}
\begin{aligned}
\mathrm{Fe}_1 &= (0,\,\half,\,\half), &\quad
\mathrm{Fe}_2 &= (\half,\,0,\,\half), \\
\mathrm{Fe}_3 &= (\half,\,0,\,0), &\quad
\mathrm{Fe}_4 &= (0,\,\half,\,0).
\end{aligned}
\label{eq:Fepos}
\end{equation}
The four Tm sublattices occupy Wyckoff position $4c$ (mirror $m$ site
symmetry, $z=\frac{1}{4}$ or $\frac{3}{4}$):
\begin{equation}
\begin{aligned}
\mathrm{Tm}_1 &= (0.0211,\,0.9284,\,\tfrac{3}{4}), \\
\mathrm{Tm}_2 &= (0.5211,\,0.5716,\,\tfrac{1}{4}), \\
\mathrm{Tm}_3 &= (0.4789,\,0.4284,\,\tfrac{3}{4}), \\
\mathrm{Tm}_4 &= (0.9789,\,0.0716,\,\tfrac{1}{4}).
\end{aligned}
\label{eq:Tmpos}
\end{equation}
The space-group generators in the $Pbnm$ setting are
\begin{equation}
S_1=\{2_{010}\,|\,\tfrac{1}{2}\tfrac{1}{2}0\},\quad
S_2=\{2_{100}\,|\,\tfrac{1}{2}00\},\quad I=\{\bar 1\,|\,000\},
\label{eq:Pbnmgens}
\end{equation}
generating the $D_{2h}$ point group $\{E,2_x,2_y,2_z,I,m_x,m_y,m_z\}$ with
the screw and glide translations as listed.

\subsection{Sublattice permutations}

Acting on the four Fe and four Tm sublattices, the generators
permute as
\begin{equation}
\begin{aligned}
\sigma^{\rm Fe}_{S_1}&=(14)(23), &\quad
\sigma^{\rm Fe}_{S_2}&=(12)(34), &\quad
\sigma^{\rm Fe}_{I}&=e,\\
\sigma^{\rm Tm}_{S_1}&=(14)(23), &\quad
\sigma^{\rm Tm}_{S_2}&=(12)(34), &\quad
\sigma^{\rm Tm}_{I}&=(14)(23).
\end{aligned}
\label{eq:sigmag}
\end{equation}
Inversion is non-trivial on Tm because $4c$ does not sit at an inversion
center; it is trivial on Fe because $4b$ does.

\subsection{Local sublattice frames}
\label{sec:localframes}

\paragraph{Definition.} The \emph{local frame} at sublattice $\mu$ is
the orthonormal triad $\{\hat e^a_\mu\}_{a=x,y,z}$ obtained by rotating
the global crystal triad $\{\hat e^a\}=\{\hat a,\hat b,\hat c\}$ via the
proper-rotation part of the Pbnm space-group element that maps site $1$
to site $\mu$:
\begin{equation}
\hat e^a_\mu = (R_\mu)^a{}_b\,\hat e^b,\qquad
R_\mu = \text{point-group part of }g_\mu,
\label{eq:Rmudef}
\end{equation}
where $g_\mu\in Pbnm$ is any operation with $g_\mu\,\bfr_1 \equiv \bfr_\mu$
mod lattice. Because the four sites in a Wyckoff orbit are
related by the four cosets of the trivial site-symmetry stabilizer in
$D_{2h}/\{E,I\} \simeq D_2$, exactly four such rotational parts arise:
\begin{equation}
\{R_\mu\}_{\mu=1}^{4} = \{E,\;C_{2x},\;C_{2y},\;C_{2z}\}.
\label{eq:RmuKlein}
\end{equation}
This is the Klein four-group of proper $\pi$-rotations about the three
crystal axes, isomorphic to $\mathbb Z_2\times\mathbb Z_2$.

\paragraph{Derivation for Fe ($4b$).} The Fe$^{3+}$ ions occupy site
symmetry $\bar 1$, so the stabilizer of $\mathrm{Fe}_1$ is just
$\{E,I\}$ and the four cosets of $\{E,I\}$ in $D_{2h}$ are represented
by $\{E,\,C_{2x},\,C_{2y},\,C_{2z}\}$. Denote by $g_\mu\in Pbnm$ a
representative whose point-group part is the $\mu$-th of these four
rotations; we label the four sublattices by
\begin{equation}
\mathrm{Fe}_\mu \equiv g_\mu\,\mathrm{Fe}_1\!\!\!\pmod{\mathrm{lattice}},
\qquad \text{point-part}(g_\mu)=R_\mu,
\label{eq:FemapPbnm}
\end{equation}
which is the convention used in Eq.~\eqref{eq:Fepos}. As diagonal
matrices in the crystal basis,
\begin{equation}
\begin{aligned}
R_1&=\diag(+,+,+), & R_2&=\diag(+,-,-),\\
R_3&=\diag(-,+,-), & R_4&=\diag(-,-,+),
\end{aligned}
\label{eq:Rmu}
\end{equation}
where $\eta^a_\mu=(R_\mu)^a{}_a$ is the eigenvalue of $S^a$ under the
proper rotation that maps site $1\to\mu$. The screw \emph{translations}
are absorbed in the position labels $\bfr_\mu$ and do not act on the
spin.

\paragraph{Spin transformation.} Spin is an axial vector. Under any
\emph{proper} rotation $R\in SO(3)$ both polar and axial vectors
transform as $V^a\to R^a{}_b V^b$. Since each $R_\mu$ is proper
($\det R_\mu = +1$ for all $\mu$), the spin at site $\mu$ in the
\emph{global} frame is related to the spin in the \emph{local} frame by
\begin{equation}
S^a_\mu\bigr|_{\rm global} = (R_\mu)^a{}_b\,\tilde S^b_\mu\bigr|_{\rm local},
\label{eq:Smutransform}
\end{equation}
i.e.\ a per-component sign flip of the chosen $\eta^a_\mu$. The screw
translations $\tfrac12\bfa+\tfrac12\bfb$, $\tfrac12\bfa$, etc., are
absorbed into the position labels $\bfr_\mu$ and do not act on the spin.

\paragraph{Derivation for Tm ($4c$).} The Tm$^{3+}$ ions sit at site
symmetry $.m.$ (the mirror $m_z$ at $z=\tfrac14,\tfrac34$). Its
stabilizer is $\{E,m_z\} = \{E,I\!\cdot\!C_{2z}\}$, so the four cosets
in $D_{2h}/\{E,m_z\}$ are
$\{E,m_z\},\{C_{2x},m_y\},\{C_{2y},m_x\},\{C_{2z},I\}$. Choosing the
proper-rotation representative in each coset gives \emph{the same}
Klein four-group Eq.~\eqref{eq:RmuKlein}. Concretely, the same
generators map $\mathrm{Tm}_1\to\mathrm{Tm}_\mu$ in the $4c$ orbit (up
to the $z=\tfrac14\leftrightarrow\tfrac34$ shift absorbed in $\bfr_\mu$):
\begin{equation}
R^{\rm Tm}_\mu = R^{\rm Fe}_\mu \quad(\mu=1,\ldots,4),
\label{eq:RTmequalsRFe}
\end{equation}
so the same matrices Eq.~\eqref{eq:Rmu} act on the projected dipoles
$\Jeff_\mu,\deff_\mu$ at every Tm site. The mirror $m_z$ in the
stabilizer acts trivially on the $\hat z$-axis chosen as a principal
quantization axis at the Tm site, so it imposes no further restriction
on the local triad beyond fixing the gauge $\hat e^z_\mu = \pm\hat c$.

\paragraph{Group structure.} The four matrices form the Klein
four-group under multiplication,
\begin{equation}
\Pi_{\mu\nu} \equiv R_\mu R_\nu = \diag(\eta_\mu^x\eta_\nu^x,\,
\eta_\mu^y\eta_\nu^y,\,\eta_\mu^z\eta_\nu^z) = R_{\mu\cdot\nu},
\label{eq:Pimu}
\end{equation}
where $\mu\cdot\nu$ is the Klein product (each component multiplied
mod $2$). This product controls the sign pattern of the
bond-projected Heisenberg exchange: a bond
$\langle\mathrm{Fe}_\mu,\mathrm{Fe}_\nu\rangle$ contributes
$J\,\bfS_\mu^{\rm loc}\!\cdot\!\Pi_{\mu\nu}\bfS_\nu^{\rm loc}$ in the
\emph{local} frame, which generates the four sign patterns
$\sigma^{F,G,C,A}$ used in the Bertaut decomposition
(Sec.~\ref{sec:Bertautalg}).

\subsection{Polar versus axial vectors under \texorpdfstring{$R_\mu$}{R\_mu}}
\label{sec:polaraxial}

\fix{The spin $\bfS$ and the projected magnetic dipole $\Jeff$ are
\emph{axial} vectors; the projected electric dipole $\deff$ is
\emph{polar}.} Under a proper rotation both transform identically (as
$R$); under inversion the polar vector picks up an extra sign while the
axial does not. The sublattice frames $R_\mu$ in
Eq.~\eqref{eq:Rmu} are improper $\det R_\mu = \pm 1$, and one must
distinguish:
\begin{equation}
R_\mu^{\rm axial} = \mathrm{diag}(\eta_\mu^{x},\eta_\mu^{y},\eta_\mu^{z}),
\qquad
R_\mu^{\rm polar} = \det(R_\mu)\,R_\mu^{\rm axial}.
\label{eq:RpolarRaxial}
\end{equation}
For $R_2=R_x(\pi)$ (a proper rotation, $\det=+1$), $R_\mu^{\rm polar}=
R_\mu^{\rm axial}$. For $R_3=R_x(\pi)R_z(\pi)$ ($\det=+1$ also, product
of two proper rotations), the same holds. So in the present case
$R_\mu^{\rm polar}=R_\mu^{\rm axial}$ for all four sublattices, and the
distinction does not bite. We retain the bookkeeping for clarity in the
emission formula (Sec.~\ref{sec:emission}), where $\partial_t\bfP$ and
$\partial_t\mathbf{M}$ enter together and the local-to-global rotation
of $\deff$ versus $\Jeff$ must be applied with the right $R_\mu^{\rm
polar/axial}$.

\subsection{Action on the projected qutrit subspace}
\label{sec:Rxnonunitarity}

A subtle point that affects the validity of the truncation to the lowest
three Tm singlets: the rotation $R_x(\pi) = e^{-i\pi J_x}$ acts on a
single $J=6$ multiplet exactly via Wigner $d$-functions, $R_x(\pi)|J,m\rangle
= (-1)^{J-m}|J,-m\rangle$, and therefore does \emph{not} leak out of the
$J=6$ multiplet. \fix{The leakage out of the truncated three-level subspace
$\{|E_1\rangle,|E_2\rangle,|E_3\rangle\}$ arises because the CEF
parameters $B_l^{|m|}$ and $B_l^{-|m|}$ are both nonzero (Sec.~\ref{sec:CEF}),
breaking the $m\!\to\!-m$ symmetry as a separate good quantum number.}
Hence eigenstates $|E_\alpha\rangle$ do not have definite $m\to-m$ parity,
and $R_x(\pi)|E_1\rangle$ is a linear combination of all even-$m_J$
eigenstates including $|E_4\rangle, |E_5\rangle,\ldots$ at energies
$\gtrsim 14$\,meV. Truncating to $\{E_1,E_2,E_3\}$ then makes the projected
matrix
\begin{equation}
[U_{R_x(\pi)}]_{\alpha\beta} \equiv \langle E_\alpha|e^{-i\pi J_x}|E_\beta\rangle
\qquad(\alpha,\beta\in\{1,2,3\})
\label{eq:URxproj}
\end{equation}
non-unitary by an amount $\|U^\dagger U - \mathds{1}\|_{\rm op} \approx 1.3$
in operator norm (numerically, with the CEF parameters of
Sec.~\ref{sec:CEF}); the missing weight sits in the discarded singlets.

For modeling purposes this means we should not use $R_x(\pi)$ as a
unitary symmetry on the qutrit. Instead, we use $R_\mu$ only at the
\emph{vector} level, where it acts as a rotation on the projected
dipoles $\Jeff,\deff$ themselves (3-vectors), not on the underlying
Hilbert space. The qutrit Hilbert space has the same Gell-Mann basis on
every sublattice and is sublattice-independent in our convention.

% =====================================================================
\section{Crystal-field model for Tm\texorpdfstring{$^{3+}$}{3+}}
\label{sec:CEF}

\subsection{Stevens-operator Hamiltonian}

\paragraph{Local-frame convention for the Tm sector.} All
single-ion operators introduced in this section --- the Stevens
operators $\hat O_l^m(J)$, the angular-momentum components
$J_x,J_y,J_z$, the projected $\Jeff,\deff$, and the Gell-Mann
$\lambda^a$ in Sec.~\ref{sec:Tgauge} below --- are defined
\emph{in the local triad $\{\hat e^a_\mu\}$ of
Sec.~\ref{sec:localframes}} at a representative Tm site, taken to be
$\mu=1$. Because the four Tm sublattices are related by the proper
rotations $R_\mu\in\{E,C_{2x},C_{2y},C_{2z}\}$
(Eq.~\eqref{eq:RTmequalsRFe}), the same Stevens parameters $B_l^m$
diagonalize the same operator $H_{\rm CEF}^{\rm Tm}$ at every
sublattice when one uses the local triad on that sublattice. The full
four-sublattice CEF Hamiltonian is therefore
\begin{equation}
H_{\rm CEF}^{\rm Tm,\,4-sl} = \sum_{\mu=1}^{4}\sum_{l,m} B_l^m\,
\hat O_l^m\!\bigl(\tilde J_\mu\bigr),
\qquad \tilde J^a_\mu = (R_\mu^{-1})^a{}_b\,J^b_\mu,
\label{eq:HCEF4sl}
\end{equation}
where $\tilde J^a_\mu$ is the local-frame component of the global
angular momentum $J^b_\mu$ at site $\mu$. Equivalently, in the local
basis the spectrum, the eigenstates $|E_n\rangle$, and the matrix
elements $M_{nn'}^{(a)}$ tabulated in Eq.~\eqref{eq:Mnumeric} are
$\mu$-independent; the $\mu$-dependence enters only through $R_\mu$
when one re-expresses local-frame matrix elements as global tensor
components --- exactly the construction used in
Sec.~\ref{sec:LambdaBertaut} to assemble Bertaut combinations $\Lambda^a_X$.

In the orthorhombic $Pbnm$ environment, the Tm$^{3+}$ ion (ground
multiplet $^3H_6$, $J=6$) sits at a site of $C_s$ symmetry, with the
mirror $m_z$ perpendicular to the $c$ axis at $z=\tfrac14$. \fix{Acting on
the $J=6$ multiplet, $m_z = R_z(\pi)\cdot I = R_z(\pi)$ on axial $\mathbf J$
(since inversion is trivial on axial vectors), so $m_z|J,m_J\rangle =
(-1)^{m_J}|J,m_J\rangle$. A Stevens operator $\hat O_l^m$ has matrix
elements with $\Delta m_J = \pm|m|$, so $m_z\hat O_l^m m_z^{-1} =
(-1)^{|m|}\hat O_l^m$. Invariance under the site mirror therefore
requires $|m|$ \emph{even}; both $\cos(m\phi)$ and $\sin(m\phi)$ partners
are independently allowed because $m_z$ does not relate $+|m|$ and
$-|m|$ (this is the role of the additional $C_2$ axes of the full
$D_{2h}$, which are absent here). The CEF Hamiltonian therefore has the
restricted form}
\begin{equation}
H_{\rm CEF}^{\rm Tm} = \sum_{l=2,4,6}\;\sum_{\substack{m=-l\\ |m|\,\rm even}}^{+l} B_l^m\,\hat O_l^m(J),
\label{eq:HCEF}
\end{equation}
where $\hat O_l^m$ are the Stevens operator equivalents in the convention
of Hutchings~\cite{Hutchings1964}, with the Stevens factors $\alpha_J=
1/99$, $\beta_J=8/49\,005$, $\gamma_J=-5/891\,891$ for Tm$^{3+}$
\emph{absorbed} into the quoted $B_l^m$. The $O_l^{|m|}$ are the
$\cos(m\phi)$ partners (real combinations of $J_+^m$ and $J_-^m$) and
$O_l^{-|m|}$ are the $\sin(m\phi)$ partners (imaginary combinations);
explicit formulas are tabulated in~\cite{Hutchings1964}, e.g.,
$O_2^0=3J_z^2-J(J+1)$, $O_2^2=\half(J_+^2+J_-^2)$, $O_2^{-2}=
\frac{1}{2i}(J_+^2-J_-^2)$. \fix{Note that the presence of
the $\sin$-partners $O_l^{-|m|}$ (with purely imaginary matrix elements
in the $|J,m_J\rangle$ basis) makes $H_{\rm CEF}^{\rm Tm}$ complex
Hermitian, \emph{not} real symmetric, in this basis whenever any
$B_l^{-|m|}\neq 0$; in particular the 4K parameters of
Tab.~\ref{tab:Stevens} satisfy $B_2^{-2}\neq 0$. Eigenstates of
Eq.~\eqref{eq:HCEF} are therefore generically complex in the
$|J,m_J\rangle$ basis, and the ``real-wavefunction'' gauge invoked in
Sec.~\ref{sec:Tgauge} below should be understood as a $T$-eigenstate
gauge (which is always achievable because $[H_{\rm CEF},T]=0$, see
Sec.~\ref{sec:Tgauge}), not as a literal reality condition on the
expansion coefficients $c_{\alpha,m}$.}

The fitted Stevens parameters from inelastic neutron
spectroscopy~\cite{Skorobogatov2020} are reproduced in
Table~\ref{tab:Stevens}. They are intended to reproduce the lowest four
singlet energies $\epsilon_1=0$, $\epsilon_2=1.920$\,meV,
$\epsilon_3=7.844$\,meV, $\epsilon_4=14.27$\,meV. \fix{Direct
diagonalization of Eq.~\eqref{eq:HCEF} on the 13-dim $J=6$ multiplet
using Tab.~\ref{tab:Stevens} (Hutchings convention, Stevens factors
absorbed; verified by the script in App.~\ref{app:script}) yields
$\epsilon_{1\ldots 6}-\epsilon_1 = (0,\,1.920,\,7.844,\,14.269,\,21.222,\,33.661)$\,meV,
reproducing the three lowest singlets within rounding.}

\begin{table}[t]
\caption{Stevens parameters $B_l^m$ for Tm$^{3+}$ in TmFeO$_3$ at low
temperature, in meV, in the Hutchings convention with Stevens factors
absorbed. \fix{These are the values that, fed into
Eq.~\eqref{eq:HCEF} and diagonalized on the full 13-dim $J=6$ space,
reproduce the cited $\epsilon_{2,3,4}$ and the wavefunction amplitudes
of Tab.~\ref{tab:CFwfs} verbatim (App.~\ref{app:script}).}}
\label{tab:Stevens}
\begin{ruledtabular}
\begin{tabular}{lr|lr|lr}
$B_2^0$  & $-0.529$ & $B_4^0$  & $-1.3\times 10^{-4}$ & $B_6^0$  & $+2.0\times 10^{-6}$ \\
$B_2^2$  & $-0.135$ & $B_4^2$  & $-1.7\times 10^{-3}$ & $B_6^2$  & $-1.1\times 10^{-5}$ \\
$B_2^{-2}$ & $+1.279$ & $B_4^{-2}$ & $+3.29\times 10^{-3}$ & $B_6^{-2}$ & $-9.0\times 10^{-6}$ \\
         &        & $B_4^4$  & $-1.22\times 10^{-3}$ & $B_6^4$  & $+6.1\times 10^{-5}$ \\
         &        & $B_4^{-4}$ & $-9.57\times 10^{-3}$ & $B_6^{-4}$ & $+3.0\times 10^{-6}$ \\
         &        &          &           & $B_6^6$  & $-9.0\times 10^{-6}$ \\
         &        &          &           & $B_6^{-6}$ & $0$ \\
\end{tabular}
\end{ruledtabular}
\end{table}

\subsection{Three-singlet manifold and mirror parity}
\label{sec:mJparity}

Diagonalizing Eq.~\eqref{eq:HCEF} on the 13-dimensional $J=6$ multiplet
yields three nondegenerate singlets in the bottom 8\,meV. The site
symmetry $C_s$ has only two irreducible representations, $A_1$ (mirror
$m_z$ even) and $A_2$ (mirror $m_z$ odd). The mirror $m_z =
R_z(\pi)$ acts on $|J,m_J\rangle$ as $(-1)^{m_J}$, so the unique
selection rule is $m_J \bmod 2$: even-$m_J$ states are $A_1$ and
odd-$m_J$ states are $A_2$. Diagonalization gives
\begin{equation}
\begin{aligned}
|E_1(A_1)\rangle &= \sum_{m\,{\rm even}} c_{1,m}|m\rangle,\\
|E_2(A_1)\rangle &= \sum_{m\,{\rm even}} c_{2,m}|m\rangle,\\
|E_3(A_2)\rangle &= \sum_{m\,{\rm odd}} c_{3,m}|m\rangle.
\end{aligned}
\label{eq:wfparity}
\end{equation}
The map $m\to-m$ is \emph{not} an element of
$C_s$ (only $m_z$ is, which gives $m\bmod 2$, not $m\to-m$); hence
the two $A_1$ states $|E_{1,2}\rangle$ are general real linear
combinations of $|m\rangle$ with even $m$, with no further parity
constraint between $c_{\alpha,m}$ and $c_{\alpha,-m}$. The matrix
element $\langle E_1|J_z|E_2\rangle\neq 0$ does not require any such
discrete parity assignment between them; it only requires the two
states to be linearly independent within the 7-dimensional even-$m$
block. Likewise $\langle E_2|J_x|E_3\rangle\neq 0$ requires that
$|E_2\rangle$ have non-vanishing amplitude on $|m\rangle$ for some $m$
adjacent (in the $J_\pm$ ladder) to a non-vanishing amplitude of
$|E_3\rangle$, which is generic.

The numerical wavefunction amplitudes, with the parity assignment of
Eq.~\eqref{eq:wfparity}, are listed in Table~\ref{tab:CFwfs}. The signs
of $c_{2,\pm m}$ are opposite (antisymmetric) and we list only $c_{2,m>0}$
with the convention $c_{2,-m}=-c_{2,m}$. \fix{In particular, antisymmetry
forces the $m=0$ amplitude $c_{2,0}=-c_{2,0}=0$ identically; any
nonzero entry there in the table below is to be read as numerical
rounding noise from the diagonalizer in a non-$T$-symmetric gauge, not
as a physical amplitude.}

\begin{table}[t]
\caption{Wavefunction amplitudes for the lowest three Tm singlets (from
diagonalizing Eq.~\eqref{eq:HCEF}, in the $T$-symmetric gauge of
Sec.~\ref{sec:Tgauge}). $|E_1(A_1)\rangle$ and $|E_3(A_2)\rangle$ are
symmetric in $m\to-m$ ($t_1=+1$, $t_3=-1$ with the $(-1)^m$-sign rule);
$|E_2(A_1)\rangle$ is antisymmetric ($t_2=-1$), so $c_{2,-m}=-c_{2,m}$
and in particular $c_{2,0}\equiv 0$.}
\label{tab:CFwfs}
\begin{ruledtabular}
\begin{tabular}{c|ccc}
$m_J$ & $c_{1,m}$ & $c_{2,m}$ & $c_{3,m}$ \\
\hline
$\pm 6$ & $+0.580$ & $+0.636$ & ---     \\
$\pm 5$ & ---      & ---      & $+0.484$ \\
$\pm 4$ & $+0.322$ & $+0.286$ & ---     \\
$\pm 3$ & ---      & ---      & $+0.402$ \\
$\pm 2$ & $+0.212$ & $+0.118$ & ---     \\
$\pm 1$ & ---      & ---      & $+0.323$ \\
$0$     & $+0.168$ & $\;\;\;0$\footnote{Forced to vanish by $t_2=-1$ antisymmetry; any nonzero value is rounding noise.} & ---     \\
\end{tabular}
\end{ruledtabular}
\end{table}

Normalization: $\sum_m |c_{\alpha,m}|^2 = 1$ within rounding for all
three states (verified to 0.5\% accuracy; deviations come from neglected
$|c|<0.02$ admixtures at $|m|=4,6$ and from the $\epsilon_4$ singlet
contribution that is excluded from the truncation).

\subsection{Projected angular-momentum matrix elements: \texorpdfstring{$T$}{T}-symmetric gauge}
\label{sec:Tgauge}

For a non-Kramers ion ($J=6$, integer) every state satisfies $T^2=+1$,
and $H_{\rm CEF}$ commutes with $T$ because every Stevens operator of
even rank $l=2,4,6$ is built from products of two angular-momentum
factors (each $T$-odd) and is therefore $T$-even. \fix{Even though
$H_{\rm CEF}$ is generically complex Hermitian in the $|J,m\rangle$
basis when $B_l^{-|m|}\neq 0$, the commutation $[H_{\rm CEF},T]=0$
guarantees that every non-degenerate eigenstate $|E_\alpha\rangle$ can
be taken to be a $T$-eigenstate by a global phase rotation:}
\begin{equation}
T|E_\alpha\rangle = t_\alpha|E_\alpha\rangle,\qquad t_\alpha\in\{+1,-1\}.
\label{eq:Teigvals}
\end{equation}
\fix{In any gauge in which the expansion $|E_\alpha\rangle =
\sum_m c_{\alpha,m}|m\rangle$ has $c_{\alpha,m}\in\mathbb R$ (which can
always be arranged \emph{after} the $T$-eigenstate phase fix is
performed; the resulting reality of $c_{\alpha,m}$ is then a
\emph{consequence}, not an assumption),} direct computation using
$T|J,m\rangle = (-1)^{J-m}|J,-m\rangle$ gives
\begin{equation}
T|E_\alpha\rangle = \sum_m c_{\alpha,m}^{*}\,(-1)^{J-m}|{-m}\rangle
= t_\alpha\sum_m c_{\alpha,m}|m\rangle,
\end{equation}
whence $c_{\alpha,m} = t_\alpha\,(-1)^{J-m}\,c_{\alpha,-m}^{*}$. With
real $c_{\alpha,m}$ and $J=6$ this collapses to
$c_{\alpha,m}=t_\alpha(-1)^{m}c_{\alpha,-m}$, so $t_\alpha = +1$ requires
the $m$-parity rule $c_{\alpha,m}=(-1)^m c_{\alpha,-m}$ (symmetric for
even $m$, antisymmetric for odd $m$), and $t_\alpha = -1$ flips both
rules.  For the model wavefunctions of Tab.~\ref{tab:CFwfs}, the script
in App.~\ref{app:script} verifies $t_1 = +1$, $t_2 = -1$, $t_3 = -1$:
$|E_1\rangle$ has only even $m$ and is $m\to -m$ symmetric (consistent
with $t_1 (-1)^{\rm even} = +1$); $|E_2\rangle$ has only even $m$ and is
antisymmetric (consistent with $t_2(-1)^{\rm even}=-1$); $|E_3\rangle$
has only odd $m$ and is symmetric (consistent with
$t_3 (-1)^{\rm odd} = +1$, i.e.\ $t_3 = -1$).

Because two of the three states are $T$-anti-invariant, the
$T$-symmetric gauge in which all projected operators $\Jeff_a$ are
purely imaginary requires multiplying \emph{both} $|E_2\rangle$ and
$|E_3\rangle$ by $i$:
\begin{equation}
|E_2\rangle\longrightarrow i|E_2\rangle,\qquad
|E_3\rangle\longrightarrow i|E_3\rangle,\qquad
|E_1\rangle\longrightarrow |E_1\rangle.
\label{eq:Tgauge}
\end{equation}
With this gauge fix every off-diagonal matrix element
$\langle E_\alpha|J_a|E_\beta\rangle$ ($\alpha\neq\beta$, $a=x,y,z$) is
purely imaginary (verified to $10^{-16}$ in App.~\ref{app:script}, Sec.~4).

In this $T$-symmetric gauge the projected operators are
\begin{equation}
\begin{aligned}
J_z^{\rm eff} &= +5.264\,i\,|E_2\rangle\langle E_1| + \mathrm{H.c.}
= +5.264\,\lambda^2,\\
J_x^{\rm eff} &= +2.392\,i\,|E_3\rangle\langle E_1| + 0.913\,i\,|E_3\rangle\langle E_2| + \mathrm{H.c.}\\
&= +2.392\,\lambda^5 + 0.913\,\lambda^7,\\
J_y^{\rm eff} &= -2.787\,i\,|E_3\rangle\langle E_1| + 0.466\,i\,|E_3\rangle\langle E_2| + \mathrm{H.c.}\\
&= -2.787\,\lambda^5 + 0.466\,\lambda^7,
\end{aligned}
\label{eq:Jeff}
\end{equation}
where the Gell-Mann matrices are in the ordered basis
$(|E_1\rangle,|E_2\rangle,|E_3\rangle)$. The numerical
$|\langle E_1|J_z|E_2\rangle|^2 = 5.264^2 = 27.71$ matches the value
extracted from inelastic neutron scattering~\cite{Skorobogatov2020} to
better than 1\%.

\subsection{Projected electric dipole}
The projected electric dipole on the Tm site (from admixture of opposite-parity 
$5d/4f^{n+1}5d^{-1}$ states, treated phenomenologically~\cite{JuddOfelt}) 
has matrix elements with the same selection rules as $J_x$ and $J_z$ at the
$C_s$ site:
\begin{equation}
\begin{aligned}
&d_x^{\rm eff} \in \mathrm{span}(\lambda^1),\qquad
d_z^{\rm eff} \in \mathrm{span}(\lambda^4,\lambda^6),\\
&d_y^{\rm eff} \in \mathrm{span}(\lambda^1)\\
&\quad(\text{site mirror reverses sign of }d_y).
\end{aligned}
\label{eq:deff}
\end{equation}
Quantitative magnitudes from a Judd-Ofelt analysis of orthoferrite
oscillator strengths are typically $\sim 10^{-2}$\,debye~\cite{Wood1968}; we
treat the coefficients $p_x, p_{z4}, p_{z6}$ as phenomenological,
constrained by the polarization-resolved 2DCS amplitudes.

\subsection{Mirror classification of Gell-Mann generators}
The site mirror $m_z = R_z(\pi)$ acts on the qutrit as
$\mathrm{diag}(+1,+1,-1)$ in the $(E_1,E_2,E_3)$ ordering (because $E_{1,2}$
have even $m_J$ and $E_3$ has odd $m_J$, and $R_z(\pi)|m\rangle=(-1)^m|m\rangle$).
\fix{We adopt the standard $C_s\times\{T\}$ labeling convention:
the Mulliken letter denotes the spatial (mirror) parity
($A_1\equiv$ mirror-even, $A_2\equiv$ mirror-odd), and the superscript
$\pm$ denotes time-reversal parity ($+\equiv T$-even, $-\equiv T$-odd).
This is the convention used throughout Sec.~\ref{sec:vanvleck} and
beyond for the Fe Bertaut bilinears
$\mathcal I^{\rm Fe}_{A_1^+},\mathcal I^{\rm Fe}_{A_2^+}$.}
Using this convention, the eight Gell-Mann generators classify as
\begin{equation}
\boxed{
\begin{aligned}
A_1^+ \text{ (mirror-even, $T$-even):} &\quad \lambda^1,\lambda^3,\lambda^8,\\
A_1^- \text{ (mirror-even, $T$-odd):} &\quad \lambda^2,\\
A_2^+ \text{ (mirror-odd, $T$-even):} &\quad \lambda^4,\lambda^6,\\
A_2^- \text{ (mirror-odd, $T$-odd):} &\quad \lambda^5,\lambda^7.
\end{aligned}}
\label{eq:lambdaclass}
\end{equation}
\fix{(Earlier drafts of this manuscript contained a labeling typo in
which $\lambda^2$ was written as ``$A_2^+$'' and $\lambda^{4,6}$ as
``$A_1^-$''; the assignments above are the corrected ones, internally
consistent with the Fe-bilinear classification of Eq.~\eqref{eq:IA1}
and Eq.~\eqref{eq:IA2}.)}
This classification organizes every coupling in the Hamiltonian and
every selection rule in the 2DCS analysis below. Direct verification:
in the ordered $(E_1,E_2,E_3)$ basis the diagonal mirror $U_{m_z}=
\mathrm{diag}(1,1,-1)$ commutes with $\lambda^{1,2,3,8}$ (mirror-even)
and anticommutes with $\lambda^{4,5,6,7}$ (mirror-odd), while complex
conjugation in this basis fixes the imaginary Gell-Mann matrices
$\lambda^{2,5,7}$ as $T$-odd and the real ones
$\lambda^{1,3,4,6,8}$ as $T$-even (in the $T$-symmetric gauge of
Sec.~\ref{sec:Tgauge}, where $T$ acts as plain complex conjugation
on the qutrit basis).

\subsection{Bertaut algebra for the Tm Gell-Mann operators}
\label{sec:LambdaBertaut}

The classification Eq.~\eqref{eq:lambdaclass} is a \emph{single-site}
$C_s$ classification under the local site mirror $m_z$. The full
crystal symmetry $Pbnm$ further organizes the four-sublattice sums of
$\lambda^a$ into Bertaut combinations transforming as definite irreps
of the global point group --- exactly as the Fe spins are organized
into $\bfF,\bfG,\bfC,\bfA$. Carrying out this step is essential
because (i) it tells us \emph{which} sublattice combination of
$\lambda^a$ is excited by a uniform field, and (ii) it tells us which
combination couples to each Fe Bertaut mode in
Eqs.~\eqref{eq:Hanisogen}--\eqref{eq:VVcoherence}. We do this here.

\paragraph{Convention.} As stated in Sec.~\ref{sec:CEF}, the
qutrit Hilbert space at every Tm sublattice $\mu$ uses the \emph{same}
ordered basis $\{|E_1\rangle,|E_2\rangle,|E_3\rangle\}$ defined in the
\emph{local} frame at sublattice 1; the Gell-Mann operator $\lambda^a$
then has the same $3\times 3$ matrix at every sublattice. The
sublattice dependence enters not in $\lambda^a$ itself but in the
\emph{interpretation} of $\lambda^a$ as a global-frame tensor
component. Concretely, the projected operator that on sublattice~1
represents the local-frame component $J_a^{\rm loc}$ represents on
sublattice~$\mu$ the global-frame component $(R_\mu^{-1})_{ab}\,J_b^{\rm glob}$
(equivalently, $J_a^{\rm loc}|_\mu = (R_\mu)_{ab}\,J_b^{\rm glob}$),
because $R_\mu$ is the orthogonal map from local to global at sublattice $\mu$.

\paragraph{$\sigma$-vectors of axial / polar generators.} For the
Klein-four sublattice frames $\{R_\mu\}=\{E,C_{2x},C_{2y},C_{2z}\}$
of Eq.~\eqref{eq:Rmu}, the diagonal of $R_\mu$ in the cubic axes gives
the sign carried by each axial component~$a\in\{x,y,z\}$ at sublattice
$\mu$:
\begin{equation}
\eta^a_\mu \equiv (R_\mu)_{aa},\qquad
\begin{aligned}
\eta^x &= (+,+,-,-)\equiv\sigma_C,\\
\eta^y &= (+,-,+,-)\equiv\sigma_A,\\
\eta^z &= (+,-,-,+)\equiv\sigma_G,
\end{aligned}
\label{eq:sigmaTm}
\end{equation}
together with the trivial $\sigma_F=(+,+,+,+)$. (All $R_\mu$ are
proper rotations, so polar and axial vectors carry the \emph{same}
$\sigma$-vectors here, cf.\ Sec.~\ref{sec:polaraxial}.) Eq.~\eqref{eq:sigmaTm}
recovers the standard Bertaut assignment for $Pbnm$ orthoferrites:
the $\Gamma_2$ irrep contains $F_x, C_y, G_z$ --- exactly the three
axial Bertaut variables whose $\sigma$-vector is $\sigma_C, \sigma_A, \sigma_G$
weighted by the local-axis label $\{x,y,z\}$.

\paragraph{$\sigma$-vectors of the eight Gell-Mann generators.}
Using the projected-operator decomposition
Eqs.~\eqref{eq:Mnumeric}--\eqref{eq:deff}, each $\lambda^a$ inherits
the $\sigma$-vector of the local-frame Cartesian operator that it
represents. \fix{The ``represents (local)'' entries below are the
\emph{inverse} of the matrix-element decomposition
$J_x^{\rm eff} = a\lambda^5 + b\lambda^7$,
$J_y^{\rm eff} = c\lambda^5 + d\lambda^7$ of
Eq.~\eqref{eq:Mnumeric}: solving for $\lambda^{5,7}$ in terms of
$J_{x,y}^{\rm eff}$ gives $\lambda^5 = (dJ_x - bJ_y)/(ad-bc)$ and
$\lambda^7 = (-cJ_x + aJ_y)/(ad-bc)$, with the common denominator
$ad-bc \approx 3.66$ for $(a,b,c,d)=(2.392,0.913,-2.787,0.466)$. The
entries below are written without that overall normalization, which
does not affect the $\sigma$-vector decomposition (only the absolute
scale).}
\begin{widetext}
\begin{equation}
\boxed{
\renewcommand{\arraystretch}{1.15}
\begin{array}{c|c|c|l}
\lambda^a & \text{represents (local), $\propto$} & \sigma\text{-vector} & \text{couples (uniform field) to}\\
\hline
\lambda^2 & J_z & \sigma_G=(+,-,-,+) & G_z\text{ (Fe N\'eel)}\\
\lambda^5 & d J_x - b J_y & d\,\sigma_C - b\,\sigma_A & \text{lin.\ comb of } C_x, A_y\\
\lambda^7 & -c J_x + a J_y & -c\,\sigma_C + a\,\sigma_A & \text{lin.\ comb of } C_x, A_y\\
\lambda^4 & d_z\text{ (polar)} & \sigma_G & G_z^{\rm pol}\\
\lambda^6 & d_z\text{ (polar)} & \sigma_G & G_z^{\rm pol}\\
\lambda^1 & d_{x,y}\text{ (polar)} & \text{lin.\ comb of }\sigma_C,\sigma_A & \text{electric drive}\\
\lambda^3,\lambda^8 & \text{populations} & \sigma_F & \text{uniform } n_a\\
\end{array}}
\label{eq:LambdaBertaut}
\end{equation}
\end{widetext}
with $(a,b,c,d)=(2.392, 0.913,-2.787,0.466)$ from
Eqs.~\eqref{eq:Mnumeric}.

\paragraph{Bertaut sublattice combinations.} Define the four
Tm-sublattice combinations of any $\lambda^a$ in direct analogy with
Eq.~\eqref{eq:Bertautdef}:
\begin{equation}
\Lambda^a_X = \frac{1}{4}\sum_{\mu=1}^{4}\sigma^X_\mu\,\lambda^a_\mu,\qquad
X\in\{F,G,C,A\}.
\label{eq:LambdaXdef}
\end{equation}
Then a uniform external perturbation that couples to the global axial
component $J_a^{\rm glob}$ acts only on the Bertaut combination whose
$\sigma$-vector matches $\eta^a$:
\begin{equation}
\sum_\mu (R_\mu \mathbf h)_{a}\,(J_a^{\rm loc})|_\mu
= \sum_\mu \eta^a_\mu\,h_a\,(J_a^{\rm loc})|_\mu
\;\Longrightarrow\; h_a\cdot 4\Lambda^{[J_a^{\rm loc}]}_{X(a)},
\label{eq:uniformsel}
\end{equation}
where $X(a)\in\{F,C,A,G\}$ is the irrep of $\eta^a$. In particular,
\begin{itemize}[leftmargin=*]
\item Uniform $h_z$ excites $\Lambda^2_G$ and (via $\lambda^{4,6}$
electric channel for $e_z$) $\Lambda^{4,6}_G$ --- not $\Lambda^2_F$.
This is the precise meaning of the row $\lambda^2\to G_z$ in
Eq.~\eqref{eq:LambdaBertaut}.
\item Uniform $h_x$ excites the $\sigma_C$-component of
$d\Lambda^5 - b\Lambda^7$, and uniform $h_y$ the $\sigma_A$-component of
$-c\Lambda^5 + a\Lambda^7$. The four orthogonal Bertaut combinations
$\Lambda^{5,7}_{C,A}$ are independent dynamical variables but only
specific linear combinations are dipole-active under uniform drive.
\item Populations $\lambda^{3,8}$ have $\sigma_F$, so only the uniform
combination $\Lambda^{3,8}_F$ is generated by uniform pump (the
homogeneous Tm population grating).
\end{itemize}

\paragraph{Polar (electric) drive --- same $R_\mu$, different parity.}
The polar generators $\lambda^1$ ($d_x, d_y$ channel) and
$\lambda^{4,6}$ ($d_z$ channel) carry through the local-frame
construction \emph{identically} to the axial case, because the
Klein-four sublattice rotations $R_\mu\in\{E,C_{2x},C_{2y},C_{2z}\}$
are all proper ($\det R_\mu = +1$, Sec.~\ref{sec:polaraxial}): polar
and axial vectors share the same $\sigma$-vectors $\sigma_C, \sigma_A,
\sigma_G$. Concretely, the projected polar-dipole operator obeys the
same transport law as $J_a$,
\begin{equation}
d_a^{\rm loc}\big|_\mu = (R_\mu)_{ab}\,d_b^{\rm glob},\qquad
\eta^{a,\rm pol}_\mu = \eta^{a,\rm axial}_\mu = (R_\mu)_{aa},
\label{eq:dlocal}
\end{equation}
with no extra inversion sign because $R_\mu$ does not contain $I$
(the inversion has been pushed into the site stabilizer / position
labels in Sec.~\ref{sec:localframes}). The electric-field analogue
of Eq.~\eqref{eq:uniformsel} is therefore
\begin{equation}
\boxed{\;-\sum_\mu E_a\,d_a^{\rm glob}|_\mu
= -E_a\,\sum_\mu \eta^a_\mu\,d_a^{\rm loc}|_\mu
\;\Longrightarrow\; -E_a\cdot 4\Lambda^{[d_a^{\rm loc}]}_{X(a)},\;}
\label{eq:uniformselE}
\end{equation}
with the same Klein-four irrep map $X(a)\in\{F,C,A,G\}$. Reading off
the rows of Eq.~\eqref{eq:LambdaBertaut} now from the polar side:
\begin{itemize}[leftmargin=*]
\item Uniform $E_z$ excites $\Lambda^4_G$ and $\Lambda^6_G$
(the $G_z$-staggered electric channel between $|E_{1,2}\rangle$ and
$|E_3\rangle$). This is the electric-dipole counterpart of the
$\Lambda^2_G$ Zeeman channel and shares the same $\sigma_G$ pattern,
so $E_z$ and $h_z$ probe the \emph{same} Bertaut irrep --- the
distinguishing label is the matrix element ($\propto \deff$ for $E_z$
vs.\ $\propto\Jeff$ for $h_z$), not the sublattice symmetry.
\item Uniform $E_x$ excites $\sigma_C$-projected $\Lambda^1$, i.e.\
$E_x\cdot p_x\,\Lambda^1_C$, since the local-frame matrix element of
$d_x$ coincides (up to the phenomenological coefficient $p_x$) with
the $\lambda^1$ generator. Uniform $E_y$ similarly excites
$E_y\cdot p_y\,\Lambda^1_A$. Crucially the \emph{same} $\lambda^1$
operator carries both: the orthogonal $\sigma_{C}, \sigma_A$
combinations of the four $\lambda^1_\mu$ are independent dynamical
variables, and $E_x, E_y$ select different ones. (This is the polar
dual of the $\lambda^{5,7}$ statement above.)
\item No uniform $E$ component excites $\Lambda^a_F$ for $a=1,4,6$,
because no $\eta^{a,\rm pol}_\mu = +\!+\!+\!+$ exists in the Klein
four-group --- exactly as for axial drives.
\end{itemize}
The $|E_n\rangle$ basis-parity assignment from Sec.~\ref{sec:mJparity}
($|E_1\rangle,|E_2\rangle\in A^+$ and $|E_3\rangle\in A^-$ under
$m_z$) is what forces $d_z$ to populate \emph{only}
$\lambda^4,\lambda^6$ (off-diagonal between opposite-parity states)
and $d_{x,y}$ to populate \emph{only} $\lambda^1$ (off-diagonal
between same-parity states $|E_1\rangle\leftrightarrow|E_2\rangle$);
the local-frame map $R_\mu$ then promotes this single-site
selection rule to the four-sublattice $\sigma_{F,G,C,A}$ structure.

\paragraph{Why Fe has no polar dipole, but Tm does.} The same
local-frame analysis applied to the Fe sublattice would give
$d_a^{\rm loc}|_\mu = (R_\mu)_{ab}\,d_b^{\rm glob}$ as well; however,
the Fe site stabilizer $\{E,I\}$ at $4b$ contains \emph{inversion},
which is polar-odd: $I\,d_a I^{-1} = -d_a$. The on-site projected
electric dipole at Fe is therefore strictly zero by site symmetry,
and only the rare-earth (Tm $4c$, stabilizer $\{E,m_z\}$, no
inversion) carries an on-site electric dipole. This is why the
electric channel of Eqs.~\eqref{eq:LambdaBertaut} and \eqref{eq:uniformselE}
acts entirely through the Tm $\lambda^a$ sector --- there is no
parallel ``polar Bertaut amplitude'' on Fe to populate. Magnetic
dipole, by contrast, is even under $I$ and therefore survives at both
Fe (4b) and Tm (4c) sites, which is why the axial Bertaut variables
$\bfF,\bfG,\bfC,\bfA$ exist on Fe and the parallel $\Lambda^{2,5,7}_X$
combinations exist on Tm.

\paragraph{Cross-coupling to Fe.} The Van Vleck back-action
Eq.~\eqref{eq:Hanisogen}, written explicitly with sublattice indices,
becomes
\begin{widetext}
\begin{equation}
H_{\rm aniso\text{-}Tm} = -4\!\sum_{\mathbf n}\Bigl[
u_1\,\Lambda^1_{X_1}\!\cdot\!\mathcal I^{\rm Fe}_{X_1}
+ u_3\,\Lambda^3_F\,\mathcal I^{\rm Fe}_{A_1^+}
+ u_8\,\Lambda^8_F\,\mathcal I^{\rm Fe}_{A_1^+}
+ \cdots\Bigr],
\label{eq:HanisogenBertaut}
\end{equation}
\end{widetext}
i.e., for each $\lambda^a$ only the $\Lambda^a_X$ Bertaut combination
whose $\sigma$-vector matches the irrep of the Fe bilinear $\mathcal I^{\rm Fe}_{X}$
is allowed. In the $\Gamma_2$ phase the only Fe bilinears with nonzero
expectation value are in $A_1^+$ ($G_z^2, F_x^2, C_y^2,\ldots$, all
$\sigma_F$-pattern), so the coupled Tm combinations are
$\Lambda^{3,8}_F$ (population channel) and the $\sigma_F$-projection of
$\Lambda^1$ (polar coherence channel). The $\bar\chi^{2z}\,G_z\,\lambda^2$
linear coupling of Eq.~\eqref{eq:VVcoherence} acts on
$\Lambda^2_G$ specifically (matching $\sigma_G\!\cdot\!\sigma_G=\sigma_F$).

\paragraph{Consequence for the scalar EOMs of Sec.~\ref{sec:EOMscalar}.}
The amplitudes $\rho_{12}, \rho_{13}, \rho_{23}, n_1, n_2$ in
Eqs.~\eqref{eq:dotrho12}--\eqref{eq:dotn2} should be read as
\emph{Bertaut-projected} amplitudes:
$\rho_{12}\equiv\Lambda^2_G$ (the $G_z$-staggered combination), and
similarly $\rho_{13},\rho_{23}$ are the appropriate
$\sigma_C, \sigma_A$ linear combinations of $\Lambda^{5,7}$ that couple
to uniform $h_x, h_y$. With this identification, the scalar EOMs are
the correct Pbnm-allowed equations of motion at $\mathbf q=0$: every
amplitude carries a definite global irrep, every coupling vertex
respects the four-sublattice $\sigma$-vector algebra, and the
six-peak intrinsic count of Sec.~\ref{sec:MEcatalogue} is preserved.

% =====================================================================
\section{Fe four-sublattice Hamiltonian}
\label{sec:Fe}

\subsection{Microscopic Hamiltonian}

The Fe sector is the standard four-sublattice orthoferrite Hamiltonian
with Heisenberg exchange ($J_1$ in-plane NN, $J_{1c}$ out-of-plane NN,
$J_2$ NNN), Dzyaloshinskii--Moriya (DM) interaction with two independent
DM components $D_1, D_2$ allowed by $Pbnm$, and biaxial single-ion
anisotropy with axes locked to $\hat a, \hat b, \hat c$:
\begin{equation}
\begin{aligned}
H_{\rm Fe} =\;& \sum_{\langle ij\rangle_{\rm NN,\parallel}}\!\! J_1\, \bfS_i\cdot\bfS_j
 + \sum_{\langle ij\rangle_{\rm NN,c}}\!\! J_{1c}\,\bfS_i\cdot\bfS_j\\
 &+ \sum_{\langle\!\langle ij\rangle\!\rangle_{\rm NNN}}\!\! J_2\,\bfS_i\cdot\bfS_j\\
 &+ \sum_{\langle ij\rangle_{\rm NN,\parallel}}\!\! \bfD_{ij}\cdot(\bfS_i\times\bfS_j)\\
 &- \sum_i \bigl[K_a (S_i^x)^2 + K_b (S_i^y)^2 + K_c (S_i^z)^2\bigr],
\end{aligned}
\label{eq:HFe}
\end{equation}
with the standard hierarchy $J_1\sim 4.5$\,meV $\gg|D_1|,|D_2|\sim 0.1$\,meV
$\gg K_{a,b,c}\sim 10\,\mu$eV. The DM bond pattern dictated by $Pbnm$
is~\cite{Treves1965,Bertaut1968}:
\begin{equation}
\mathbf D^{(z=\frac{1}{2})}_{ij} = (0,+D_1,+D_2),\qquad
\mathbf D^{(z=0)}_{ij} = (0,-D_1,+D_2),
\label{eq:Dpattern}
\end{equation}
on in-plane NN bonds, with the partial sign reversal between $z=\half$
and $z=0$ planes required by the $S_1$ screw mapping the two planes (which
flips $\hat y$ but preserves $\hat z$). Out-of-plane and NNN couplings
have $\mathbf D=0$ by inversion symmetry of the bond midpoint.

\subsection{Local-frame form of \texorpdfstring{$H_{\rm Fe}$}{HFe}}
\label{sec:HFelocal}

For the Bertaut decomposition (Sec.~\ref{sec:Bertautalg}) and for
the linearization around the $\Gamma_2$ ground state in
Sec.~\ref{sec:EOMexpansion} it is essential to recast
Eq.~\eqref{eq:HFe} in the \emph{local} sublattice frames defined by
Eq.~\eqref{eq:Rmu}. Substituting
$\bfS_\mu = R_\mu\,\tilde\bfS^{\rm loc}_\mu$
(Eq.~\eqref{eq:Smutransform}) into each term gives:

\paragraph{Heisenberg term.}
\begin{equation}
\bfS_\mu\!\cdot\!\bfS_\nu
= \tilde\bfS^{\rm loc}_\mu\cdot\bigl(R_\mu^\top R_\nu\bigr)\,\tilde\bfS^{\rm loc}_\nu
= \tilde\bfS^{\rm loc}_\mu\!\cdot\!\Pi_{\mu\nu}\,\tilde\bfS^{\rm loc}_\nu,
\label{eq:Hloc1}
\end{equation}
where $\Pi_{\mu\nu}=R_\mu R_\nu = \diag(\eta^x_\mu\eta^x_\nu,
\eta^y_\mu\eta^y_\nu,\eta^z_\mu\eta^z_\nu)$
(Eq.~\eqref{eq:Pimu}; we used $R_\mu^\top=R_\mu^{-1}=R_\mu$ for
involutions). The diagonal sign pattern of $\Pi_{\mu\nu}$ is exactly
the per-component selection rule that produces the Bertaut sign vectors
$\sigma^{F,G,C,A}_{\mu}\sigma^{F,G,C,A}_{\nu}$ on quadratic mode
combinations: a NN bond $\langle\mu\nu\rangle$ between sublattices in
the same $z$-plane has $\Pi_{\mu\nu}=\diag(+,-,-)$ for the in-plane
pair $(\mu,\nu)=(1,4)$ etc., enforcing the antiferromagnetic
$G$-staggering of $S^z$ and the $A$/$C$-pattern of $S^{x,y}$ at the
classical $\Gamma_2$ minimum (uniform $\tilde S^z_\mu = +S$ in every
local frame).

\paragraph{DM term.} Under proper rotations the cross product
transforms as $R(\mathbf u\times\mathbf v)=R\mathbf u\times R\mathbf v$, so
\begin{equation}
\bfD_{\mu\nu}\!\cdot\!(\bfS_\mu\times\bfS_\nu)
= \tilde\bfD^{\rm loc}_{\mu\nu}\!\cdot\!\bigl(\tilde\bfS^{\rm loc}_\mu
   \times \Pi_{\mu\nu}\,\tilde\bfS^{\rm loc}_\nu\bigr),
\label{eq:Hloc2}
\end{equation}
with $\tilde\bfD^{\rm loc}_{\mu\nu}=R_\mu^{-1}\bfD_{\mu\nu}$. The
sign reversal of $D_1$ between the $z=0$ and $z=\half$ planes in
Eq.~\eqref{eq:Dpattern} is the nontrivial part of the $Pbnm$ pattern
that survives in the local coordinates, but it does \emph{not} imply a
standalone $F_xG_z$ Lifshitz invariant in the Bertaut convention of
Eq.~\eqref{eq:sigmavecs}. Direct bond-sum projection of the written DM
Hamiltonian gives the channel structure quoted in
Eq.~\eqref{eq:HDMBertaut}: $D_1$ mixes the $(A,F)$ and $(G,C)$ sectors,
whereas $D_2$ mixes the $(A,C)$ and $(F,G)$ sectors.

\paragraph{Anisotropy.} Because $K_{a,b,c}$ are diagonal in the
\emph{global} cubic basis and $R_\mu$ permutes signs but not axes,
\begin{equation}
\sum_a K_a (S^a_\mu)^2
= \sum_a K_a\bigl[(R_\mu)^a{}_b\,\tilde S^b_\mu\bigr]^2
= \sum_a K_a (\tilde S^a_\mu)^2,
\label{eq:Hloc3}
\end{equation}
i.e.\ the single-ion anisotropy is \emph{$\mu$-independent and uniform}
in the local frame. This is the formal reason every Fe sublattice has
its easy axis along the same local $\hat e^z_\mu$ (the global
$\hat c$-axis times $\eta^z_\mu$), so the $\Gamma_2$ phase has uniform
$\tilde S^z_\mu = +S$ in local frames.

\paragraph{Local-frame Hamiltonian.} Combining
Eqs.~\eqref{eq:Hloc1}--\eqref{eq:Hloc3} the Fe Hamiltonian Eq.~\eqref{eq:HFe}
reads, in local frames,
\begin{equation}
\begin{aligned}
\tilde H_{\rm Fe} =\;& \sum_{\langle\mu\nu\rangle} J_{\mu\nu}\,
   \tilde\bfS^{\rm loc}_\mu\!\cdot\!\Pi_{\mu\nu}\,\tilde\bfS^{\rm loc}_\nu\\
 &+ \sum_{\langle\mu\nu\rangle}\tilde\bfD^{\rm loc}_{\mu\nu}\!\cdot\!\bigl(
   \tilde\bfS^{\rm loc}_\mu\times\Pi_{\mu\nu}\,\tilde\bfS^{\rm loc}_\nu\bigr)\\
 &-\sum_\mu\bigl[K_a(\tilde S^x_\mu)^2 + K_b(\tilde S^y_\mu)^2 + K_c(\tilde S^z_\mu)^2\bigr],
\end{aligned}
\label{eq:HFelocal}
\end{equation}
in which the exchange-dominated reference configuration is the uniform
local state $\tilde\bfS^{\rm loc}_\mu = +S\hat e^z$ for all $\mu$.
The weak cantings that generate the orthoferrite order parameters then
appear as transverse, generally non-uniform fluctuations about this
reference state. Linearizing about it (Holstein--Primakoff or
$\mathbf L$-and-$\bfM$ method) is equivalent to the linearization of the
Bertaut variables $\bfF,\bfG,\bfC,\bfA$ used in
Sec.~\ref{sec:Bertautalg}, with the Klein-product $\Pi_{\mu\nu}$
playing the role of the bond-projector of
Eq.~\eqref{eq:Bertautinv}: the Fourier transform on the four-site
Klein index diagonalizes Eq.~\eqref{eq:HFelocal} into the four
$X\in\{F,G,C,A\}$ blocks.

\subsection{Bertaut decomposition}
\label{sec:Bertautalg}

Define the four Bertaut sublattice combinations
\begin{equation}
\begin{aligned}
\bfF_{\mathbf n} &= \quart\sum_\mu\bfS_{\mathbf n,\mu}, &
\bfG_{\mathbf n} &= \quart\sum_\mu\sigma^G_\mu\bfS_{\mathbf n,\mu},\\
\bfC_{\mathbf n} &= \quart\sum_\mu\sigma^C_\mu\bfS_{\mathbf n,\mu}, &
\bfA_{\mathbf n} &= \quart\sum_\mu\sigma^A_\mu\bfS_{\mathbf n,\mu},
\end{aligned}
\label{eq:Bertautdef}
\end{equation}
with sign vectors
\begin{equation}
\begin{aligned}
\sigma^F&=(+,+,+,+), & \sigma^G&=(+,-,-,+),\\
\sigma^C&=(+,+,-,-), & \sigma^A&=(+,-,+,-).
\end{aligned}
\label{eq:sigmavecs}
\end{equation}
These satisfy the Klein four-group multiplication
$\sigma^X_\mu\sigma^Y_\mu = \sigma^{XY}_\mu$ with the table
\begin{equation}
\begin{array}{c|cccc}
\cdot & F & G & C & A \\
\hline
F & F & G & C & A \\
G & G & F & A & C \\
C & C & A & F & G \\
A & A & C & G & F
\end{array}
\label{eq:Klein4}
\end{equation}
and the orthogonality $\sum_\mu \sigma^X_\mu\sigma^Y_\mu = 4\delta^{XY}$.

The inverse of Eq.~\eqref{eq:Bertautdef} is
\begin{equation}
\bfS_{\mathbf n,\mu} = \sum_X \sigma^X_\mu\,\bfX_{\mathbf n}\quad
\text{for }X\in\{F,G,C,A\},
\label{eq:Bertautinv}
\end{equation}
i.e., $\bfX$ are mode amplitudes in units of \emph{average sublattice
spin}, not summed spin: $|\bfF|^2+|\bfG|^2+|\bfC|^2+|\bfA|^2 = S^2$ at
the spin-length constraint $|\bfS_\mu|=S$.

\subsection{Effective fields and the per-cell normalization}
\label{sec:Heffnorm}

\fix{We adopt one normalization throughout the manuscript and propagate
it consistently}: the effective field on a single Fe sublattice spin is
\begin{equation}
\bfH^{\rm eff}_{\mathbf n,\mu} = -\frac{\partial H_{\rm Fe}}{\partial\bfS_{\mathbf n,\mu}}
\qquad\text{(per-site definition)},
\label{eq:Heffsite}
\end{equation}
which has units of energy (energy per spin component). The Bertaut
projection of the LLG equation gives, after substituting
Eq.~\eqref{eq:Bertautinv} and using the orthogonality
$\sum_\mu\sigma^X_\mu\sigma^Y_\mu=4\delta^{XY}$,
\begin{equation}
\begin{aligned}
\dot\bfX_{\mathbf n} &= -\frac{1}{\hbar}\sum_{Y,Z:\;Y\cdot Z=X}\bfY_{\mathbf n}\times\bfH^{Z,{\rm site}}_{\mathbf n},\\
\bfH^{Z,{\rm site}}_{\mathbf n} &\equiv \quart\sum_\mu\sigma^Z_\mu\bfH^{\rm eff}_{\mathbf n,\mu},
\end{aligned}
\label{eq:BertautEOM}
\end{equation}
where the Klein-product structure $Y\cdot Z = X$ ensures only the
allowed mode mixings appear. The auxiliary field $\bfH^{Z,{\rm site}}$
is the average per-site effective field weighted by $\sigma^Z_\mu$. With
this convention the precession prefactor is $1/\hbar$ (and equals unity
in the $\hbar=1$ dynamical equations used elsewhere in the manuscript),
not $\gamma_e=g_e\mu_B/\hbar$, which would apply only if
$\bfH^{\rm eff}$ were a magnetic field rather than an energy
derivative. There is likewise no extra prefactor in front of the
Heisenberg or anisotropy contributions:
\begin{equation}
\bfH^{F,{\rm site}}_{\rm aniso} = +2\mathbf K\bfF,\quad
\bfH^{G,{\rm site}}_{\rm aniso} = +2\mathbf K\bfG,\quad
\text{etc.,}
\label{eq:HXaniso_correct}
\end{equation}
where $\mathbf K = \mathrm{diag}(K_a,K_b,K_c)$. \emph{Derivation}: per
site, $-\partial[-K_a(S_\mu^x)^2]/\partial S_\mu^x = +2K_a S_\mu^x$, so
$\bfH^{\rm eff,aniso}_{\mathbf n,\mu} = 2\mathbf K\bfS_{\mathbf n,\mu}$.
Substituting $\bfS_\mu = \sum_X\sigma^X_\mu\bfX$ and projecting via
Eq.~\eqref{eq:BertautEOM},
\begin{equation}
\begin{aligned}
\bfH^{Z,{\rm site}}_{\rm aniso} &= \quart\sum_\mu\sigma^Z_\mu\cdot 2\mathbf K\sum_Y\sigma^Y_\mu\bfY\\
&= \quart\cdot 2\mathbf K\cdot 4\bfZ = 2\mathbf K\bfZ,
\end{aligned}
\label{eq:HXanisoderiv}
\end{equation}
matching Eq.~\eqref{eq:HXaniso_correct} cleanly.

Similarly the Heisenberg contribution gives, per Klein channel,
\begin{equation}
\begin{aligned}
\bfH^{F,{\rm site}}_{\rm Heis} &= -2 J^F\bfF, &
\bfH^{G,{\rm site}}_{\rm Heis} &= -2 J^G\bfG,\\
\bfH^{C,{\rm site}}_{\rm Heis} &= -2 J^C\bfC, &
\bfH^{A,{\rm site}}_{\rm Heis} &= -2 J^A\bfA,
\end{aligned}
\label{eq:HXHeis}
\end{equation}
with the Bertaut-projected exchange constants
\begin{equation}
\begin{aligned}
J^F &= +(z_1 J_1 + z_{1c} J_{1c} + z_2 J_2),\\
J^G &= -(z_1 J_1 + z_{1c} J_{1c} - z_2 J_2),\quad\text{etc.,}
\end{aligned}
\label{eq:JXdef}
\end{equation}
where $z_1, z_{1c}, z_2$ are coordination numbers (Sec.~\ref{sec:bondsymmetry})
and the relative signs follow from $\sigma^X$ on the bond endpoints.

\subsection{DM Hamiltonian in Bertaut form}

We derive the DM Bertaut Lifshitz invariants by direct bond-sum
projection. The DM term in Eq.~\eqref{eq:HFe} is a sum over in-plane NN
bonds. Using the directed bond list implicit in Eq.~\eqref{eq:Dpattern}
(four $\mathrm{Fe}_1\!\to\!\mathrm{Fe}_2$ bonds in the $z=\half$ plane with
$\mathbf D=(0,+D_1,+D_2)$ and four $\mathrm{Fe}_3\!\to\!\mathrm{Fe}_4$
bonds in the $z=0$ plane with $\mathbf D=(0,-D_1,+D_2)$) and the Bertaut
convention of Eq.~\eqref{eq:sigmavecs}, the $\mathbf q\!=\!0$ reduction of
the written bond Hamiltonian is (App.~\ref{app:DMderivation})
\begin{equation}
H_{\rm DM}^{\mathbf q\approx 0} = N_c\Bigl[
-8 D_1\bigl((\bfF\times\bfG)_y + (\bfC\times\bfA)_y\bigr)
+ 8 D_2\bigl((\bfG\times\bfC)_z - (\bfF\times\bfA)_z\bigr)
\Bigr],
\label{eq:HDMBertaut}
\end{equation}
where $N_c$ is the number of unit cells. The $D_1$ invariant
$(\bfF\times\bfG)_y = F_zG_x-F_xG_z$ is the standard
Treves--Bertaut weak-ferromagnetic Lifshitz invariant of $Pbnm$
orthoferrites~\cite{Treves1965,Bertaut1968,Herrmann1963}: evaluated on
the $\Gamma_2$ background it reduces to $-F_xG_z$, so the DM
energy at the background is $+8 N_c D_1 F_x G_z$. The $D_2$ invariants
$(\bfG\times\bfC)_z = G_xC_y-G_yC_x$ and $(\bfF\times\bfA)_z = F_xA_y-F_yA_x$
are the corresponding mixing terms between the $\Gamma_2$ partner modes
$(G_z,C_y)$ and the $\Gamma_4$ excited modes. The overall sign is tied
to the directed-bond convention; reversing every bond orientation is
equivalent to $D_{1,2}\to-D_{1,2}$ and leaves the physics unchanged.

The DM contribution to the Bertaut effective fields
$\bfH^{X,\rm site}_{\rm DM} = -\tfrac14\,\partial(H_{\rm DM}/N_c)/\partial\bfX$
is then
\begin{equation}
\begin{aligned}
\bfH^{F,{\rm site}}_{\rm DM} &= (\,2D_1G_z - 2D_2A_y,\;-2D_2A_x,\;2D_1G_x\,),\\
\bfH^{G,{\rm site}}_{\rm DM} &= (\,2D_1F_z - 2D_2C_y,\;2D_2C_x,\;-2D_1F_x\,),\\
\bfH^{C,{\rm site}}_{\rm DM} &= (-2D_1A_z + 2D_2G_y,\;-2D_2G_x,\;2D_1A_x),\\
\bfH^{A,{\rm site}}_{\rm DM} &= (\,2D_1C_z - 2D_2F_y,\;2D_2F_x,\;-2D_1C_x).
\end{aligned}
\label{eq:HDMexplicit}
\end{equation}

\subsection{Ground state in the \texorpdfstring{$\Gamma_2$}{Gamma2} phase}
\label{sec:Gamma2GS}

The $\Gamma_2$ phase is defined by the order parameter pattern
\begin{equation}
\bfG = G_z\hat z,\qquad
\bfF = F_x\hat x,\qquad
\bfC = C_y\hat y,\qquad
\bfA \approx 0,
\label{eq:Gamma2OP}
\end{equation}
as the phenomenological identification of the ordered phase. However,
for the Hamiltonian written in Eqs.~\eqref{eq:HFe} and
\eqref{eq:HDMBertaut}, the DM sector does \emph{not} generate a direct
$D_1F_xG_z$ or $D_2C_yF_x$ energy inside the reduced
$(\bfF,\bfG,\bfC,\bfA)$ manifold. Instead it couples the nominal
$\Gamma_2$ variables to the exchange-stiff partner modes
$A_z$, $C_x$, $G_y$, and $A_x$, so a consistent static reduction
requires either integrating out those auxiliary modes or adopting a
different Bertaut-label convention. Accordingly we replace the naive
canting estimate by the weaker, but correct, statement
\begin{equation}
\begin{aligned}
F_x/G_z,\; C_y/G_z
&\text{ must be obtained from a full static minimization of }H_{\rm Fe}\\
&\text{(or from an explicit elimination of }A_z,\,C_x,\,G_y,\,A_x\text{),}
\end{aligned}
\label{eq:cantingangles}
\end{equation}
not from the truncated ansatz $F_x\sim D_1G_z/H_E^{\rm exch}$ alone.

\section{Tm--Fe coupling}
\label{sec:TmFe}

\subsection{Most general bilinear coupling}

Before imposing time reversal, the most general bilinear Tm--Fe coupling
allowed by $Pbnm$ would be a sum over all eight qutrit generators.
Because the Hamiltonian itself must be $T$-even while the Fe spin is
$T$-odd, only the $T$-odd Tm generators
$\lambda^2,\lambda^5,\lambda^7$ [Eq.~\eqref{eq:lambdaclass}] can appear
linearly. The most general bilinear Tm--Fe coupling consistent with both
$Pbnm$ and time reversal is therefore
\begin{equation}
H_{\rm Tm-Fe}^{\rm lin} = \sum_{\mathbf n,\mathbf n'}\sum_{j,i}\sum_{a\in\{2,5,7\}}\sum_{b=x,y,z}
\chi^{ab}_{ji}(\mathbf n - \mathbf n')\,\lambda^a_{\mathbf n,j}\,S^b_{\mathbf n',i},
\label{eq:Hlingen}
\end{equation}
where $\lambda^a_{\mathbf n,j}$ is the $a$th Gell-Mann generator on
Tm sublattice $j$ in cell $\mathbf n$, $S^b_{\mathbf n',i}$ is the $b$th
spin component on Fe sublattice $i$ in cell $\mathbf n'$, and
$\chi^{ab}_{ji}(\mathbf R)$ is a bond-resolved coupling tensor. The space
group $Pbnm$ constrains $\chi$ via the symmetry-orbit relations of
Sec.~\ref{sec:bondsymmetry}; time reversal excludes any coefficient that
would multiply the $T$-even generators $\lambda^{1,3,4,6,8}$ linearly.

For the on-site $\mathbf q\!=\!0$ projection that controls the
Bertaut-mode coupling, the relevant combinations are
\begin{equation}
\bar\chi^{ab} = \sum_{\mathbf R, j, i}\chi^{ab}_{ji}(\mathbf R),
\label{eq:chibardef}
\end{equation}
i.e., the uniform sum over all bond shells. Because Tm sits at a
non-centrosymmetric site (4c), and the inversion image of the bond from
Tm$_j$ to Fe$_i$ at offset $\mathbf R$ goes from Tm$_{(14)(23)j}$ to
Fe$_i$ at offset $-\mathbf R$. Because the Fe spin is an axial vector,
it is \emph{unchanged} by inversion; the inversion-pair cancellation at
$\mathbf q=0$ therefore acts only through the mirror parity of the Tm
operator, exactly as in App.~\ref{app:bondcancel}. The inversion-pair
sums in Eq.~\eqref{eq:chibardef} cancel
several entries of $\bar\chi^{ab}$.

\subsection{Surviving \texorpdfstring{$\mathbf q\!=\!0$}{q=0} channels}

After time reversal and the inversion-pair cancellation, the surviving
$\mathbf q=0$ linear couplings are:
\begin{itemize}[leftmargin=*]
\item $\bar\chi^{2z}\neq 0$: couples $\lambda^2$ ($A_1^-$, mirror-even,
$T$-odd; cf.\ Eq.~\eqref{eq:lambdaclass}) to $G_z$ (the $\Gamma_2$ N\'eel order parameter). \emph{This is
the dominant linear Fe--Tm exchange vertex used in Geometries~II and~III.}
\item All off-diagonal mirror-odd couplings $\bar\chi^{5x}, \bar\chi^{7x}$,
$\bar\chi^{5y}, \bar\chi^{7y}$ \emph{cancel
on the $\mathbf q=0$ projection} due to the bond-pair structure (App.~\ref{app:bondcancel}).
\item The $T$-even channels $\lambda^{1,3,4,6,8}$ therefore do not appear
in the microscopic linear exchange at $\mathbf q=0$; they enter only
through the second-order macroscopic back-action discussed in
Sec.~\ref{sec:vanvleck}.
\end{itemize}
This is the key statement that motivates the literature-aligned
``anisotropy-quench'' route (Sec.~\ref{sec:vanvleck}) for the apparent
mirror-odd-to-$F$ coupling: it does not exist as a $\mathbf q=0$ linear
exchange, but it does emerge at second order via the symmetric
($\lambda^1, \lambda^3, \lambda^8$) Tm density coordinates.

\subsection{Bare CEF field on the Tm qutrit}
\label{sec:bareCEFfield}

The on-site Tm Hamiltonian in the truncated three-level basis is
$H_{\rm Tm}^{(0)} = \sum_\alpha\epsilon_\alpha|E_\alpha\rangle\langle E_\alpha|$,
which we expand in Gell-Mann basis as
\begin{equation}
H_{\rm Tm}^{(0)} = \bar\epsilon\,\mathds{1} + h^{(0)}_3\lambda^3 + h^{(0)}_8\lambda^8,
\label{eq:HTm0}
\end{equation}
with
\begin{equation}
\begin{aligned}
\bar\epsilon &= (\epsilon_1+\epsilon_2+\epsilon_3)/3,\\
h^{(0)}_3 &= (\epsilon_1-\epsilon_2)/2,\\
h^{(0)}_8 &= (\epsilon_1+\epsilon_2-2\epsilon_3)/(2\sqrt{3}).
\end{aligned}
\label{eq:hbareCEF}
\end{equation}
Numerically, $h^{(0)}_3 = -0.96$\,meV, $h^{(0)}_8 = -3.97$\,meV.

The \emph{first-order static energy shift} from
$\bar\chi^{2z}G_z\lambda^2$ vanishes in the $\Gamma_2$ ground state
because $\langle\lambda^2\rangle_0 = 0$ (the ground state $|E_1\rangle$
is real in the $T$-symmetric gauge). The operator itself is of course
off-diagonal in the qutrit basis and mixes $|E_1\rangle$ with
$|E_2\rangle$, so it contributes to static renormalization only from
second order onward. There is
no time-reversal-allowed linear $G_z\lambda^{3,8}$ term in
$H_{\rm Tm-Fe}^{\rm lin}$. Any static renormalization of the
$\epsilon_\alpha$ from the Fe background is therefore absorbed into the
effective CEF parameters or into the second-order Van Vleck functional,
not written as a first-order diagonal exchange term.

\subsection{Anisotropy modulation: Van Vleck back-action}
\label{sec:vanvleck}

The dominant Tm-induced contribution to the Fe anisotropy at
$\mathbf q=0$ comes from second-order perturbation in the linear Tm--Fe
coupling. The object introduced below is not the most general microscopic
Tm--Fe Hamiltonian; it is the most general $\mathbf q=0$ macroscopic
effective back-action operator on the \emph{retained} qutrit sector, to
quadratic order in the Fe Bertaut amplitudes, generated by second-order
virtual processes in the microscopic Tm--Fe coupling. For the Tm ground
state, the same mechanism reduces to the familiar Van Vleck
energy correction. At fixed $\bfG=
G_z\hat z, \bfF=F_x\hat x$ is
\begin{equation}
\delta E^{\rm VV}_1 = -\sum_{\alpha\neq 1}\frac{|\langle E_\alpha|H_{\rm Tm-Fe}^{\rm lin}|E_1\rangle|^2}
{\epsilon_\alpha - \epsilon_1},
\label{eq:VVE1}
\end{equation}
which depends on $G_z, F_x$ through the projected operators in
$H_{\rm Tm-Fe}^{\rm lin}$. This remains consistent even though
$H_{\rm Tm-Fe}^{\rm lin}$ contains only the $T$-odd generators
$\lambda^{2,5,7}$ linearly: at second order, products of two $T$-odd
operators are $T$-even, and the SU(3) algebra closes onto the even sector,
with $\{\lambda^2,\lambda^5\}=\lambda^6$,
$\{\lambda^2,\lambda^7\}=-\lambda^4$,
$\{\lambda^5,\lambda^7\}=\lambda^1$, while
$(\lambda^2)^2,(\lambda^5)^2,(\lambda^7)^2$ contain the diagonal
generators $\lambda^{3,8}$ (plus the identity). Expanding to quadratic order in the Bertaut
amplitudes and decomposing in the mirror sector classification of
Eq.~\eqref{eq:lambdaclass}, the resulting $\mathbf q=0$ macroscopic
back-action functional is
\begin{equation}
\begin{aligned}
H_{\rm aniso\text{-}Tm} = &-\sum_{\mathbf n,j}\Bigl[
u_{1,j}\lambda^1_{\mathbf n,j} + u_{3,j}\lambda^3_{\mathbf n,j}\\
&\qquad\quad{} + u_{8,j}\lambda^8_{\mathbf n,j}
\Bigr]\,\mathcal I^{\rm Fe}_{A_1^+}(\mathbf n)\\
&- \sum_{\mathbf n,j}\Bigl[
v_{4,j}\lambda^4_{\mathbf n,j} + v_{6,j}\lambda^6_{\mathbf n,j}
\Bigr]\,\mathcal I^{\rm Fe}_{A_2^+}(\mathbf n),
\end{aligned}
\label{eq:Hanisogen}
\end{equation}
where the most general $A_1^+$ (mirror-even) Fe bilinears are
\begin{equation}
\mathcal I^{\rm Fe}_{A_1^+} = \alpha_1\,G_z^2 + \alpha_2\,F_x^2 + \alpha_3\,C_y^2
+ \alpha_4\,(G_x^2 - G_y^2) + \alpha_5\,F_x C_y + \cdots
\label{eq:IA1}
\end{equation}
and the $A_2^+$ (mirror-odd) ones are
\begin{equation}
\mathcal I^{\rm Fe}_{A_2^+} = \beta_1\,F_x G_z + \beta_2\,G_z C_y + \cdots
\label{eq:IA2}
\end{equation}
with $\alpha_i, \beta_i$ phenomenological $\mathbf q=0$ coefficients
inherited from the microscopic $\chi^{ab}_{ji}$ tensor and the Tm CEF
matrix elements. We keep at least the two
leading combinations $\alpha_1 G_z^2$ and $\alpha_2 F_x^2$ as independent;
collapsing $\mathcal I^{\rm Fe}_{A_1^+}$ to a single combination such as
$G_z^2 - F_x^2$ would be a one-parameter truncation, not a symmetry
constraint.

The $A_2^+$ on-site coupling $v_4 \lambda^4 F_x G_z$ vanishes in the
inversion-pair cancellation analysis (App.~\ref{app:bondcancel}) at
$\mathbf q=0$, leaving the $A_1^+$ Van Vleck channels as the dominant
back-action mechanism.

\subsection{Coherence contribution to Van Vleck}
The Van Vleck back-action Eq.~\eqref{eq:Hanisogen} carries both
\emph{population} and \emph{coherence} contributions in the SU(3) Bloch
language. Writing $\rho_{\rm Tm} = \sum_a \rho_a\lambda^a/2 + \mathds{1}/3$ and
extracting the coefficient of each $A_1^+$ generator,
\begin{equation}
\begin{aligned}
u_{1,j}\langle\lambda^1_{\mathbf n,j}\rangle &\sim \mathrm{Re}[\rho^{(j)}_{12}(\mathbf n)],\\
u_{3,j}\langle\lambda^3_{\mathbf n,j}\rangle &\sim p_1^{(j)}(\mathbf n)-p_2^{(j)}(\mathbf n),\\
u_{8,j}\langle\lambda^8_{\mathbf n,j}\rangle &\sim \tfrac{1}{\sqrt 3}(p_1+p_2-2 p_3).
\end{aligned}
\label{eq:VVcoherence}
\end{equation}
The coherence term $\mathrm{Re}[\rho_{12}(t)]$ oscillates at $\omega_{12}$
during the pump-delay coordinate and therefore modulates the Fe
anisotropy through the Van Vleck channel in ideal Geometry~I of the
low-temperature $\Gamma_2$ phase.

% =====================================================================
\section{Drive Hamiltonian}
\label{sec:drive}

The single-cycle THz pulse couples to both Fe and Tm sectors. The Fe
Zeeman vertex (drives the qFM and qAFM modes through their magnetic-dipole
matrix elements) and the Tm magnetic-dipole vertex (drives the
$\omega_{12,13,23}$ CF coherences through $\Jeff$):
\begin{equation}
H_{\rm drive}^{\rm m}(t) = -g_e\mu_B\sum_{\mathbf n,i}\mathbf h(t)\cdot\bfS_{\mathbf n,i}
- g_J\mu_B\sum_{\mathbf n,j}\mathbf h(t)\cdot\Jeff_{\mathbf n,j}.
\label{eq:Hdrivem}
\end{equation}
The Tm electric-dipole vertex (drives the same CF coherences through
$\deff$):
\begin{equation}
H_{\rm drive}^{\rm e}(t) = -\sum_{\mathbf n,j}\mathbf e(t)\cdot\deff_{\mathbf n,j}.
\label{eq:Hdrivee}
\end{equation}
For THz propagation along $\hat b$, the two relevant field assignments are
\begin{equation}
\begin{aligned}
&H\!\parallel\! a,\, E\!\parallel\! c:\ h_x(t)\neq 0,\ e_z(t)\neq 0\\
&\qquad\Longrightarrow\ \text{couples }\lambda^{5,7}\text{ and }\lambda^{4,6};\\
&H\!\parallel\! c,\, E\!\parallel\! a:\ h_z(t)\neq 0,\ e_x(t)\neq 0\\
&\qquad\Longrightarrow\ \text{couples }\lambda^2\text{ and }\lambda^1.
\end{aligned}
\label{eq:geometries}
\end{equation}

\subsection{Drive Hamiltonian in Bertaut form}
\label{sec:driveBertaut}

Substituting the local-frame transport laws Eq.~\eqref{eq:Smutransform}
(spin/$\Jeff$, axial) and Eq.~\eqref{eq:dlocal} ($\deff$, polar) into
Eqs.~\eqref{eq:Hdrivem}--\eqref{eq:Hdrivee} converts each four-sublattice
sum into the Bertaut combination $\Lambda^a_{X(a)}$ matched to
the $\sigma$-vector of the global field component (Eqs.~\eqref{eq:uniformsel}
and \eqref{eq:uniformselE}):
\begin{widetext}
\begin{equation}
\boxed{\;
\begin{aligned}
H_{\rm drive}^{\rm m,Tm}(t) &= -4 g_J\mu_B\sum_{\mathbf n}\Bigl[
h_z\,M^{(z)}_{12}\,\Lambda^2_G
+ h_x\bigl(d\,\Lambda^5_C - b\,\Lambda^7_C\bigr)
+ h_y\bigl(-c\,\Lambda^5_A + a\,\Lambda^7_A\bigr)
\Bigr],\\[2pt]
H_{\rm drive}^{\rm e,Tm}(t) &= -4\sum_{\mathbf n}\Bigl[
e_z\bigl(p_{z4}\,\Lambda^4_G + p_{z6}\,\Lambda^6_G\bigr)
+ e_x\,p_x\,\Lambda^1_C
+ e_y\,p_y\,\Lambda^1_A
\Bigr],\\[2pt]
H_{\rm drive}^{\rm m,Fe}(t) &= -4 g_e\mu_B\sum_{\mathbf n}\Bigl[
h_x\,F_x + h_y\,F_y + h_z\,F_z
\Bigr],
\end{aligned}\;}
\label{eq:HdriveBertaut}
\end{equation}
\end{widetext}
with $(a,b,c,d)$ from Eq.~\eqref{eq:Mnumeric} and the Bertaut
combinations $\Lambda^a_X$ defined in Eq.~\eqref{eq:LambdaXdef}. Three
features of Eq.~\eqref{eq:HdriveBertaut} are worth pointing out:
\begin{itemize}[leftmargin=*]
\item Every Tm coupling sits on the Bertaut combination whose
$\sigma$-vector matches the field's global axis (Klein-four selection
rule). \emph{No uniform Tm-sector coupling carries $\sigma_F$}, so the
homogeneous $\Lambda^a_F$ combinations cannot be excited by the bare
linear drive.
\item Fe couples to the uniform $F_a$ combination only (per
Eq.~\eqref{eq:HFelocal} the global Fe Zeeman acts as
$-g_e\mu_B\sum_a h_a F_a$ once the local-frame sum is performed,
because each $h_a S^a_\mu = h_a \eta^a_\mu \tilde S^a_\mu$ projects
onto the $\sigma_F$-symmetric piece of $\tilde S^a_\mu$). At
$\mathbf q=0$ this is exactly the $F_x, F_y, F_z$ ferromagnetic dipole.
\item The ratio of Tm-electric to Tm-magnetic coupling at fixed wave
impedance $|\bfE|/|\bfH|=Z_0$ is $|p|/(g_J\mu_B Z_0 |\Jeff|)$, controlling
which channel dominates each pathway.
\end{itemize}
The two propagation geometries Eq.~\eqref{eq:geometries} now read off
directly:
\begin{itemize}[leftmargin=*]
\item $H\!\parallel\!\hat a, E\!\parallel\!\hat c$:
the drive populates the four \emph{distinct} Bertaut combinations
$\Lambda^5_C, \Lambda^7_C$ (axial $h_x$) and $\Lambda^4_G, \Lambda^6_G$
(polar $e_z$). All four sit in $\Gamma_3+\Gamma_4$ (off the $\Gamma_2$
ground state), hence drive the $\omega_{13},\omega_{23}$ Tm coherences.
\item $H\!\parallel\!\hat c, E\!\parallel\!\hat a$:
the drive populates $\Lambda^2_G$ (axial $h_z$, the unique $\Gamma_2$
admixture --- this is the qAFM channel) and $\Lambda^1_C$ (polar $e_x$,
in $\Gamma_4$). The $\Lambda^2_G$ coupling is the Tm partner of the Fe
$G_z$ N\'eel coordinate, which is why this geometry mixes Tm and Fe
through $\bar\chi^{2z}$.
\end{itemize}

% =====================================================================
\section{Microscopic Hamiltonian and \texorpdfstring{$\mathbf q=0$}{q=0} Reduction}
\label{sec:Htotal}

At this point the microscopic Hamiltonian derivation is complete. The
bond-resolved time-dependent Hamiltonian is
\begin{equation}
\boxed{\;
H_{\rm micro}(t) = H_{\rm Fe} + H_{\rm CEF}^{\rm Tm,\,4\text{-}sl}
+ H_{\rm Tm\text{-}Fe}^{\rm lin}
+ H_{\rm drive}^{\rm m}(t) + H_{\rm drive}^{\rm e}(t)\;,
}
\label{eq:Hmicro}
\end{equation}
with the Fe four-sublattice Hamiltonian $H_{\rm Fe}$ from
Eq.~\eqref{eq:HFe}, the four-sublattice Tm crystal-field Hamiltonian
$H_{\rm CEF}^{\rm Tm,\,4\text{-}sl}$ from Eq.~\eqref{eq:HCEF4sl}, the
time-reversal-even linear Fe--Tm coupling $H_{\rm Tm\text{-}Fe}^{\rm lin}$
from Eq.~\eqref{eq:Hlingen}, and the THz magnetic/electric drive terms
$H_{\rm drive}^{\rm m,e}(t)$ from Eqs.~\eqref{eq:Hdrivem}
and \eqref{eq:Hdrivee}.

For the uniform low-temperature spectroscopy problem studied here, we then
integrate out the virtual Tm processes at second order and keep the
resulting $\mathbf q=0$ Van Vleck back-action functional
$H_{\rm aniso\text{-}Tm}$. The working uniform-sector Hamiltonian used in the
rest of the manuscript is therefore
\begin{equation}
\boxed{\;
H_{\rm total}(t) = H_{\rm micro}(t) + H_{\rm aniso\text{-}Tm}\;,
}
\label{eq:Htotal}
\end{equation}
where $H_{\rm aniso\text{-}Tm}$ is the effective $\mathbf q=0$ second-order
back-action functional of Eq.~\eqref{eq:Hanisogen}. Equation
\eqref{eq:Hmicro} is the microscopic all-bond Hamiltonian; Eq.~\eqref{eq:Htotal}
is the derived working Hamiltonian for the uniform spectroscopy problem.

For the uniform $\mathbf q=0$ response analyzed below, the drive sector is
further reduced to the Bertaut form Eq.~\eqref{eq:HdriveBertaut}, and the
static Tm part is projected to the three-level qutrit Hamiltonian
$H_{\rm Tm}^{(0)}$ of Eq.~\eqref{eq:HTm0}. The subsequent low-temperature
$\Gamma_2$ analysis then linearizes Eq.~\eqref{eq:Htotal} around the
background $G_z+F_x$, keeps the low-frequency Fe block
$M^{\Gamma_2}_{4\times 4}$, and combines it with the Tm SU(3) block and
the surviving $\mathbf q=0$ coupling channels. Dissipation is added only
at the equation-of-motion level through the Lindblad superoperator and is
not part of $H_{\rm total}(t)$.

% =====================================================================
\section{2DCS framework and experimental geometries}
\label{sec:2dcsframework}

\subsection{Experimental nonlinear field, normal coordinates, and geometry channels}

Experimentally the nonlinear field is isolated by subtracting the
single-pulse controls from the full two-pulse response,
\begin{equation}
E_{\rm sig}^{\rm NL}(t_3,t_2,t_1)
= E_{AB}(t_3,t_2,t_1) - E_A(t_3,t_2,t_1) - E_B(t_3,t_2,t_1).
\label{eq:ENL}
\end{equation}
Equation~\eqref{eq:ENL} is the directly measured time-domain quantity;
the response-function formalism below is its perturbative evaluation in
terms of the Hamiltonian Eq.~\eqref{eq:Htotal} and the emission operator
Eq.~\eqref{eq:emissionslab}.

For a standard magnetic-2DCS presentation it is convenient to use the
collective normal coordinates of the low-temperature $\Gamma_2$ phase,
\begin{equation}
\begin{aligned}
Q_{\rm qFM}(t)&\equiv q_{\rm qFM}(t), & \Omega_{\rm qFM}&=\omega_{\rm qFM},\\
Q_{\rm qAFM}(t)&\equiv q_{\rm qAFM}(t), & \Omega_{\rm qAFM}&=\omega_{\rm qAFM},\\
Q_{12}(t)&\equiv 2\,\Re\rho_{12}(t), & \Omega_{12}&=\omega_{12},\\
Q_{13}(t)&\equiv 2\,\Re\rho_{13}(t), & \Omega_{13}&=\omega_{13},\\
Q_{23}(t)&\equiv 2\,\Re\rho_{23}(t), & \Omega_{23}&=\omega_{23},
\end{aligned}
\label{eq:Qcoords}
\end{equation}
where $Q_{13}$ and $Q_{23}$ mean the Bertaut-projected $13$ and $23$
coherences selected by the active geometry, i.e. the appropriate member
of the set written explicitly in Eq.~\eqref{eq:dynBertaut}. At linear
order these coordinates obey the standard coupled-oscillator form
\begin{equation}
\ddot Q_\mu + 2\Gamma_\mu \dot Q_\mu + \Omega_\mu^2 Q_\mu
+ \sum_\nu \kappa_{\mu\nu} Q_\nu = F_\mu(t),
\label{eq:Qlinear}
\end{equation}
with the retained Fe--Tm mixing contained in
$\kappa_{12,{\rm qAFM}}\sim\kappa_{{\rm qAFM},12}\sim
\bar\chi^{2z}\sin\phi_{\rm qAFM}$ and the remaining off-diagonal
entries fixed by the symmetry analysis of Secs.~\ref{sec:driveBertaut}
and \ref{sec:linEOM}.

The geometry-resolved drive and detection channels are then:
\begin{itemize}[leftmargin=*]
\item Geometry I ($H\parallel a, E\parallel c$, VV): the pump drives the
mirror-odd Tm sector $\Lambda_C^5,\Lambda_C^7,\Lambda_G^4,\Lambda_G^6$;
VV detection projects onto the same mirror-odd channel, so the directly
bright linear coordinates are $Q_{13}$ and $Q_{23}$.
\item Geometry II ($H\parallel a, E\parallel c$, VH): the pump is
identical to Geometry I, but VH detection projects onto the mirror-even
channel $\Lambda_G^2,\Lambda_C^1$; therefore $Q_{12}$ is read out only
after Fe-to-Tm conversion through the qAFM admixture.
\item Geometry III ($H\parallel c, E\parallel a$, VH): the pump drives
$\Lambda_G^2$, $\Lambda_C^1$, and the Fe $F_z$ Zeeman channel; the
directly addressed linear coordinates are $Q_{12}$ and $Q_{\rm qFM}$,
while VH detection reads the orthogonal qAFM emission channel.
\end{itemize}
Each observed peak must therefore be explained twice: once as an allowed
Liouville pathway and once as a symmetry-allowed coupling among the
normal coordinates $Q_\mu$.

\subsection{The third-order response function}

The 2DCS signal in the rephasing/non-rephasing direction is
\begin{equation}
\begin{aligned}
S(\omega_t, t_2, \omega_\tau) =\;& \int_0^\infty\!dt_3\,e^{i\omega_t t_3}\!\int_0^\infty\!dt_1\,e^{\pm i\omega_\tau t_1}\\
&\times R^{(3)}(t_3, t_2, t_1),
\end{aligned}
\label{eq:S2DCS}
\end{equation}
where the $\pm$ selects rephasing ($-$) or non-rephasing ($+$) and
\begin{equation}
\begin{aligned}
R^{(3)}(t_3, t_2, t_1) =\;& \Bigl(\tfrac{i}{\hbar}\Bigr)^3 \mathrm{Tr}\Bigl\{
\hat O_{\rm det}\,\mathcal G(t_3)\,\mathcal V_3\,\mathcal G(t_2)\\
&\times\mathcal V_2\,\mathcal G(t_1)\,\mathcal V_1\,\rho_{\rm eq}
\Bigr\},
\end{aligned}
\label{eq:R3}
\end{equation}
with $\mathcal V_i = [\hat V_i,\,\cdot\,]$ the dipole-coupling
superoperator at interaction $i$, $\mathcal G(t) = e^{-i\mathcal Lt}$ the
free-evolution superoperator (with damping), and $\hat O_{\rm det}$ the
dipole detected in emission. Each Liouville pathway is one of the eight
double-sided Feynman diagrams that arise from the choice of
left/right action of $\mathcal V_i$ on the density matrix.

% =====================================================================
\section{SU(3) Bloch dynamics with CPTP dissipation}
\label{sec:bloch}

\subsection{Coherent evolution}
The Tm density matrix on each sublattice obeys
\begin{equation}
\dot\rho^{(j)}_{\mathbf n} = -i[H^{(j)}_{\mathbf n}, \rho^{(j)}_{\mathbf n}] + \mathcal L^{\rm diss}[\rho^{(j)}_{\mathbf n}],
\label{eq:rhoEOM}
\end{equation}
where $H^{(j)}_{\mathbf n}$ is the local Tm projection of the complete
Hamiltonian Eq.~\eqref{eq:Htotal} (bare CEF + linear Fe coupling + drive
+ Van Vleck back-action). In the SU(3) Bloch
representation $\rho = \mathds{1}/3 + \half\sum_a n^a\lambda^a$ with
$n^a = \langle\lambda^a\rangle$, the coherent equation reads
\begin{equation}
\dot n^a = \sum_{b,c} f^{abc}\,h^b\,n^c,\qquad
H = h_0\mathds{1} + \half\sum_a h^a\lambda^a.
\label{eq:nEOM}
\end{equation}
With this convention, the bare CEF field is $h^3 = 2 h^{(0)}_3$, $h^8 =
2 h^{(0)}_8$ (factor of 2 from the $\half$ in the convention), giving the
correct precession frequencies $\omega_{12} = 2|h^{(0)}_3|/\hbar
= |\epsilon_2-\epsilon_1|/\hbar$, etc. \fix{The factor-of-2 chain is
collected in App.~\ref{app:conventions} and propagated consistently.}

\subsection{Lindblad dissipator}
A naive pure-dephasing dissipator $\dot n^a\supset -\Gamma_a n^a$ is not
in general CPTP-compliant on a qutrit because it can drive $\rho$ outside
the positive cone. We therefore use a secular Lindblad form in the CF
eigenbasis with explicit relaxation jumps and diagonal dephasing
projectors:
\begin{equation}
\begin{aligned}
\mathcal L^{\rm diss}[\rho] = &\sum_{\alpha>\beta}\Gamma_{\alpha\to\beta}\Bigl[
L_{\beta\alpha}\rho L_{\beta\alpha}^\dagger\\
&\qquad\qquad{} - \half\{L_{\beta\alpha}^\dagger L_{\beta\alpha}, \rho\}
\Bigr]\\
&{} + \sum_{\alpha=1}^3\gamma^\phi_{\alpha}\Bigl[
P_{\alpha}\rho P_{\alpha} - \half\{P_{\alpha}, \rho\}
\Bigr],
\end{aligned}
\label{eq:Lindblad}
\end{equation}
with $L_{\beta\alpha} = |E_\beta\rangle\langle E_\alpha|$ for inelastic
relaxation and $P_{\alpha}=|E_\alpha\rangle\langle E_\alpha|$ for pure
dephasing. The relaxation rates obey detailed balance
$\Gamma_{\alpha\to\beta}/\Gamma_{\beta\to\alpha} =
e^{-(\epsilon_\alpha-\epsilon_\beta)/k_BT}$, while the coherence
$|E_\alpha\rangle\langle E_\beta|$ acquires the additional dephasing rate
$(\gamma^\phi_{\alpha}+\gamma^\phi_{\beta})/2$. CPTP is automatic for any
$\Gamma_{\alpha\to\beta}, \gamma^\phi_{\alpha}\geq 0$.

\subsection{Linearization around the \texorpdfstring{$|E_1\rangle$}{|E1>} vacuum}
\label{sec:Tmlin}

At $T\ll(\epsilon_2-\epsilon_1)/k_B \approx 22$\,K the Tm population is
concentrated in $|E_1\rangle$, $\rho \approx |E_1\rangle\langle E_1|$,
and the equilibrium SU(3) Bloch vector is
\begin{equation}
n^{3,\rm eq} = +1,\qquad n^{8,\rm eq} = +1/\sqrt 3,\qquad
n^{a,\rm eq} = 0\text{ otherwise.}
\label{eq:Tmeq}
\end{equation}
Linearizing Eq.~\eqref{eq:nEOM} around Eq.~\eqref{eq:Tmeq} with the bare
CEF field of Eq.~\eqref{eq:HTm0} block-diagonalizes into three $2\times 2$
oscillator blocks:
\begin{equation}
\begin{aligned}
(\delta n^1, \delta n^2):&\quad \omega_{12} = |\epsilon_2-\epsilon_1|/\hbar,\\
(\delta n^4, \delta n^5):&\quad \omega_{13} = |\epsilon_3-\epsilon_1|/\hbar,\\
(\delta n^6, \delta n^7):&\quad \omega_{23} = |\epsilon_3-\epsilon_2|/\hbar,
\end{aligned}
\label{eq:omegabare}
\end{equation}
plus the two zero-frequency population channels $(\delta n^3, \delta n^8)$.

\paragraph*{Validity regime.} Eq.~\eqref{eq:Tmeq} is well-justified for
$T\lesssim 20$\,K. At the spin-reorientation temperatures
$T_{\rm SR1,2}\approx 81$--$91$\,K the populations are
$p_1\approx 0.50$, $p_2\approx 0.40$, $p_3\approx 0.10$ from
Boltzmann weighting, and the linearization around $|E_1\rangle$ is
\emph{not valid}; the model must be re-linearized around a thermally
mixed reference state, with the equilibrium $n^a$ shifted accordingly.
We restrict the present analysis to $T\lesssim 20$\,K where the standard
linearization holds.

% =====================================================================
\section{Effective Fe magnon sector around the \texorpdfstring{$\Gamma_2$}{Gamma2} ground state}
\label{sec:magnons}

\subsection{Linearization}
We write
\begin{equation}
\begin{aligned}
\bfF &= F_x\hat x + \delta\bfF, & \bfG &= G_z\hat z + \delta\bfG,\\
\bfC &= C_y\hat y + \delta\bfC, & \bfA &= \delta\bfA,
\end{aligned}
\label{eq:linearizationFe}
\end{equation}
with the spin-length constraint $|\bfS_\mu|=S$ enforced through
$\bfF\cdot\delta\bfF + \bfG\cdot\delta\bfG + \bfC\cdot\delta\bfC = 0$
order by order. To leading order this kills three of the twelve
components $(\delta\bfF,\delta\bfG,\delta\bfC,\delta\bfA)$, leaving
nine independent linearized degrees of freedom. Of these, four split
into the high-frequency optical branches (involving $\delta\bfA$,
$\delta C_z$, etc.) at the exchange scale, and the remaining
\emph{five} couple via the DM and anisotropy terms.
\fix{With the corrected DM Bertaut Hamiltonian Eq.~\eqref{eq:HDMBertaut},
the canonical Treves--Bertaut Lifshitz invariant $D_1 F_xG_z$ \emph{is}
carried inside the reduced $(\bfF,\bfG,\bfC,\bfA)$ manifold, so the
low-frequency block closes on the four-component subspace below to
leading order in $D_{1,2}/J_1$; the exchange-stiff $\Gamma_4$ partners
($\delta A_x,\delta A_z,\delta C_x,\delta C_z$) only renormalize the
$\kappa$ coefficients of Eq.~\eqref{eq:M4x4}.} For spectroscopy we
introduce the \emph{effective} low-frequency coordinate vector
\begin{equation}
\mathbf v \equiv (\delta F_y,\;\delta G_y;\;\delta F_z,\;\delta G_x)^T
\label{eq:lowfreqsubspace}
\end{equation}
whose last two entries represent the dominant projection of the qFM-like
resonance (the soft $ac$-plane rotation), while the first two represent
the dominant projection of the qAFM-like hard-axis fluctuation. Any
additional components of the exact eigenvectors are absorbed into the
renormalized coefficients of the reduced model or into small
phenomenological mixing angles introduced later.

\subsection{Explicit \texorpdfstring{$4\times 4$}{4x4} linearized matrix}
\label{sec:M4x4}

In the effective basis Eq.~\eqref{eq:lowfreqsubspace}, we therefore use
an effective reduced Fe magnon block,
\begin{equation}
\dot{\mathbf v} = M^{\Gamma_2}_{4\times 4}\,\mathbf v,
\qquad
M^{\Gamma_2}_{4\times 4} \equiv
\begin{pmatrix}
0 & \kappa_{\rm qAFM}^{(1)} & 0 & 0 \\
-\kappa_{\rm qAFM}^{(2)} & 0 & 0 & 0 \\
0 & 0 & 0 & \kappa_{\rm qFM}^{(1)} \\
0 & 0 & -\kappa_{\rm qFM}^{(2)} & 0
\end{pmatrix},
\label{eq:M4x4}
\end{equation}
whose coefficients are the renormalized frequency scales obtained after
eliminating the exchange-stiff partner modes omitted from
Eq.~\eqref{eq:lowfreqsubspace}. \fix{The previously printed
$\alpha_i,\beta_i$ expressions are withdrawn: they are not the
linearization of the corrected Section-IV Hamiltonian.} In practice the
$\kappa$'s are determined either from a numerical linearization of the
full Fe sector or by fitting the two observed Fe resonance frequencies.
The two $2\times 2$ blocks encode, by construction, the dominant qAFM
and qFM resonances used in the later reduced joint model.

\subsection{qFM frequency}

The qFM-like block $(\delta F_z, \delta G_x)$ represents the soft
$ac$-plane rotation of the $\Gamma_2$ state. Its reduced-model
frequency is
\begin{equation}
\omega_{\rm qFM}^2 = \kappa_{\rm qFM}^{(1)}\kappa_{\rm qFM}^{(2)}.
\label{eq:omqFMexact}
\end{equation}
After the exchange-stiff partner modes are eliminated, the familiar
orthoferrite scaling is recovered only at the level of a schematic
two-mode reduction,
\begin{equation}
\omega_{\rm qFM}^2 \propto H_E^{\rm exch}\,H_A^{\rm qFM},
\label{eq:omqFM}
\end{equation}
where $H_E^{\rm exch}$ is the exchange scale and $H_A^{\rm qFM}$ the
effective anisotropy scale controlling the soft $ac$-plane rotation,
with $H_A^{\rm qFM}\to 0$ at $T_{\rm SR1}$. Thus qFM is the soft mode
of the $\Gamma_2\to\Gamma_{24}$ spin reorientation. In what follows we
use $\omega_{\rm qFM}$ as the experimentally calibrated or numerically
linearized low-energy Fe resonance, not as the value of the withdrawn
raw $4\times4$ formulas.

\subsection{qAFM frequency: with the DM gap}
\label{sec:qAFMfreq}

The qAFM-like block $(\delta F_y, \delta G_y)$ represents the harder
$b$-axis Fe resonance in the reduced model, with
\begin{equation}
\omega_{\rm qAFM}^2 = \kappa_{\rm qAFM}^{(1)}\kappa_{\rm qAFM}^{(2)}.
\label{eq:omqAFMexact}
\end{equation}
After the same elimination of stiff partner modes, its gap is governed
schematically by hard-axis anisotropy plus a DM-induced contribution,
\begin{equation}
\omega_{\rm qAFM}^2 \sim H_E^{\rm exch}\!\left(H_A^{\rm qAFM} + c_D\,\frac{H_D^2}{H_E^{\rm exch}}\right),
\label{eq:omqAFM}
\end{equation}
where $H_A^{\rm qAFM}$ is the effective hard-axis anisotropy scale,
$H_D$ the effective DM scale, and $c_D=O(1)$ depends on the precise
reduction convention. The important qualitative point is that the qAFM
resonance remains finite across $T_{\rm SR1}$ because the DM sector
produces a nonzero effective gap even when the bare hard-axis
anisotropy is small. Later sections encode the small projection of this
mode onto the particular $G$-staggered coordinate that couples to the
Tm $\lambda^2$ channel by a phenomenological mixing angle
$\phi_{\rm qAFM}$, rather than reading it directly off the withdrawn raw
block formulas.

\subsection{Spin-reorientation transitions}

The signs of the anisotropy differences $K_c - K_a$, $K_c - K_b$ depend
on temperature through the rare-earth contribution
$K^{\rm Tm}_{ab}(T)$ to the effective Fe single-ion anisotropy
(Sec.~\ref{sec:vanvleck}). In TmFeO$_3$, $K_c - K_a$ changes sign at
$T_{\rm SR1}\approx 81$\,K and $K_c - K_b$ changes sign at
$T_{\rm SR2}\approx 91$\,K, giving:
\begin{itemize}[leftmargin=*]
\item $T < T_{\rm SR1}$: $\Gamma_2$ phase, $\bfG\parallel\hat z$,
soft mode is qFM ($\omega_{\rm qFM}\to 0$ as $T\to T_{\rm SR1}^-$);\
\item $T_{\rm SR1} < T < T_{\rm SR2}$: $\Gamma_{24}$ mixed phase, with
$\bfG$ continuously rotating in the $ac$ plane; both the qFM-like and
qAFM-like modes have intermediate frequencies;\
\item $T > T_{\rm SR2}$: $\Gamma_4$ phase, $\bfG\parallel\hat x$, soft
mode is qFM-like with respect to $\Gamma_4$ ($\omega\to 0$ as $T\to T_{\rm SR2}^+$).
\end{itemize}
\fix{This is the canonical orthoferrite spin-reorientation phenomenology
with two second-order transitions bracketing one continuous rotation},
distinct from a single SR transition picture~\cite{ShapiroAxeBirgeneau1974,Yamaguchi1974,Belov1976}.

% =====================================================================
\section{Coherent emission}
\label{sec:emission}

\fix{In SI units, the radiated electric field at the detector contains
both the electric-dipole polarization $\mathbf P$ and the magnetic-dipole
magnetization $\mathbf M$, related by the impedance of free space:}
\begin{equation}
\mathbf E_{\rm sig}(t,\mathbf r) \propto
\partial_t\mathbf P_\perp(t,\mathbf r) + \frac{Z_0}{c}\,\partial_t\mathbf M_\perp(t,\mathbf r)\times\hat{\mathbf k},
\label{eq:emission}
\end{equation}
where $Z_0 = \sqrt{\mu_0/\epsilon_0}\approx 377\,\Omega$ is the
free-space impedance, and the perpendicular projection is with respect to
the propagation direction $\hat{\mathbf k}$. In SI the prefactor on the
$\mathbf M$ contribution is $Z_0/c = \mu_0$, but writing it as $Z_0/c$
emphasizes that the relative weight of electric- and magnetic-dipole
emission in the detected field is set by the impedance.

The signal is the projection of Eq.~\eqref{eq:emission} onto the detector
polarization $\hat{\mathbf e}_{\rm det}$, summed coherently over the
illuminated volume. For a slab of thickness $L$ (with $L\ll\lambda_{\rm THz}$
in the thin-sample limit) and uniform pump,
\begin{equation}
E_{\rm sig}(t)\propto L\,\hat{\mathbf e}_{\rm det}\cdot\Bigl[
\partial_t\mathbf P(t) + \frac{Z_0}{c}\,\partial_t\mathbf M(t)\times\hat{\mathbf k}
\Bigr],
\label{eq:emissionslab}
\end{equation}
where the cell-averaged polarization and magnetization are
\begin{equation}
\begin{aligned}
\mathbf P &= \frac{1}{V_c}\sum_j \langle\deff_{\mathbf n,j}\rangle,\\
\mathbf M &= \frac{1}{V_c}\Bigl[g_e\mu_B\sum_i\langle\bfS_{\mathbf n,i}\rangle\\
&\qquad\quad{} + g_J\mu_B\sum_j\langle\Jeff_{\mathbf n,j}\rangle\Bigr],
\end{aligned}
\label{eq:PMcell}
\end{equation}
and $V_c = abc$ is the unit-cell volume. \fix{In the local-frame
formulation, $\langle\Jeff_{\mathbf n,j}\rangle$ and
$\langle\deff_{\mathbf n,j}\rangle$ must be rotated back to the global
frame using $R_\mu^{\rm axial}$ and $R_\mu^{\rm polar}$ respectively
before summation} (Sec.~\ref{sec:localframes}).

% =====================================================================
\section{Linearized joint EOM}
\label{sec:linEOM}

\subsection{Joint variable vector}
The joint dynamics combines the Fe magnon block (4 dofs in the
low-frequency subspace, Eq.~\eqref{eq:lowfreqsubspace}) with the Tm SU(3)
Bloch block (6 dofs in the three coherence pairs, Eq.~\eqref{eq:omegabare}):
\begin{equation}
\mathbf w = \bigl(\delta F_y, \delta G_y, \delta F_z, \delta G_x;\;
\delta n^1, \delta n^2, \delta n^4, \delta n^5, \delta n^6, \delta n^7\bigr)^T.
\label{eq:wjoint}
\end{equation}
The two Tm population channels $(\delta n^3, \delta n^8)$ decouple at
zeroth order and enter only via second-order Van Vleck back-action
(Sec.~\ref{sec:vanvleck}).

\subsection{Joint matrix}
The linearized EOM is $\dot{\mathbf w} = M^{\Gamma_2,{\rm joint}}\,\mathbf w$
with $M^{\Gamma_2,{\rm joint}}$ block-structured as
\begin{equation}
M^{\Gamma_2,{\rm joint}} = \begin{pmatrix}
M^{\Gamma_2}_{4\times 4} & M^{\rm Fe\to Tm}_{4\times 6} \\
M^{\rm Tm\to Fe}_{6\times 4} & M^{\rm Tm}_{6\times 6}
\end{pmatrix},
\label{eq:Mjoint}
\end{equation}
with:
\begin{itemize}[leftmargin=*]
\item $M^{\Gamma_2}_{4\times 4}$: the Fe magnon matrix of Eq.~\eqref{eq:M4x4};
\item $M^{\rm Tm}_{6\times 6}$: block-diagonal in the three CF doublets,
with $\omega_{12,13,23}$ on the diagonals;
\item $M^{\rm Fe\to Tm}_{4\times 6}$: after projecting the omitted
$\delta G_z$ coordinate onto the low-frequency qAFM eigenvector, the only
surviving reduced-basis entry is the effective
$\mathrm{qAFM}\to\delta n^2$ coupling with strength
$\propto\bar\chi^{2z}\sin\phi_{\rm qAFM}$;
\item $M^{\rm Tm\to Fe}_{6\times 4}$: the reciprocal
$\delta n^2\to\mathrm{qAFM}$ entry generated by the same projected
exchange vertex.
\end{itemize}
Thus the reduced 10-component model contains no bare $\delta G_z$ row or
column. The Fe--Tm hybridization enters only through the projected
coupling
\begin{equation}
\delta n^2(t)\leftrightarrow \bar\chi^{2z}\sin\phi_{\rm qAFM}\,\mathrm{qAFM}(t)
\label{eq:n2qAFM}
\end{equation}
with $\sin\phi_{\rm qAFM}=\mathcal O(D_1/J_1)\ll 1$. This is the
parametric origin of the $(\mathrm{qAFM}, E_{12})$ peak in Geometry~II.

% =====================================================================
\section{Detailed Liouville-pathway analysis of the six 2DCS peaks}
\label{sec:2dcs}

\subsection{General peak amplitude expression}

For a peak generated by a Liouville pathway $\mathcal P$ that carries a
coherence during $t_1$, an intermediate coherence or population
$\eta_{\mathcal P}$ during $t_2$, and a radiating coherence during $t_3$,
the pathway contribution has the generic form
\begin{widetext}
\begin{equation}
A_{\mathcal P}(\omega_\tau,t_2,\omega_t) = \Bigl(\frac{i}{\hbar}\Bigr)^3\rho_{00}\,
\frac{\langle d|\hat O_{\rm det}|c\rangle\,\mathcal V^{(3)}_{\mathcal P}\,\mathcal V^{(2)}_{\mathcal P}\,\mathcal V^{(1)}_{\mathcal P}}
{\mathcal D_{\tau,\mathcal P}(\pm\omega_\tau)\,\mathcal D_{t,\mathcal P}(\omega_t)}
\,e^{-[\Gamma_{\eta,\mathcal P}+i\Omega_{\eta,\mathcal P}]t_2},
\label{eq:Apeak}
\end{equation}
\end{widetext}
where $\mathcal V^{(i)}_{\mathcal P}$ are the pathway-dependent matrix
elements or effective bilinear vertices selected by the three
interactions, $\mathcal D_{\tau,\mathcal P}$ and
$\mathcal D_{t,\mathcal P}$ are the first- and third-interval coherence
denominators, and $\Omega_{\eta,\mathcal P}=0$ when the intermediate
object is a population grating. The explicit peak amplitudes written below
are specializations of Eq.~\eqref{eq:Apeak} to the six symmetry-allowed
pathways.
There is also a formal population branch at
$(\omega_\tau,\omega_t)=(\omega_{12},0)$: the first interaction creates
$\rho_{12}$, the second converts it into the diagonal SU(3) population
channels $\Lambda_F^{3,8}\sim(n_1,n_2)$ through
Eqs.~\eqref{eq:dotn1}--\eqref{eq:dotn2}, and the third interval then
evolves at zero frequency. We keep this $(E_{12},0)$ pathway in mind as a
real population feature of the Liouville problem, but do not include it in
the radiative peak count because $\hat O_{\rm det}$ is a dipole operator and
the detected field is proportional to $\partial_t\mathbf P$ and
$\partial_t\mathbf M$ [Eqs.~\eqref{eq:emission}--\eqref{eq:emissionslab}];
a static population channel therefore carries no coherent far-field THz
emission at $\omega_t=0$ unless one extends the measurement model to direct
population readout or relaxation-tail detection.
Rephasing pathways use the minus sign in
$\mathcal D_{\tau,\mathcal P}(\pm\omega_\tau)$ and produce diagonal-axis
peaks; non-rephasing pathways use the plus sign and produce
cross-axis peaks. We work out the six peaks of the corrected list
explicitly.

\subsection{Bertaut bookkeeping for every pathway}
\label{sec:LiouvilleBertaut}

Before working out the six peaks individually, we record the
$\sigma$-vector that each vertex carries, and the global Klein-four
selection rule that constrains every double-sided Feynman diagram.
Each interaction vertex $\hat V_i$ in Eq.~\eqref{eq:Apeak} carries a
$\sigma$-vector inherited from the field component and the Bertaut
combination it excites (Eqs.~\eqref{eq:HdriveBertaut}); a coherent
pathway is allowed only if the product of $\sigma$-vectors along the
diagram returns to $\sigma_F$ (the trivial irrep), since the
detected one-cell observable transforms as $\sigma_F$.

\paragraph{$\sigma$-rule for a third-order pathway.}
Writing $\sigma_{V_i}$ for the $\sigma$-vector of vertex $\hat V_i$
and $\sigma_{O_{\rm det}}$ for the detector,
\begin{equation}
\boxed{\;\sigma_{V_1}\!\cdot\!\sigma_{V_2}\!\cdot\!\sigma_{V_3}\!\cdot\!\sigma_{O_{\rm det}}
\stackrel{!}{=}\sigma_F,\;}
\label{eq:sigmaRule}
\end{equation}
where the product is component-wise (Klein-four multiplication, Eq.~\eqref{eq:Klein4}).
This is the four-sublattice generalization of the conventional
$\Gamma_{\rm det}\subset\Gamma_{V_1}\otimes\Gamma_{V_2}\otimes\Gamma_{V_3}$
group-theoretic selection rule; for $Pbnm$ orthoferrites at $\mathbf q=0$
it reduces to the requirement that the four-cycle of $\sigma$-vectors
multiply to $(+,+,+,+)$.

\paragraph{$\sigma$-vectors of the six peaks.}
Reading off Eq.~\eqref{eq:HdriveBertaut} and the matrix-element
catalogue Eq.~\eqref{eq:Mnumeric} (each $\Jeff,\deff$ vertex inherits
the $\sigma$ of its field index; each Fe Bertaut excitation
$F_a, G_a$, etc., carries its own $\sigma$):
\begin{widetext}
\begin{equation}
\renewcommand{\arraystretch}{1.15}
\begin{array}{l|cccc|c|c}
\text{Peak} & \sigma_{V_1} & \sigma_{V_2} & \sigma_{V_3} & \sigma_{O_{\rm det}} & \prod\sigma & \text{Bertaut path}\\
\hline
(E_{12},E_{23}) & \sigma_C & \sigma_G & \sigma_C & \sigma_G & \sigma_F & \Lambda^5_C\Lambda^2_G\Lambda^7_C/\Lambda^2_G \\
(E_{12},E_{13}) & \sigma_C & \sigma_G & \sigma_G & \sigma_C & \sigma_F & \Lambda^5_C\Lambda^2_G\Lambda^2_G/\Lambda^5_C\\
(E_{12},\mathrm{qAFM}) & \sigma_G & \sigma_G & \sigma_G & \sigma_G & \sigma_F & \Lambda^2_G\Lambda^2_G\,G_z \\
(\mathrm{qAFM},E_{12}) & \sigma_F & \sigma_F & \sigma_G & \sigma_G & \sigma_F & F_x\,F_x\,(\sin\phi_{\rm qAFM}\,G_z)\,\Lambda^2_G\\
(\mathrm{qAFM},E_{13}) & \sigma_F & \sigma_F & \sigma_C & \sigma_C & \sigma_F & F_x\,F_x\,(\sin\phi_{\rm qAFM}\,G_z)\,\Lambda^5_C\\
(\mathrm{qFM},\mathrm{qAFM}) & \sigma_F & \sigma_F & \sigma_F & \sigma_G & \sigma_G \!\to\!\sigma_F\;\text{via}\;T^{(2)}\!&\!F_y\,F_y\,(T^{(2)}\!\!:\!\sigma_F\!\to\!\sigma_G)\,G_z
\end{array}
\label{eq:peakSigmaTable}
\end{equation}
\end{widetext}
For the first five rows the four-cycle of $\sigma$-vectors collapses
directly to $\sigma_F$ via the Klein algebra
$\sigma_C\!\cdot\!\sigma_C=\sigma_G\!\cdot\!\sigma_G=\sigma_F$, etc. For
the last row $({\rm qFM},{\rm qAFM})$ the $T^{(2)}$ vertex carries an
explicit $\sigma_F\!\to\!\sigma_G$ transformation (DM-induced
$F\leftrightarrow G$ admixture, App.~\ref{app:Bertaut2}), supplying the
final $\sigma_G$ needed to balance the detector.

\paragraph{Consequences for the pathway analysis.} The
selection-rule Eq.~\eqref{eq:sigmaRule} explains, without explicit
computation, why every cross-sector peak that links Fe to Tm requires
either (a) the small DM-induced
$\sin\phi_{\rm qAFM}$ admixture of $G_z$ into the qAFM eigenvector
($(\mathrm{qAFM},E_{12,13})$ and $(E_{12},\mathrm{qAFM})$), or (b) the
$T^{(2)}$ DM-driven $F\leftrightarrow G$ Bertaut mixing
$(\mathrm{qFM},\mathrm{qAFM})$. No ``$\sigma$-orphan'' peak --- one whose
$\sigma$-vector product would fail to close to $\sigma_F$ --- appears in
the six-peak list, and the
selection rule predicts that such peaks have strictly zero amplitude
at $\mathbf q=0$.

\subsection{Geometry I: \texorpdfstring{$H\parallel a, E\parallel c$, VV}{H||a, E||c, VV}}

The pump and detect interactions are $\hat V = h_x\Jeff_x + e_z\deff_z$,
which couples to $\lambda^{4,5,6,7}$ (mirror-odd Tm sector, $A_{1,2}^-$).
VV detection reads out the same combination.

\paragraph{$(E_{12}, E_{23})$.} Two pump interactions $V_1 = V_2 \in
\{\lambda^4,\lambda^5,\lambda^6,\lambda^7\}$ generate the
$|E_2\rangle\langle E_1|$ sector only at \emph{second order}: the odd-sector
bilinears $(4,6)$ and $(5,7)$ close onto $\lambda^1$ by anticommutation,
while $(4,7)$ and $(5,6)$ close onto $\lambda^2$ by the complementary
commutator/anticommutator channel. We denote the resulting effective
quadratic Geometry-I source on the $E_{12}$ coherence by
$\mathcal S_{12}^{\rm odd}$. The third odd-sector interaction then promotes
$|E_2\rangle$ to $|E_3\rangle$, producing emission at $\omega_{23}$.
Amplitude:
\begin{equation}
A_{(E_{12},E_{23})} \propto
\frac{\mathcal S_{12}^{\rm odd}\cdot M_{13}\cdot M_{23}}
{(\Gamma_{12}+i\omega_{12})(\Gamma_{23}+i\omega_{23})}\,e^{-\Gamma_{23}t_2}\,\rho_{00}^{\rm Tm},
\label{eq:A_E12E23}
\end{equation}
where $M_{\alpha\beta}\equiv\langle E_\alpha|\Jeff_x|E_\beta\rangle$ and
$\mathcal S_{12}^{\rm odd}$ is the effective second-order source built from
the allowed Geometry-I odd-sector vertices
$\{h_x J_x^{\rm eff}, e_z d_z^{\rm eff}\}$.

\paragraph{$(E_{12}, E_{13})$.} Same $E_{12}$ excitation as above; the
third interaction $V_3\sim\lambda^5$ promotes $|E_2\rangle$ to
$|E_3\rangle$ via the alternative odd-sector matrix element. In the
magnetic-$x$ readout channel the amplitude is analogous to
Eq.~\eqref{eq:A_E12E23} with $M_{23}\to M_{13}$:
\begin{equation}
A_{(E_{12},E_{13})}/A_{(E_{12},E_{23})} = M_{13}/M_{23} = 2.392/0.913 \approx 2.6.
\label{eq:E13overE23}
\end{equation}
This ratio is a parameter-free prediction of the model.

\paragraph{$(\mathrm{qAFM}, E_{13})$.} Fe-pumped pathway: the magnetic
field $h_x$ excites the qAFM mode via the Fe Zeeman vertex; the qAFM
eigenvector has a $\sin\phi_{\rm qAFM}\sim D_1/J_1$ admixture along
$\delta G_z$ (Sec.~\ref{sec:qAFMfreq}); the $\delta G_z$ coordinate
sources the linear $\bar\chi^{2z}\delta G_z\lambda^2$ exchange field on
the Tm qutrit, populating $|E_2\rangle\langle E_1|$ through the SU(3)
generator $\lambda^2$. A subsequent odd-sector interaction promotes to
$|E_3\rangle$, producing $E_{13}$ emission. Amplitude is
suppressed by the small ratio $D_1/J_1$ relative to the purely Tm
pathways:
\begin{equation}
A_{(\mathrm{qAFM},E_{13})} \propto (D_1/J_1)\cdot\bar\chi^{2z}\cdot\cdots,
\label{eq:A_qAFME13}
\end{equation}
making it the weakest of the four Geometry-I peaks.

\subsection{Geometry II: \texorpdfstring{$H\parallel a, E\parallel c$, VH}{H||a, E||c, VH}}

The pump is identical to Geometry I (mirror-odd Tm sector); VH
detection reads out the orthogonal mirror-even sector via
$\lambda^1\sim d_x^{\rm eff}$ and $\lambda^2\sim J_z^{\rm eff}$.

\paragraph{$(\mathrm{qAFM}, E_{12})$.} The cleanest linear-exchange
pathway in the dataset: $h_x\to$ qAFM $\to (D_1/J_1)\delta G_z \to
\bar\chi^{2z}\delta n^2 \to E_{12}$ emission via $J_z^{\rm eff}\sim
\lambda^2$. Amplitude:
\begin{equation}
\begin{aligned}
A_{(\mathrm{qAFM}, E_{12})} \propto\;& (g_e\mu_B)^2 (D_1/J_1)\bar\chi^{2z} |M_{12}|\\
&\times\frac{1}{(\Gamma_{\rm qAFM}+i\omega_{\rm qAFM})(\Gamma_{12}+i\omega_{12})}.
\end{aligned}
\label{eq:A_qAFME12}
\end{equation}
This is a third-order signal (two Fe Zeeman interactions and one Tm
emission) with three causal time orderings. The pure-Fe-pump character
makes it especially clean to identify experimentally.

\subsection{Geometry III: \texorpdfstring{$H\parallel c, E\parallel a$, VH}{H||c, E||a, VH}}

The pump directly addresses the even Tm sector via $h_z J_z^{\rm eff}\sim
h_z\lambda^2$ and $e_x d_x^{\rm eff}\sim e_x\lambda^1$.

\paragraph{$(E_{12}, \mathrm{qAFM})$.} Reciprocal of Geometry-II
pathway: the pump creates $\delta n^2$ directly; this drives $\delta G_z$
via $\bar\chi^{2z}$; $\delta G_z$ couples into the qAFM eigenvector
through its DM admixture; qAFM emits during $t_3$. Amplitude:
\begin{equation}
A_{(E_{12}, \mathrm{qAFM})} \propto
(g_J\mu_B)^2\cdot |M_{12}|^2\cdot\bar\chi^{2z}\cdot (D_1/J_1)\cdot (g_e\mu_B),
\label{eq:A_E12qAFM}
\end{equation}
which is the time-reversed pair of Eq.~\eqref{eq:A_qAFME12}.

\paragraph{$(\mathrm{qFM}, \mathrm{qAFM})$.} Fe nonlinearity: the
magnetic field $h_z$ couples directly to the qFM coordinate via the
Zeeman vertex; the qFM and qAFM modes are coupled at second order in the
LLG dynamics through the spin-length constraint and the DM-induced
mixing of $\bfF$ and $\bfG$. The bilinear coupling tensor
$T^{(2)}_{\rm qFM\to qAFM}$ is the effective quadratic Fe vertex inherited
from the second-order Bertaut expansion (App.~\ref{app:Bertaut2}):
\begin{equation}
T^{(2)}_{\rm qFM\to qAFM} \equiv \Omega^{(2)}_{\rm qFM,qAFM}
\label{eq:TqFMqAFM}
\end{equation}
whose symmetry is fixed by the DM-induced $F\leftrightarrow G$ mixing and
whose magnitude scales linearly with the small canting ratio $D_1/J_1$ in
a reduced two-sublattice description. The resulting peak amplitude is
$\propto T^{(2)}\cdot (g_e\mu_B)^2\cdot$ Fe matrix elements. This peak is
purely a Fe-sector nonlinearity and does not require Tm involvement.

\subsection{Summary peak table}

\begin{table*}[t]
\caption{Standard 2DCS pathway table for the six intrinsic peaks. Each
row lists the excitation frequency $\omega_\tau$, the detection
frequency $\omega_t$, the symmetry or dynamical coupling that must be
present, and the leading scaling of the corresponding peak amplitude.}
\label{tab:peaks}
\begin{ruledtabular}
\begin{tabular}{l|l|l|l|p{0.29\textwidth}|l}
Geometry & Peak & $\omega_\tau$ & $\omega_t$ & Required coupling & Leading scaling \\
\hline
I (VV) & $(E_{12},E_{23})$ & $\omega_{12}$ & $\omega_{23}$ & Odd-sector Tm-only pathway with source $\mathcal S_{12}^{\rm odd}$ and odd-sector $23$ readout & $\propto M_{13}M_{23}|M_{12}|^2$ \\
I (VV) & $(E_{12},E_{13})$ & $\omega_{12}$ & $\omega_{13}$ & Odd-sector Tm-only pathway with source $\mathcal S_{12}^{\rm odd}$ and odd-sector $13$ readout & $\propto M_{13}^2|M_{12}|^2$ \\
I (VV) & $(\mathrm{qAFM},E_{13})$ & $\omega_{\rm qAFM}$ & $\omega_{13}$ & DM admixture $\sin\phi_{\rm qAFM}$ plus linear exchange $\bar\chi^{2z}$ converting qAFM motion into the even $12$ channel before odd-sector promotion to $E_{13}$ & $\propto (D_1/J_1)\bar\chi^{2z} M_{13}$ \\
\hline
II (VH) & $(\mathrm{qAFM},E_{12})$ & $\omega_{\rm qAFM}$ & $\omega_{12}$ & Linear Fe$\to$Tm conversion through $\bar\chi^{2z}\sin\phi_{\rm qAFM}$ with VH readout in the mirror-even channel & $\propto (D_1/J_1)\bar\chi^{2z}|M_{12}|$ \\
\hline
III (VH) & $(E_{12},\mathrm{qAFM})$ & $\omega_{12}$ & $\omega_{\rm qAFM}$ & Reciprocal Tm$\to$Fe conversion through the same $\bar\chi^{2z}\sin\phi_{\rm qAFM}$ channel & $\propto |M_{12}|^2\bar\chi^{2z}(D_1/J_1)$ \\
III (VH) & $(\mathrm{qFM},\mathrm{qAFM})$ & $\omega_{\rm qFM}$ & $\omega_{\rm qAFM}$ & Fe-only quadratic mode-mixing vertex $T^{(2)}_{\rm qFM\to qAFM}$ from the second-order Bertaut expansion & $\propto T^{(2)}_{\rm qFM\to qAFM}$ \\
\end{tabular}
\end{ruledtabular}
\end{table*}

The reciprocal pair Geometry~II $(\mathrm{qAFM},E_{12})$ and
Geometry~III $(E_{12},\mathrm{qAFM})$ measure the same exchange vertex
$\bar\chi^{2z}$ in the two opposite time orderings; their amplitude ratio
in the rephasing vs.\ non-rephasing channels is a stringent self-consistency
test of the model.

\subsection{Explicit matrix-element catalogue}
\label{sec:MEcatalogue}

We collect here, for every one of the six Liouville pathways, the
\emph{four} operator factors that enter Eq.~\eqref{eq:Apeak} (three
vertices $V_1,V_2,V_3$ plus one detection $\hat O_{\rm det}$), the two
coherence denominators and the $t_2$ decay, in the same notation. Define
\begin{equation}
\begin{aligned}
g_e\!\!\!\!&\equiv g_e\mu_B,\quad g_J\equiv g_J\mu_B,\\
\mathcal D_{ab}(\omega)\!\!\!\!&\equiv \Gamma_{ab}+i(\omega-\omega_{ab}),\quad
\mathcal D_{\rm m}(\omega)\equiv \Gamma_{\rm m}+i(\omega-\omega_{\rm m}),
\end{aligned}
\label{eq:catalogueshorthand}
\end{equation}
for $\rm m\in\{qFM,qAFM\}$. The numerical values of the projected Tm
matrix elements (T-symmetric gauge, Eq.~\eqref{eq:Jeff}) are
\begin{equation}
\begin{aligned}
M_{12}^{(z)}&\equiv\langle E_1|J_z^{\rm eff}|E_2\rangle = -5.264\,i,\\
M_{13}^{(x)}&\equiv\langle E_1|J_x^{\rm eff}|E_3\rangle = -2.392\,i,\\
M_{23}^{(x)}&\equiv\langle E_2|J_x^{\rm eff}|E_3\rangle = -0.913\,i,\\
M_{13}^{(y)}&\equiv\langle E_1|J_y^{\rm eff}|E_3\rangle = +2.787\,i,\\
M_{23}^{(y)}&\equiv\langle E_2|J_y^{\rm eff}|E_3\rangle = -0.466\,i,
\end{aligned}
\label{eq:Mnumeric}
\end{equation}
all in units of $\hbar$. Diagonal-Cartan elements
$\langle E_\alpha|J_a|E_\alpha\rangle$ vanish identically in the
$T$-symmetric gauge for $\alpha=1$ and are absorbed into $h^{(0)}_{3,8}$
for $\alpha=2,3$ (Eq.~\eqref{eq:hbareCEF}).
For the electric-dipole channel we likewise write
$D_{12}^{(x)}\equiv\langle E_1|d_x^{\rm eff}|E_2\rangle=p_x$,
$D_{13}^{(z)}\equiv\langle E_1|d_z^{\rm eff}|E_3\rangle=p_{z4}$, and
$D_{23}^{(z)}\equiv\langle E_2|d_z^{\rm eff}|E_3\rangle=p_{z6}$.

\paragraph{Pathway notation.} For each peak we write the rephasing
($R$) and non-rephasing ($NR$) double-sided diagrams as
\begin{equation}
|0\rangle\langle 0|\xrightarrow{V_1}|a\rangle\langle 0|\xrightarrow{V_2}\bigl(|c\rangle\langle 0|\text{ or }|a\rangle\langle d|\bigr)\xrightarrow{V_3}|e\rangle\langle f|\xrightarrow{\hat O_{\rm det}},
\label{eq:Rdiag}
\end{equation}
where each arrow is a left- or right-multiplication by the indicated
vertex. The Tm states are $\{|E_1\rangle\!\equiv\!|0\rangle,|E_2\rangle,|E_3\rangle\}$;
the Fe magnon states are bosonic Fock states of the qFM and qAFM
oscillators; the joint vacuum is $|0\rangle_{\rm Tm}\otimes|{\rm vac}\rangle_{\rm Fe}$.

\paragraph{Geometry-I peaks.} The pump-detect channels are
$V_i\in\{h_x J_x^{\rm eff}, e_z d_z^{\rm eff}\}\propto\{\lambda^{5,7},\lambda^{4,6}\}$
during $t_1,t_3$. Because these odd-sector generators do not contain a
direct $\lambda^2$ vertex, the $E_{12}$ coherence in Geometry~I is
generated only at second order. We denote the resulting effective source by
$\mathcal S_{12}^{\rm odd}$; it collects the allowed odd-sector bilinears
$(4,6)$, $(5,7)$, $(4,7)$, and $(5,6)$ that close onto $\lambda^{1,2}$.
The four pathways:

\begin{widetext}
\begin{equation}
\begin{aligned}
A_{(E_{12},E_{23})}^{R}\!\!\!\! &= \tfrac{i}{\hbar^3}\,(g_J)^4\,
\underbrace{\langle E_1|J_x^{\rm eff}|E_3\rangle}_{=M_{13}^{(x)}}\,
\underbrace{\langle E_3|J_x^{\rm eff}|E_2\rangle^*}_{=M_{23}^{(x)*}}\,
\underbrace{\langle E_2|J_x^{\rm eff}|E_1\rangle^*}_{=M_{12}^{(x)*}=0}\,
\underbrace{\langle E_1|J_x^{\rm eff}|E_2\rangle}_{=M_{12}^{(x)}=0}\\
&\quad\Longrightarrow\;\text{$x$-only channel forbidden;}
\\[4pt]
A_{(E_{12},E_{23})}^{R({\rm odd})} &= \tfrac{i}{\hbar^3}\,\mathcal S_{12}^{\rm odd}\,(g_J h_x)^2\,
\bigl[M_{13}^{(x)}\bigr]\bigl[M_{23}^{(x)}\bigr]^*
\,\frac{e^{-\Gamma_{23}t_2}}{\mathcal D_{12}(-\omega_\tau)\,\mathcal D_{23}(\omega_t)},\\
&= \tfrac{i}{\hbar^3}\,\mathcal S_{12}^{\rm odd}\,(g_J h_x)^2\,(2.392)(-0.913)\,
\frac{e^{-\Gamma_{23}t_2}}{\mathcal D_{12}^*\mathcal D_{23}}\\
&\equiv -\tfrac{i}{\hbar^3}\,\mathcal S_{12}^{\rm odd}\,(g_J h_x)^2\cdot 2.18\cdot
\frac{e^{-\Gamma_{23}t_2}}{\mathcal D_{12}^*\mathcal D_{23}};
\end{aligned}
\label{eq:A_E12E23_full}
\end{equation}

\begin{equation}
\begin{aligned}
A_{(E_{12},E_{13})}^{R} &= \tfrac{i}{\hbar^3}\,\mathcal S_{12}^{\rm odd}(g_J h_x)^2\,
\bigl[M_{13}^{(x)}\bigr]^2\,
\frac{e^{-\Gamma_{13}t_2}}{\mathcal D_{12}^*\mathcal D_{13}}\\
&= \tfrac{i}{\hbar^3}\,\mathcal S_{12}^{\rm odd}(g_J h_x)^2\,(2.392)^2\,
\frac{e^{-\Gamma_{13}t_2}}{\mathcal D_{12}^*\mathcal D_{13}}\\
&\equiv -\tfrac{i}{\hbar^3}\,\mathcal S_{12}^{\rm odd}(g_J h_x)^2\cdot 5.72\cdot
\frac{e^{-\Gamma_{13}t_2}}{\mathcal D_{12}^*\mathcal D_{13}};
\end{aligned}
\label{eq:A_E12E13_full}
\end{equation}

\begin{equation}
\begin{aligned}
A_{(\mathrm{qAFM},E_{13})}^{R}&=\tfrac{i}{\hbar^3}(g_e h_x)\bigl[(g_e S)\sin\phi_{\rm qAFM}\bar\chi^{2z}\bigr](g_J h_x)^2\,
\bigl[M_{13}^{(x)}\bigr]\bigl[M_{12}^{(z)}\bigr]^*\\
&\quad\times\langle{\rm qAFM}|\delta G_z|{\rm vac}\rangle\;
\frac{e^{-\Gamma_{13}t_2}}{\mathcal D_{\rm qAFM}(-\omega_\tau)\,\mathcal D_{13}(\omega_t)},
\end{aligned}
\label{eq:A_qAFME13_full}
\end{equation}
with $\sin\phi_{\rm qAFM}=\mathcal O(D_1/J_1)$ the small qAFM
$\delta G_z$ admixture (Sec.~\ref{sec:qAFMfreq}).

\paragraph{Geometry-II peak.} VH detection along $\hat z$:
\begin{equation}
\begin{aligned}
A_{(\mathrm{qAFM},E_{12})}^{R} &= \tfrac{i}{\hbar^3}(g_e h_x)^2\,
\bigl[\sin\phi_{\rm qAFM}\,\bar\chi^{2z}\bigr]\,(g_J h_z)\,
\bigl[M_{12}^{(z)}\bigr]\bigl[M_{12}^{(z)}\bigr]^*\\
&\quad\times \frac{e^{-\Gamma_{12}t_2}}{\mathcal D_{\rm qAFM}(-\omega_\tau)\,\mathcal D_{12}(\omega_t)}\\
&= \tfrac{i}{\hbar^3}(g_e)^2 g_J h_x^2 h_z\,
\sin\phi_{\rm qAFM}\,\bar\chi^{2z}\,(5.264)^2\,
\frac{e^{-\Gamma_{12}t_2}}{\mathcal D_{\rm qAFM}^*\mathcal D_{12}}.
\end{aligned}
\label{eq:A_qAFME12_full}
\end{equation}

\paragraph{Geometry-III peaks.}
\begin{equation}
\begin{aligned}
A_{(E_{12},\mathrm{qAFM})}^{R} &= \tfrac{i}{\hbar^3}(g_J h_z)^2\,
\bigl[M_{12}^{(z)}\bigr]^*\bigl[M_{12}^{(z)}\bigr]\,
\bar\chi^{2z}\,\bigl[\sin\phi_{\rm qAFM}\bigr]\,(g_e h_z)\,
\langle{\rm qAFM}|\delta G_z|{\rm vac}\rangle\\
&\quad\times \frac{e^{-\Gamma_{\rm qAFM}t_2}}{\mathcal D_{12}(-\omega_\tau)\,\mathcal D_{\rm qAFM}(\omega_t)},
\end{aligned}
\label{eq:A_E12qAFM_full}
\end{equation}
the time-reverse of Eq.~\eqref{eq:A_qAFME12_full}, and
\begin{equation}
\begin{aligned}
A_{(\mathrm{qFM},\mathrm{qAFM})}^{R} &= \tfrac{i}{\hbar^3}(g_e h_z)^3\,
T^{(2)}_{\rm qFM,qAFM}\,(g_e h_x)\,
\langle{\rm qFM}|\delta F_z|{\rm vac}\rangle\,\langle{\rm qAFM}|\delta G_y|{\rm vac}\rangle^2\\
&\quad\times \frac{e^{-\Gamma_{\rm qAFM}t_2}}{\mathcal D_{\rm qFM}(-\omega_\tau)\,\mathcal D_{\rm qAFM}(\omega_t)},
\end{aligned}
\label{eq:A_qFMqAFM_full}
\end{equation}
with $T^{(2)}_{\rm qFM,qAFM}$ the effective quadratic mixing rate of
Eq.~\eqref{eq:TqFMqAFM}, and $\langle{\rm m}|\delta X|{\rm vac}\rangle =
\sqrt{S/2}\,(u_X^{\rm m} + v_X^{\rm m})$ from the Bogoliubov transformation
diagonalizing $M^{\Gamma_2}_{4\times 4}$.
\end{widetext}

\paragraph{Magnitudes and ratios.} All six amplitudes share the
overall factor $(i/\hbar^3)$ and depend on the same matrix-element
catalogue Eq.~\eqref{eq:Mnumeric}. The cleanest model-independent ratios
predicted from the catalogue are
\begin{equation}
\begin{aligned}
\frac{A_{(E_{12},E_{13})}}{A_{(E_{12},E_{23})}} &= \frac{M_{13}^{(x)}}{M_{23}^{(x)}}=\frac{-2.392\,i}{-0.913\,i}\approx +2.62,\\
\frac{A_{(\mathrm{qAFM},E_{12})}^{\rm Geo II}}{A_{(E_{12},\mathrm{qAFM})}^{\rm Geo III}} &= \frac{h_x^2 h_z\,\mathcal D_{12}^{-1}\mathcal D_{\rm qAFM}^{-1}}
{h_z^2\,\mathcal D_{12}^{-1}\mathcal D_{\rm qAFM}^{-1}}\biggr|_{\text{equal $h_x,h_z$}}=1,\\
\frac{A_{(\mathrm{qFM},\mathrm{qAFM})}}{A_{(E_{12},\mathrm{qAFM})}} &\sim
\frac{T^{(2)}\,(g_e h_z)^2}{(g_J h_z)^2 |M_{12}|^2\,\bar\chi^{2z}},
\end{aligned}
\label{eq:peakratios}
\end{equation}
the first of which is parameter-free. The other two reduce the
six-amplitude problem to the exchange vertex $\bar\chi^{2z}$ and the
Fe-sector nonlinear coefficient $T^{(2)}$ once the catalogue
Eq.~\eqref{eq:Mnumeric} and the linearized magnon sector are fixed.

% =====================================================================
\section{Perturbative EOM derivation of \texorpdfstring{$\chi^{(3)}$}{chi(3)}}
\label{sec:EOMexpansion}

The Liouville-pathway analysis of Sec.~\ref{sec:2dcs} can be re-derived
from a complementary, completely classical (Heisenberg) viewpoint by
expanding the joint Bertaut+Bloch EOM order by order in the drive
amplitude. This provides (i) an independent check of every matrix
element in the catalogue Eq.~\eqref{eq:Mnumeric}, (ii) explicit access
to the nonlinear vertices $N_2,N_3$ that mediate the cross-sector
peaks, and (iii) a clean separation of the Fe-magnon nonlinearities
(spin-length constraint, DM\,$\times$\,$K$ mixing) from the Tm SU(3)
non-Abelian nonlinearities (Gell-Mann commutators).

\subsection{Setup: nonlinear joint EOM}
\label{sec:EOMsetup}

Including all cubic and quartic terms retained in the Hamiltonian, the
joint EOM for the 10-component variable $\bfw$ of
Eq.~\eqref{eq:wjoint} (extended by the two diagonal Tm coordinates
$\delta n^3,\delta n^8$ and the four omitted Fe components,
$\bfw\to\tilde\bfw\in\mathbb R^{16}$) has the structure
\begin{equation}
\boxed{\;
\dot{\tilde\bfw}^a = M^a{}_b\,\tilde\bfw^b
+ N_2^{a}{}_{bc}\,\tilde\bfw^b\tilde\bfw^c
+ N_3^{a}{}_{bcd}\,\tilde\bfw^b\tilde\bfw^c\tilde\bfw^d
+ f^a(t),\;}
\label{eq:joinNLEOM}
\end{equation}
with:
\begin{itemize}[leftmargin=*]
\item $M$: the linear matrix Eq.~\eqref{eq:Mjoint};
\item $N_2$: the cubic vertices generated by (a) the SU(3) commutator
$\dot\rho_a = -i f_{abc} h^{(0)}_b\rho_c$ entering at the Tm bilinear
order, (b) the Bertaut LLG $\bfX\times\bfH$ entering at the Fe bilinear
order via the spin-length constraint $\bfF^2+\bfG^2+\bfC^2+\bfA^2=4S^2$,
and (c) the cross-sector vertex $u_a\lambda^a\,\mathcal I^{\rm Fe}$
linearized in the Fe coordinate;
\item $N_3$: the quartic vertices from (a) the
$f_{abc}f_{cde}\rho_b\rho_d\rho_e$ Tm triple commutator and (b) the
Van Vleck back-action expanded to bilinear in Fe and linear in
Tm coherence;
\item $f(t)$: the drive force, linear in the field amplitude through
the Zeeman and electric-dipole vertices Eqs.~\eqref{eq:Hdrivem}--\eqref{eq:Hdrivee}.
\end{itemize}

\paragraph{Explicit form of the nonlinear vertices.}
The Tm-only $N_2$ vertex follows from $\dot n^a = f^{abc}\,h^{\rm tot}_b n^c$
with $h^{\rm tot}_b = h^{(0)}_b + h^{(1)}_b(t) + \delta h_b[\bfX^{\rm Fe}]$;
expanding,
\begin{equation}
\dot n^a\big|_{\mathcal O(\tilde\bfw^2)}
= N_2^{a,{\rm Tm}}{}_{bc}\,\tilde\bfw^b\tilde\bfw^c
= f^{abd}\,\delta h_b[\tilde\bfw]\,\delta n_d\big|_{\mathcal O(\tilde\bfw^2)},
\label{eq:N2Tm}
\end{equation}
which contains the $\bar\chi^{2z}\,\delta G_z\,\lambda^2$ vertex
contracted by the structure constants $f_{2bc}$ to drive any of the
six Tm coherences after projection onto the qAFM admixture. The Fe-only
$N_2$ vertex follows from the bilinear part of the Bertaut precession
equation,
\begin{equation}
\dot{\bfX}\big|_{\mathcal O(\tilde\bfw^2)}
= N_2^{{\bfX},{\rm Fe}}{}_{bc}\,\tilde\bfw^b\tilde\bfw^c
= -\frac{1}{\hbar}\sum_{Y\cdot Z=X}\delta\bfY\times\delta\bfH^Z\big|_{\mathcal O(\tilde\bfw^2)},
\label{eq:N2Fe}
\end{equation}
and contains, in particular, the qFM--qAFM bilinear coupling
$T^{(2)}_{\rm qFM,qAFM}$ derived in App.~\ref{app:Bertaut2}. The
cross-sector vertex
\begin{equation}
N_2^{a,{\rm Fe\leftarrow Tm}}{}_{bc} = -\frac{1}{\hbar}\,\frac{\partial^2 H_{\rm aniso\text{-}Tm}}
{\partial\tilde\bfw^b\,\partial\tilde\bfw^c}\bigg|_0
\label{eq:N2cross}
\end{equation}
encodes the Van Vleck back-action Eq.~\eqref{eq:Hanisogen}: setting
$b=\delta n^1$ and $c=F_x$ extracts the coefficient
$u_1\,\alpha_2^{\rm Fe}\,F_x$ that appears in
the population-induced Fe anisotropy shift.

\paragraph{Bertaut reading of the dynamical variables.} Throughout
this section, every Tm dynamical variable is implicitly the
\emph{Bertaut combination} $\Lambda^a_X$ of Eq.~\eqref{eq:LambdaXdef}
selected by the drive (Eqs.~\eqref{eq:HdriveBertaut}) rather than the
bare single-site coordinate. Specifically:
\begin{widetext}
\begin{equation}
\begin{aligned}
\rho_{12}(t) &\equiv \Lambda^2_G(t)
  &&\text{($\sigma_G$-staggered Tm coherence; couples to $h_z, G_z$)},\\
\rho_{13}^{(x)}(t) &\equiv \Lambda^5_C(t),\quad
\rho_{13}^{(y)}(t) \equiv \Lambda^7_A(t)
  &&\text{(separate $\sigma_C, \sigma_A$ projections of the }13\text{ coherence)},\\
\rho_{23}^{(x)}(t) &\equiv \Lambda^7_C(t),\quad
\rho_{23}^{(y)}(t) \equiv \Lambda^5_A(t),\\
n_1(t)-n_1^{\rm eq} &\equiv \Lambda^3_F(t),\qquad n_2(t)-n_2^{\rm eq} \equiv \Lambda^8_F(t),\\
a_{\rm qFM}(t) &\in \{F_z, G_x\}\text{-block (the $ac$-plane }\Gamma_2\text{ rotation)},\\
a_{\rm qAFM}(t) &\in \{F_y, G_y\}\text{-block, with } \delta G_z\text{ admixture from DM},
\end{aligned}
\label{eq:dynBertaut}
\end{equation}
\end{widetext}
i.e., the scalar amplitudes $\rho_{\alpha\beta}, n_a, a_{\rm m}$
should be read as the appropriate Bertaut-projected components
labeled by definite $Pbnm$ irreps. Every coupling vertex in
Eqs.~\eqref{eq:dotqFM}--\eqref{eq:dotn2} is then the contraction of
matched $\sigma$-vectors guaranteed by the global selection rule
Eq.~\eqref{eq:sigmaRule}: e.g., the $\xi|\rho_{12}|^2$ source on
$\dot a_{\rm qAFM}$ closes as $\sigma_G\!\cdot\!\sigma_G=\sigma_F$ on
the $|\rho_{12}|^2$ side and is projected back onto $G_z\subset$ qAFM,
which is what makes the vertex symmetry-allowed. The scalar
trees of Sec.~\ref{sec:EOMscalar} are one-to-one with the six
$\sigma$-rule rows of
Eq.~\eqref{eq:peakSigmaTable}, and the
peak count is reproduced \emph{by symmetry} alone.

\subsection{Standard second-order normal-coordinate form}

For direct comparison with the magnetic-2DCS literature one may
eliminate the conjugate momenta and rewrite the same dynamics in terms
of the real coordinates $Q_\mu$ of Eq.~\eqref{eq:Qcoords}. Extending the
linear equation Eq.~\eqref{eq:Qlinear} by the lowest nonlinear terms
gives the standard form
\begin{equation}
\ddot Q_\mu + 2\Gamma_\mu \dot Q_\mu + \Omega_\mu^2 Q_\mu
= F_\mu(t) + \sum_{\nu\rho} g_{\mu\nu\rho}\,Q_\nu Q_\rho + \cdots,
\label{eq:QmodeEOM}
\end{equation}
where the nonzero $g_{\mu\nu\rho}$ are fixed by the same symmetry rules
as the Liouville pathways. In the present problem the required quadratic
vertices are, schematically,
\begin{equation}
\begin{aligned}
Q_{\rm qAFM} &\leftarrow T^{(2)}_{\rm qAFM\leftrightarrow qFM}\,Q_{\rm qFM}Q_{\rm qAFM},\\
Q_{\rm qFM} &\leftarrow \widetilde T^{(2)}_{\rm qFM\leftarrow qAFM^2}\,Q_{\rm qAFM}^2,\\
Q_{12} &\leftarrow g_{12,13,23}\,Q_{13}Q_{23},
\end{aligned}
\label{eq:QmodeVertices}
\end{equation}
with $g_{12,13,23}\propto \mathcal S_{12}^{\rm odd}$. Keeping the slow
population coordinate $P_{12}\equiv n_1-n_2$ adds the zero-frequency
branch $P_{12}\leftarrow g_{P,12,12}Q_{12}^2$, which is the formal
$(E_{12},0)$ response discussed below. Equation~\eqref{eq:QmodeEOM} is
the standard normal-coordinate version of the first-order Bloch+LLG
system written explicitly in Sec.~\ref{sec:EOMscalar}.

\subsection{Explicit scalar EOMs in the physical mode basis}
\label{sec:EOMscalar}

It is useful to write Eq.~\eqref{eq:joinNLEOM} component by component
in the diagonal eigenmode basis of $M^{\Gamma_2}$. Define the complex
mode amplitudes
\begin{align}
a_{\rm qFM}(t) &= q_{\rm qFM}+i p_{\rm qFM}/(\hbar\omega_{\rm qFM}),\\
a_{\rm qAFM}(t)&= q_{\rm qAFM}+i p_{\rm qAFM}/(\hbar\omega_{\rm qAFM}),
\end{align}
for the two Fe magnons (so $a\sim e^{-i\omega t}$ at the linear order),
together with the six Tm coherences
$\rho_{\alpha\beta}(t)\equiv\langle E_\alpha|\hat\rho|E_\beta\rangle$
($\alpha\neq\beta$) and the two independent populations
$n_1,n_2$ (with $n_3=1-n_1-n_2$). Then
Eq.~\eqref{eq:joinNLEOM} unfolds into the following coupled system,
keeping the representative terms needed to expose each peak topology in
Sec.~\ref{sec:MEcatalogue}. For the electric-dipole channel we use the
same shorthand as in Sec.~\ref{sec:MEcatalogue}, namely
$D_{12}^{(x)}=\langle E_1|d_x^{\rm eff}|E_2\rangle$,
$D_{13}^{(z)}=\langle E_1|d_z^{\rm eff}|E_3\rangle$, and
$D_{23}^{(z)}=\langle E_2|d_z^{\rm eff}|E_3\rangle$:

\medskip
\noindent\emph{(i) Fe quasi-FM magnon.} Drive-only at first order; at
second order the Tm population $n_1$ acts as a static restoring-field
shift through the Van Vleck channel:
\begin{widetext}
\begin{equation}
\boxed{\;
\dot a_{\rm qFM} = -(i\omega_{\rm qFM}+\Gamma_{\rm qFM})\,a_{\rm qFM}
+\underbrace{\tfrac{g_e\mu_B}{\hbar}h_z(t)}_{\text{Zeeman drive}}
+\underbrace{T^{(2)}_{\rm qFM,qAFM}\,a_{\rm qAFM}^2}_{\text{spin-length}}
+\underbrace{\zeta_1\,n_1\,a_{\rm qFM}}_{\text{Van Vleck shift}}.\;}
\label{eq:dotqFM}
\end{equation}
\end{widetext}

\noindent\emph{(ii) Fe quasi-AFM magnon.} The DM-bilinear $T^{(2)}$
vertex feeds qFM into qAFM and vice versa; the Tm population $n_1$
both shifts and rotates the qAFM mode (because Van Vleck modifies
$K_2^{\rm Fe}$, which is what gives qAFM its DM-protected gap):
\begin{widetext}
\begin{equation}
\boxed{\;
\dot a_{\rm qAFM} = -(i\omega_{\rm qAFM}+\Gamma_{\rm qAFM})\,a_{\rm qAFM}
+ \sin\phi_{\rm qAFM}\,\tfrac{g_e\mu_B}{\hbar}h_x(t)
+ T^{(2)}_{\rm qAFM,qFM}\,a_{\rm qFM}\,a_{\rm qAFM}
+ \zeta_2\,n_1\,a_{\rm qAFM}
+ \xi\,|\rho_{12}|^2.\;}
\label{eq:dotqAFM}
\end{equation}
\end{widetext}
The very last term, bilinear in the Tm coherence, is the source of
the $(E_{12},\,{\rm qAFM})$ rephasing peak: a coherent grating in
$\rho_{12}$ at frequency $\omega_2-\omega_1$ launches a qAFM magnon
at the same frequency through the cross vertex
$\xi = \partial^2 H_{\rm aniso\text{-}Tm}/\partial|\rho_{12}|^2\,\partial G_z|_0$.

\medskip
\noindent\emph{(iii) Tm $|1\rangle\!\leftrightarrow\!|2\rangle$ coherence.}
Driven by $h_z$ via $M_{12}^{(z)}$; modulated nonlinearly by Fe magnons
through $\delta G_z$ and by populations through SU(3) symmetric vertex:
\begin{widetext}
\begin{equation}
\boxed{\;
\dot\rho_{12} = -(i\omega_{12}+\Gamma_{12})\,\rho_{12}
+ \frac{i}{\hbar}M_{12}^{(z)}\,g_J h_z(t)\,(n_1-n_2)
+ \frac{i}{\hbar}D_{12}^{(x)}\,e_x(t)\,(n_1-n_2)
+ \frac{i}{\hbar}\bar\chi^{2z}\,\delta G_z(t)\,(n_1-n_2)
+ \frac{i}{\hbar}\bigl[M_{13}^{(x)}g_J h_x(t)+D_{13}^{(z)}e_z(t)\bigr]\,\rho_{32}
- \frac{i}{\hbar}\bigl[M_{23}^{(x)*}g_J h_x(t)+D_{23}^{(z)*}e_z(t)\bigr]\,\rho_{13}.\;}
\label{eq:dotrho12}
\end{equation}
\end{widetext}
The $h_x/e_z$ odd-sector rotation terms are the Liouville $R$-pathway
junctions that rotate $\rho_{13}\to\rho_{12}$ during $t_2$ and produce
the $(E_{12},E_{23})$ and $(E_{12},E_{13})$ peaks.

\medskip
\noindent\emph{(iv) Tm $|1\rangle\!\leftrightarrow\!|3\rangle$ and
$|2\rangle\!\leftrightarrow\!|3\rangle$ coherences.}
\begin{widetext}
\begin{align}
\dot\rho_{13} &= -(i\omega_{13}+\Gamma_{13})\,\rho_{13}
+\tfrac{i}{\hbar}\bigl[M_{13}^{(x)}g_J h_x+M_{13}^{(y)}g_J h_y+D_{13}^{(z)}e_z\bigr](n_1-n_3)
+\tfrac{i}{\hbar}\bigl[M_{12}^{(z)}g_J h_z+D_{12}^{(x)}e_x\bigr]\,\rho_{23}
-\tfrac{i}{\hbar}\bigl[M_{23}^{(x)*}g_J h_x+D_{23}^{(z)*}e_z\bigr]\,\rho_{12},
\label{eq:dotrho13}\\
\dot\rho_{23} &= -(i\omega_{23}+\Gamma_{23})\,\rho_{23}
+\tfrac{i}{\hbar}\bigl[M_{23}^{(x)}g_J h_x+M_{23}^{(y)}g_J h_y+D_{23}^{(z)}e_z\bigr](n_2-n_3)
+\tfrac{i}{\hbar}\bigl[M_{12}^{(z)*}g_J h_z+D_{12}^{(x)*}e_x\bigr]\,\rho_{13}
-\tfrac{i}{\hbar}\bigl[M_{13}^{(x)*}g_J h_x+D_{13}^{(z)*}e_z\bigr]\,\rho_{12}.
\label{eq:dotrho23}
\end{align}
\end{widetext}

\medskip
\noindent\emph{(v) Tm populations.} Driven only at second order,
forming the diagonal SU(3) gratings that mediate every cross-sector
peak:
\begin{widetext}
\begin{align}
\dot n_1 &= -\Gamma_{1}^{\rm pop}(n_1-n_1^{\rm eq})
+\tfrac{2}{\hbar}{\rm Im}\!\bigl[M_{12}^{(z)*}g_J h_z(t)\,\rho_{12}\bigr]
+\tfrac{2}{\hbar}{\rm Im}\!\bigl[D_{12}^{(x)*}e_x(t)\,\rho_{12}\bigr]
+\tfrac{2}{\hbar}{\rm Im}\!\bigl[(M_{13}^{(x)*}g_J h_x(t)+D_{13}^{(z)*}e_z(t))\,\rho_{13}\bigr],
\label{eq:dotn1}\\
\dot n_2 &= -\Gamma_{2}^{\rm pop}(n_2-n_2^{\rm eq})
-\tfrac{2}{\hbar}{\rm Im}\!\bigl[M_{12}^{(z)*}g_J h_z(t)\,\rho_{12}\bigr]
+\tfrac{2}{\hbar}{\rm Im}\!\bigl[D_{12}^{(x)*}e_x(t)\,\rho_{12}\bigr]
+\tfrac{2}{\hbar}{\rm Im}\!\bigl[(M_{23}^{(x)*}g_J h_x(t)+D_{23}^{(z)*}e_z(t))\,\rho_{23}\bigr].
\label{eq:dotn2}
\end{align}
\end{widetext}
The $h_z/e_x$ source on $\dot n_1$ and $\dot n_2$ is precisely the
second-order vertex that creates the $\delta n^1$ grating used in
the Van Vleck cascade $\zeta_1 n_1 a_{\rm qFM}$ of
Eq.~\eqref{eq:dotqFM} and $\xi|\rho_{12}|^2$ of
Eq.~\eqref{eq:dotqAFM}. The same source is the formal origin of the
population branch $(E_{12},0)$ discussed in Sec.~\ref{sec:2dcs}.

\medskip
\noindent\emph{Reading off the six peaks.} Iterating
Eqs.~\eqref{eq:dotqFM}--\eqref{eq:dotn2} to third order in the drive
generates exactly the six nonlinear processes (one per peak
in Sec.~\ref{sec:MEcatalogue}). Each peak is identified by its tree
topology in the scalar EOM:
\begin{widetext}
\begin{center}
\renewcommand{\arraystretch}{1.15}
\begin{tabular}{ll}
\hline
Peak & Tree in Eqs.~\eqref{eq:dotqFM}--\eqref{eq:dotn2}\\
\hline
$(E_{12},E_{23})$  & $\{h_xJ_x^{\rm eff},e_z d_z^{\rm eff}\}^{\otimes 2}{\to}\rho_{12}^{\rm odd}$; then odd-sector readout at $\omega_{23}$\\
$(E_{12},E_{13})$  & same $\rho_{12}^{\rm odd}$ source, followed by odd-sector readout at $\omega_{13}$\\
$(E_{12},{\rm qAFM})$& $(h_z,e_x){\to}\rho_{12}$; then $\bar\chi^{2z}\rho_{12}{\to}\delta G_z{\to}a_{\rm qAFM}$ through the qAFM admixture\\
$({\rm qAFM},E_{12})$& $h_x{\to}a_{\rm qAFM}{\to}\delta G_z$; then $\bar\chi^{2z}\delta G_z(n_1-n_2)$ in \eqref{eq:dotrho12}\\
$({\rm qAFM},E_{13})$& $h_x{\to}a_{\rm qAFM}{\to}\delta G_z$; then $\bar\chi^{2z}\delta G_z{\to}\rho_{12}$, followed by odd-sector rotation into $\rho_{13}$\\
$({\rm qFM},{\rm qAFM})$& $h_z{\to}a_{\rm qFM}$, followed by the effective Fe-only quadratic mixing vertex $T^{(2)}$ before qAFM readout\\
\hline
\end{tabular}
\end{center}
\end{widetext}

In this language the $\dot a_{\rm qAFM}$ EOM contains three distinct
nonlinear sources --- the effective Fe-only vertex $T^{(2)}$,
$\zeta_2 n_1 a_{\rm qAFM}$, and $\xi|\rho_{12}|^2$ --- each one a different
2DCS peak; and the
$\dot\rho_{12}$ EOM contains four ($M^{(z)}h_z(n_1-n_2)$,
$\bar\chi^{2z}\delta G_z(n_1-n_2)$, and the two SU(3) raising terms
$h_x\rho_{32}$, $h_x\rho_{13}$). The total count of independent
trilinear terms across Eqs.~\eqref{eq:dotqFM}--\eqref{eq:dotn2} is
six --- matching the six peaks of Table~\ref{tab:peaks}.

\subsection{Order-by-order expansion}

Writing $f(t) = \lambda f^{(1)}(t)$ for a formal small parameter
$\lambda$ counting powers of the drive, expand
$\tilde\bfw = \sum_{n\geq 1}\lambda^n\tilde\bfw^{(n)}$. Substituting in
Eq.~\eqref{eq:joinNLEOM} and collecting powers:
\begin{align}
\dot{\tilde\bfw}^{(1)} &= M\,\tilde\bfw^{(1)} + f^{(1)}(t),
\label{eq:order1EOM}\\
\dot{\tilde\bfw}^{(2)} &= M\,\tilde\bfw^{(2)} + N_2[\tilde\bfw^{(1)},\tilde\bfw^{(1)}],
\label{eq:order2EOM}\\
\dot{\tilde\bfw}^{(3)} &= M\,\tilde\bfw^{(3)} + 2 N_2[\tilde\bfw^{(1)},\tilde\bfw^{(2)}]
+ N_3[\tilde\bfw^{(1)},\tilde\bfw^{(1)},\tilde\bfw^{(1)}].
\label{eq:order3EOM}
\end{align}
In Fourier space with the resolvent (Green's function)
$G(\omega)\equiv (i\omega\,\mathds{1} - M)^{-1}$,
\begin{equation}
\begin{aligned}
\tilde\bfw^{(1)}(\omega_1)\!\!\!\! &= G(\omega_1)f^{(1)}(\omega_1),\\
\tilde\bfw^{(2)}(\omega_1{+}\omega_2)\!\!\!\! &= G(\omega_1+\omega_2)\,N_2\!\bigl[\tilde\bfw^{(1)}(\omega_1),\tilde\bfw^{(1)}(\omega_2)\bigr],\\
\tilde\bfw^{(3)}(\Omega)\!\!\!\! &= G(\Omega)\Bigl\{2 N_2[\tilde\bfw^{(1)},\tilde\bfw^{(2)}]\\
&\quad{} + N_3[\tilde\bfw^{(1)},\tilde\bfw^{(1)},\tilde\bfw^{(1)}]\Bigr\},
\end{aligned}
\label{eq:GFsolutions}
\end{equation}
with $\Omega = \omega_1+\omega_2+\omega_3$. The detected polarization
$P_{\rm det}^{(3)}(t) = \langle\hat O_{\rm det},\,\tilde\bfw^{(3)}(t)\rangle$
gives the third-order signal, and $\chi^{(3)}(\omega_t,\omega_2,\omega_\tau)$
is read off by Fourier transform with $\omega_1=\pm\omega_\tau$,
$\omega_2$ slow, $\omega_3=\omega_t-\omega_1-\omega_2$.

\subsection{Numerical 2DCS simulation protocol}

For comparison with experiment one solves the full nonlinear equations
Eq.~\eqref{eq:joinNLEOM} (or the reduced normal-coordinate form
Eq.~\eqref{eq:QmodeEOM}) three times with identical pulse envelopes and
damping parameters: once with both pulses present, once with pulse $A$
only, and once with pulse $B$ only. The simulated nonlinear field is
then constructed by the same subtraction used experimentally,
Eq.~\eqref{eq:ENL}, and Fourier transformed as
\begin{equation}
S_{\rm sim}(\omega_t,t_2,\omega_\tau)
= \int_0^\infty\!dt_3\,e^{i\omega_t t_3}\int_0^\infty\!dt_1\,e^{\pm i\omega_\tau t_1}
E_{\rm sig}^{\rm NL}(t_3,t_2,t_1).
\label{eq:Ssim2DCS}
\end{equation}
In practice Eq.~\eqref{eq:Hmicro} is the microscopic starting point for
Monte Carlo or bond-resolved parameter sampling, while
Eqs.~\eqref{eq:Htotal}, \eqref{eq:Qlinear}, and \eqref{eq:QmodeEOM}
provide the reduced driven dynamics whose peak positions, amplitudes,
and geometry selectivity are compared against Table~\ref{tab:peaks}.

\subsection{Structure of \texorpdfstring{$G(\omega)$}{G(omega)}}

The resolvent factorizes into uncoupled blocks at lowest order in
$D/J$:
\begin{equation}
G(\omega) = G^{\rm Fe}(\omega)\oplus G^{\rm Tm}(\omega)
+ \mathcal O(D_1/J_1),
\label{eq:Gblock}
\end{equation}
with the Tm block diagonal in the three coherence pairs,
\begin{equation}
G^{\rm Tm}_{\alpha\beta\to\alpha\beta}(\omega)
= \frac{1}{i(\omega-\omega_{\alpha\beta}) + \Gamma_{\alpha\beta}}
\equiv \mathcal D_{\alpha\beta}^{-1}(\omega),
\label{eq:GTm}
\end{equation}
which are exactly the coherence denominators of
Eq.~\eqref{eq:catalogueshorthand}. The Fe block, after Bogoliubov
diagonalization of $M^{\Gamma_2}_{4\times 4}$, becomes diagonal in the
$\{{\rm qFM},{\rm qAFM}\}$ eigenmodes with denominators
$\mathcal D_{\rm qFM,qAFM}(\omega)$. The off-block correction
$\mathcal O(D_1/J_1)$ produces the $\sin\phi_{\rm qAFM}\sim D_1/J_1$
admixture used throughout Sec.~\ref{sec:2dcs}.

\subsection{Pathway-by-pathway recovery of \texorpdfstring{$\chi^{(3)}$}{chi(3)}}

We now show that each Liouville pathway of Sec.~\ref{sec:2dcs} arises
from a unique tree diagram in
Eqs.~\eqref{eq:order1EOM}--\eqref{eq:order3EOM}, and that the
matrix-element products are the same as in
Sec.~\ref{sec:MEcatalogue}.

\paragraph{Tm-only $\chi^{(3)}$ — pathway $(E_{12},E_{23})$.}
Pure odd-sector SU(3) cascade. Two first-order Geometry-I odd-sector
drives on $\lambda^{4,5,6,7}$ generate an effective source
$\mathcal S_{12}^{\rm odd}$ on the $\lambda^{1,2}$ sector; a third odd-sector
vertex then reads out the $23$ coherence. Schematically,
\begin{equation}
\begin{aligned}
\tilde\bfw^{(3)}_{\lambda^7}(\Omega)\!\!\!\! &= G^{\rm Tm}_{23}(\Omega)\;
\mathcal S_{12}^{\rm odd}(\omega_1,\omega_2)\,\tilde\bfw^{(1)}_{\lambda^7}(\omega_3),
\end{aligned}
\label{eq:E12E23EOM}
\end{equation}
where $\mathcal S_{12}^{\rm odd}$ collects the allowed odd-sector bilinears
$(4,6)$, $(5,7)$, $(4,7)$, and $(5,6)$ that close onto
$\lambda^{1,2}$. This yields the denominator structure and matrix-element
bookkeeping summarized in Eq.~\eqref{eq:A_E12E23_full}.

\paragraph{Fe$\to$Tm exchange — pathway $(\mathrm{qAFM},E_{12})$.}
Two-step diagram: (i)
$\tilde\bfw^{(1)}_{\delta G_z} = G^{\rm Fe}_{\rm qAFM}(\omega_1)\sin\phi_{\rm qAFM}\,(g_e h_x)$
(magnon excited then projected onto $\delta G_z$); (ii)
$\tilde\bfw^{(2)}_{\lambda^2} = G^{\rm Tm}_{12}(\omega_1+\omega_2)\,
\bar\chi^{2z}\,\tilde\bfw^{(1)}_{\delta G_z}(\omega_1)\,(g_J h_z)$;
(iii) detection at $\omega_t = \omega_{12}$ gives
\begin{widetext}
\begin{equation}
P^{(3)}_{\lambda^2}(\Omega) = \bigl[M_{12}^{(z)}\bigr]\,G^{\rm Tm}_{12}(\Omega)\,
\bar\chi^{2z}\,\sin\phi_{\rm qAFM}\,G^{\rm Fe}_{\rm qAFM}(\omega_1)\,(g_e h_x)(g_J h_z)\cdot(g_e h_x).
\label{eq:qAFME12EOM}
\end{equation}
\end{widetext}
Up to the universal $(i/\hbar^3)$ this is identical to
Eq.~\eqref{eq:A_qAFME12_full}.

\paragraph{Fe-only nonlinearity — pathway $(\mathrm{qFM},\mathrm{qAFM})$.}
Pure $N_2^{\rm Fe}$ in two cascades:
$\tilde\bfw^{(1)}_{\rm qFM}\to\tilde\bfw^{(2)}_{\rm qAFM}$ via the
$T^{(2)}_{\rm qFM,qAFM}$ vertex. Explicitly,
\begin{widetext}
\begin{equation}
\tilde\bfw^{(2)}_{\rm qAFM}(\omega_1{+}\omega_2)
= G^{\rm Fe}_{\rm qAFM}(\omega_1{+}\omega_2)\,
T^{(2)}_{\rm qFM,qAFM}\,\tilde\bfw^{(1)}_{\rm qFM}(\omega_1)\,\tilde\bfw^{(1)}_{\rm qAFM}(\omega_2),
\label{eq:qFMqAFMEOM}
\end{equation}
\end{widetext}
followed by one more drive interaction at order $\lambda^3$ to give
the detected polarization. This is the EOM realization of
Eq.~\eqref{eq:A_qFMqAFM_full} and exposes the spin-length-constraint
origin of the $T^{(2)}$ vertex (the $|\bfS|=S$ sphere curvature).

\subsection{Summary: EOM-Liouville dictionary}

The dictionary between the two derivations is the identification
\begin{widetext}
\begin{equation}
\boxed{\;\;
G^{\rm Tm}_{\alpha\beta}(\omega)\;\longleftrightarrow\;\mathcal G_{|E_\alpha\rangle\langle E_\beta|}(\omega),\qquad
N_2,N_3\;\longleftrightarrow\;\text{double-sided Feynman vertices},\;\;}
\label{eq:dict}
\end{equation}
\end{widetext}
together with the identification of $f_{abc},d_{abc}$ structure
constants with the antisymmetric/symmetric parts of the
$\hat V_i$ vertices in Eq.~\eqref{eq:Apeak}. Every $A_{\rm peak}$ in
Sec.~\ref{sec:MEcatalogue} arises from one tree diagram in
Eqs.~\eqref{eq:order1EOM}--\eqref{eq:order3EOM}; the tree topology is
fixed by the Liouville pathway. This consistency between the
operator-algebra (Liouville) and classical-EOM viewpoints is the
manuscript's strongest internal cross-check that the six peaks have
been assigned to physically allowed nonlinearities of $H_{\rm total}$.

% =====================================================================
\section{Conclusion}
\label{sec:conclusion}

We have constructed a symmetry-allowed minimal Hamiltonian for TmFeO$_3$
in the $\Gamma_2$ phase, addressing the six peaks of the corrected 2DCS
target list through explicit Liouville-pathway derivations. The
essential physics ingredients are:
\begin{enumerate}[leftmargin=*]
\item the four-sublattice Fe LLG sector with proper DM-gapped qAFM
frequency, written via an effective reduced $4\times 4$ linearized magnon matrix
$M^{\Gamma_2}$ (Sec.~\ref{sec:M4x4});
\item the three Tm singlets in the $T$-symmetric gauge with explicit
$m_J\!\to\!-m_J$ parity assignments (Sec.~\ref{sec:CEF});
\item the surviving uniform Fe--Tm exchange vertex $\bar\chi^{2z}\delta
G_z\lambda^2$ after inversion-pair cancellation of all mirror-odd channels
at $\mathbf q=0$ (Sec.~\ref{sec:TmFe});
\item the Van Vleck back-action through the $A_1^+$ Tm density
coordinates $\lambda^{1,3,8}$, which controls the population-side
nonlinear feedback entering the Fe sector (Sec.~\ref{sec:vanvleck});
\item a CPTP-compliant Lindblad dissipator (Sec.~\ref{sec:bloch});
\item a coherent emission expression with the SI impedance prefactor on
the magnetic-dipole channel (Sec.~\ref{sec:emission});
\item a Liouville-pathway derivation that gives every peak in
Tab.~\ref{tab:peaks} a definite amplitude formula in terms of microscopic
matrix elements and reduced-model nonlinear vertices (Sec.~\ref{sec:2dcs}).
\end{enumerate}

The model is restricted to the low-temperature regime $T\lesssim 20$\,K
where the linearization around $|E_1\rangle$ is justified. Extension to
the spin-reorientation window $T_{\rm SR1}<T<T_{\rm SR2}$ requires
re-linearization around a thermally mixed reference state and is left
for future work, as is the explicit treatment of magnon-polariton
dispersion in the biaxial host.

% =====================================================================
\appendix

\section{Conventions glossary}
\label{app:conventions}

A single-page reference for every prefactor used in the manuscript.
\begin{itemize}[leftmargin=*]
\item Gell-Mann normalization: $\Tr(\lambda^a\lambda^b)=2\delta^{ab}$,
$[\lambda^a,\lambda^b]=2if^{abc}\lambda^c$, $\{\lambda^a,\lambda^b\}=
\frac{4}{3}\delta^{ab}+2d^{abc}\lambda^c$.
\item Hamiltonian on Tm qutrit: $H = h_0\mathds{1}+\half\sum_a h^a\lambda^a$,
hence $h^a = \Tr(H\lambda^a)$ and the SU(3) Bloch EOM is $\dot n^a =
\sum_{bc}f^{abc}h^b n^c$.
\item Bertaut definition: $\bfX = \quart\sum_\mu\sigma^X_\mu\bfS_\mu$,
inverse $\bfS_\mu = \sum_X\sigma^X_\mu\bfX$, orthogonality
$\sum_\mu\sigma^X_\mu\sigma^Y_\mu=4\delta^{XY}$.
\item Per-site effective field on Fe: $\bfH^{\rm eff}_{\mathbf n,\mu}
= -\partial H_{\rm Fe}/\partial\bfS_{\mathbf n,\mu}$. Bertaut-projected
field: $\bfH^{X,{\rm site}}=\quart\sum_\mu\sigma^X_\mu\bfH^{\rm eff}_\mu$.
\item Anisotropy: $H_{\rm aniso}=-\sum_{\mathbf n,\mu}\sum_a K_a(S_{\mathbf n,\mu}^a)^2$,
giving $\bfH^{\rm eff,aniso}_{\mathbf n,\mu}=2\mathbf K\bfS_\mu$ and
$\bfH^{X,{\rm site}}_{\rm aniso}=2\mathbf K\bfX$ (no extra factor of 4).
\item Exchange: $H_{\rm Heis}=\sum_{\langle ij\rangle}J_{ij}\bfS_i\cdot\bfS_j$,
giving the per-site $\bfH^{\rm eff,Heis}_{\mathbf n,\mu}=-2\sum_j J_{\mu j}\bfS_j$
and the Bertaut $\bfH^{X,{\rm site}}_{\rm Heis}=-2 J^X\bfX$ with $J^X$ from
Eq.~\eqref{eq:JXdef}.
\item Larmor with energy-valued effective field: $\dot\bfS = -\hbar^{-1}\bfS\times\bfH^{\rm eff}$.
For a Zeeman term $-g_e\mu_B\,\mathbf h\cdot\mathbf S$ this becomes
$\dot\bfS = -(g_e\mu_B/\hbar)\bfS\times\mathbf h \equiv -\gamma_e\bfS\times\mathbf h$.
\item Magnetic moment of Tm$^{3+}$: $\boldsymbol\mu = -g_J\mu_B\mathbf J$ with $g_J=7/6$.
\item Stevens parameters in Tab.~\ref{tab:Stevens}: Hutchings convention
with Stevens factors absorbed.
\end{itemize}

\section{DM Bertaut prefactor derivation}
\label{app:DMderivation}

With the bond enumeration of App.~\ref{app:bonds} (4 in-plane NN bonds
per cell at $z=\half$ with $\mathbf D=(0,+D_1,+D_2)$ and 4 at $z=0$
with $\mathbf D=(0,-D_1,+D_2)$), direct symbolic projection of
$\sum_{\langle ij\rangle}\mathbf D_{ij}\cdot(\bfS_i\times\bfS_j)$ onto
the Bertaut modes (Sec.~\ref{sec:Bertautalg}) yields
\paragraph{Explicit antisymmetrization.}
For a single bond $\mu\!\to\!\nu$ in the Bertaut basis,
$(\bfS_\mu\times\bfS_\nu)_a = \sum_{X,Y}\sigma^X_\mu\sigma^Y_\nu(\bfX\times\bfY)_a
= \tfrac12\sum_{X\neq Y}(\sigma^X_\mu\sigma^Y_\nu-\sigma^Y_\mu\sigma^X_\nu)(\bfX\times\bfY)_a$.
With the four $\mathrm{Fe}_1\!\to\!\mathrm{Fe}_2$ bonds in $z=\half$ the
nonzero antisymmetric Bertaut differences are
$(\sigma^F_1\sigma^G_2-\sigma^G_1\sigma^F_2)=-2$,
$(\sigma^F_1\sigma^A_2-\sigma^A_1\sigma^F_2)=-2$,
$(\sigma^G_1\sigma^C_2-\sigma^C_1\sigma^G_2)=+2$,
$(\sigma^C_1\sigma^A_2-\sigma^A_1\sigma^C_2)=-2$;
for $\mathrm{Fe}_3\!\to\!\mathrm{Fe}_4$ in $z=0$ the corresponding
differences are $+2,-2,+2,+2$. Summing $4(\bfS_1\!\times\!\bfS_2) +
4(\bfS_3\!\times\!\bfS_4)$ with $\mathbf D=(0,+D_1,+D_2)$ and
$\mathbf D'=(0,-D_1,+D_2)$ respectively gives
\begin{widetext}
\begin{equation}
\boxed{\;H_{\rm DM}/N_c
  = -8 D_1\bigl[(\bfF\times\bfG)_y + (\bfC\times\bfA)_y\bigr]
  + 8 D_2\bigl[(\bfG\times\bfC)_z - (\bfF\times\bfA)_z\bigr].\;}
\label{eq:HDMBertautCorrect}
\end{equation}
\end{widetext}
Key structural consequences:
\begin{itemize}[leftmargin=*]
\item The $D_1$ invariant $(\bfF\times\bfG)_y = F_zG_x - F_xG_z$
reduces to $-F_xG_z$ on the $\Gamma_2$ background, recovering the
canonical Treves--Bertaut weak-FM Lifshitz energy $+8N_c D_1 F_x G_z$
and the standard canting $F_x/G_z\simeq D_1/H_E^{\rm exch}$
[Eq.~\eqref{eq:cantingangles}]. The partner $(\bfC\times\bfA)_y$
is its $\Gamma_3$-channel image.
\item The $D_2$ component generates $\bfG\times\bfC$ and $\bfF\times\bfA$
(the latter with a relative minus). $(\bfG\times\bfC)_z = G_xC_y - G_yC_x$
reduces to zero on the $\Gamma_2$ background and is the small
$\mathcal O(D_2 C_y/H_E)$ correction; it does not affect the leading
weak-FM picture.
\item The overall prefactor and sign are those of the directed-bond
convention used in Sec.~\ref{sec:Fe} and the implementation in
\verb|src/core/unitcell_builders.cpp|; reversing every bond direction
flips the overall sign but not the invariant content.
\end{itemize}

\section{Inversion-pair cancellation of mirror-odd Fe--Tm channels}
\label{sec:bondsymmetry}
\label{app:bondcancel}

Each Fe--Tm bond shell contains 8 bonds organized into 4 inversion-pair
orbits (App.~\ref{app:bonds}). Inversion at the bond midpoint sends the
Tm site $j$ to its inversion partner $\sigma^{\rm Tm}_I(j)=(14)(23)(j)$
and the Fe site to itself. Under this inversion, the spin $\bfS_i$ at
Fe$_i$ is unchanged (axial), but the Tm bilinear operator $\lambda^a$
transforms according to its $A_{1,2}^{\pm}$ class (Eq.~\eqref{eq:lambdaclass}):
\begin{equation}
\eta_a = \begin{cases} +1, & a\in\{1,2,3,8\}\;(\text{mirror-even}),\\
-1, & a\in\{4,5,6,7\}\;(\text{mirror-odd}),
\end{cases}
\label{eq:Iaction}
\end{equation}
so that $I:\;\lambda^a\to\eta_a\lambda^a$.
For a pair of inversion-related bonds, the on-site $\mathbf q=0$
projection $\sum_{\rm pair}\chi^{ab}\lambda^a S^b$ becomes
$(\eta_a + 1)\chi^{ab}\lambda^a S^b$ at $\mathbf q=0$, which vanishes
when $\eta_a=-1$. Hence:
\begin{itemize}[leftmargin=*]
\item all $\lambda^{4,5,6,7}$ on-site couplings to Fe spins cancel at
$\mathbf q=0$;
\item only $\lambda^{1,2,3,8}$ on-site couplings survive.
\end{itemize}
This is the symmetry origin of the dominant $\bar\chi^{2z}$ vertex.

\section{Bond enumeration}
\label{app:bonds}

\subsection{Fe--Fe NN, NNN, and out-of-plane bonds}

\paragraph{In-plane NN ($J_1$).} Each Fe sublattice has 4 in-plane NN
neighbors, distributed between two sublattices via the $S_1$ screw. The
in-plane NN bonds form two sets: those in the $z=\half$ plane connecting
Fe$_1\leftrightarrow$Fe$_2$, and those in $z=0$ connecting Fe$_3
\leftrightarrow$Fe$_4$. Coordination $z_1 = 4$.

\paragraph{Out-of-plane NN ($J_{1c}$).} Each Fe sublattice has 2 out-of-plane
NN neighbors connecting $z=0$ and $z=\half$ planes (Fe$_1\leftrightarrow$Fe$_4$,
Fe$_2\leftrightarrow$Fe$_3$). Coordination $z_{1c}=2$.

\paragraph{NNN ($J_2$).} Within each sublattice, 6 NNN neighbors
($\pm\hat a, \pm\hat b, \pm\hat c$). $z_2=6$.

\subsection{Fe--Tm bond shells}

Brute-force enumeration of all Fe--Tm pairs within a 4\,\AA\ cutoff gives
four distance shells:
\begin{equation}
\begin{aligned}
d_1 &= 3.054\,\text{\AA}, & d_2 &= 3.179\,\text{\AA},\\
d_3 &= 3.357\,\text{\AA}, & d_4 &= 3.711\,\text{\AA}.
\end{aligned}
\label{eq:FeTmshells}
\end{equation}
Each shell has 8 bonds per cell organized in 4 inversion pairs. The
distance spread of 21\% means the Fe--Tm exchange parameters scale
significantly between shells; we parametrize the radial dependence as
$\chi(d) \propto d^{-n}$ with $n\in\{6,8,10\}$ in the simulation.

\section{Second-order Bertaut expansion}
\label{app:Bertaut2}

The bilinear couplings entering Eq.~\eqref{eq:TqFMqAFM} are derived from
the spin-length constraint and the DM-induced cross-terms in the LLG
expansion. At second order in the linearization $\bfX = \bar\bfX + \delta
\bfX$, the cross-product structure
$\dot{\delta\bfX}^{(2)} = -\hbar^{-1}\sum_{Y\cdot Z=X}\delta\bfY\times\delta\bfH^Z$
generates the qFM--qAFM coupling. A representative contribution is
\begin{equation}
\dot{\delta F_z}^{(2)}\supset -(D_1/\hbar)\,\delta F_y\delta G_y,
\label{eq:qFMqAFMmix}
\end{equation}
which projects onto the effective reduced-mode coefficient
$T^{(2)}_{\rm qFM,qAFM}$ of Eq.~\eqref{eq:TqFMqAFM}. This appendix is
therefore a scaling diagnostic for the Fe-only quadratic vertex rather
than a closed-form formula for the reduced-mode coefficient itself.

% =====================================================================
\begin{thebibliography}{99}

\bibitem{Kimel2004}
A.~V.~Kimel, A.~Kirilyuk, A.~Tsvetkov, R.~V.~Pisarev, and Th.~Rasing,
\emph{Laser-induced ultrafast spin reorientation in the antiferromagnet TmFeO$_3$},
Nature \textbf{429}, 850 (2004).

\bibitem{Kimel2005}
A.~V.~Kimel, A.~Kirilyuk, P.~A.~Usachev, R.~V.~Pisarev, A.~M.~Balbashov, and Th.~Rasing,
\emph{Ultrafast non-thermal control of magnetization by instantaneous photomagnetic pulses},
Nature \textbf{435}, 655 (2005).

\bibitem{Baierl2016}
S.~Baierl \emph{et al.},
\emph{Nonlinear spin control by terahertz-driven anisotropy fields},
Nat.\ Photon.\ \textbf{10}, 715 (2016).

\bibitem{Schlauderer2019}
S.~Schlauderer \emph{et al.},
\emph{Temporal and spectral fingerprints of ultrafast all-coherent spin switching},
Nature \textbf{569}, 383 (2019).

\bibitem{Mashkovich2021}
E.~A.~Mashkovich \emph{et al.},
\emph{Terahertz light-driven coupling of antiferromagnetic spins to lattice},
Science \textbf{374}, 1608 (2021).

\bibitem{Zhang2016}
K.~Zhang \emph{et al.}, J.\ Magn.\ Magn.\ Mater.\ \textbf{398}, 80 (2016).

\bibitem{Grishunin2018}
K.~A.~Grishunin \emph{et al.},
ACS Photonics \textbf{5}, 1375 (2018).

\bibitem{LeendersHerrmann1971}
H.~Horner and C.~M.~Varma,
\emph{Nature of spin-reorientation transitions},
Phys.\ Rev.\ Lett.\ \textbf{20}, 845 (1968).

\bibitem{ShapiroAxeBirgeneau1974}
S.~M.~Shapiro, J.~D.~Axe, and J.~P.~Remeika,
\emph{Neutron scattering study of the spin-reorientation transition in TmFeO$_3$},
Phys.\ Rev.\ B \textbf{10}, 2014 (1974).

\bibitem{Belov1976}
K.~P.~Belov, A.~K.~Zvezdin, A.~M.~Kadomtseva, and R.~Z.~Levitin,
\emph{Spin-reorientation transitions in rare-earth magnets},
Sov.\ Phys.\ Usp.\ \textbf{19}, 574 (1976).

\bibitem{Yamaguchi1974}
T.~Yamaguchi,
\emph{Theory of spin reorientation in rare-earth orthochromites and orthoferrites},
J.\ Phys.\ Chem.\ Solids \textbf{35}, 479 (1974).

\bibitem{Skorobogatov2020}
S.~Yu.~Skorobogatov, S.~E.~Nikitin \emph{et al.},
\emph{Crystal-field excitations of Tm$^{3+}$ in TmFeO$_3$},
Phys.\ Rev.\ B \textbf{101}, 014432 (2020).

\bibitem{Hutchings1964}
M.~T.~Hutchings,
\emph{Point-charge calculations of energy levels of magnetic ions in crystalline electric fields},
Solid State Phys.\ \textbf{16}, 227 (1964).

\bibitem{NewmanNg1989}
D.~J.~Newman and B.~Ng,
\emph{The superposition model of crystal fields},
Rep.\ Prog.\ Phys.\ \textbf{52}, 699 (1989).

\bibitem{Marezio1970}
M.~Marezio, J.~P.~Remeika, and P.~D.~Dernier,
\emph{The crystal chemistry of the rare earth orthoferrites},
Acta Cryst.\ B \textbf{26}, 2008 (1970).

\bibitem{Treves1965}
D.~Treves,
\emph{Studies on orthoferrites at the Weizmann Institute},
J.\ Appl.\ Phys.\ \textbf{36}, 1033 (1965).

\bibitem{Bertaut1968}
E.~F.~Bertaut,
\emph{Representation analysis of magnetic structures},
Acta Cryst.\ A \textbf{24}, 217 (1968).

\bibitem{Herrmann1963}
G.~F.~Herrmann,
\emph{Resonance and high-frequency susceptibility in canted antiferromagnetic substances},
J.\ Phys.\ Chem.\ Solids \textbf{24}, 597 (1963).

\bibitem{Bouldi2018}
N.~Bouldi \emph{et al.}, Phys.\ Rev.\ B \textbf{98}, 174403 (2018).

\bibitem{JuddOfelt}
B.~R.~Judd, Phys.\ Rev.\ \textbf{127}, 750 (1962);
G.~S.~Ofelt, J.\ Chem.\ Phys.\ \textbf{37}, 511 (1962).

\bibitem{Wood1968}
D.~L.~Wood, L.~M.~Holmes, and J.~P.~Remeika,
Phys.\ Rev.\ \textbf{165}, 1035 (1968).

\bibitem{LuJonas2017}
J.~Lu \emph{et al.}, Phys.\ Rev.\ Lett.\ \textbf{118}, 207204 (2017).

\bibitem{WanArmitage2019}
Y.~Wan and N.~P.~Armitage, Phys.\ Rev.\ Lett.\ \textbf{122}, 257401 (2019).

\bibitem{MahmoodWan2021}
F.~Mahmood, D.~Chaudhuri, S.~Gopalakrishnan, R.~Nandkishore, and N.~P.~Armitage,
Nat.\ Phys.\ \textbf{17}, 627 (2021).

\bibitem{LiNorrisHuber2022}
Z.~Li, M.~Mootz, J.~Wang, and I.~E.~Perakis, NPJ Quantum Mater.\ \textbf{7}, 53 (2022).

\bibitem{HuberPRX2024}
R.~Huber group \emph{et al.}, Phys.\ Rev.\ X \textbf{14} (2024).

\end{thebibliography}

\appendix
\section{Executable verification status}
\label{app:script}

The current repository does \emph{not} contain a single monolithic
companion script named \verb|verify_tmfeo3.py|. Instead, the executable
checks used in this audit are distributed across several utilities, and
their coverage is only partial. The present status is:
\begin{enumerate}[leftmargin=*]
\item \verb|util/verify_tmfeo3_code.py| checks the implemented crystal
  structure, local-frame construction, Fe--Fe exchange and DM matrices,
  Fe--Tm bond tables, Tm sublattice frames, bare $h_3/h_8$ conventions,
  and the on-site/inter-site trilinear sign structure against the notes.
\item \verb|util/verify_tmfeo3_bonds.py| checks the Fe--Fe, Fe--Tm, and
  Tm--Tm bond enumerations, the inversion-pair organization of the
  Fe--Tm shells, and the DM sign convention between the $z=\half$ and
  $z=0$ planes.
\item \verb|util/verify_tmfeo3_eom.py| verifies the Gell-Mann algebra,
  Bertaut Klein-four identities, and a symbolic EOM expansion, but it is
  written in the alternate convention
  $\sigma^G=(+,-,+,-)$, $\sigma^A=(+,-,-,+)$. Its output therefore
  cannot be compared line-by-line to the corrected manuscript without a
  convention translation.
\item The Sec.~\ref{sec:CEF} qutrit spectrum, the corrected Stevens table,
  and the projected $J_a$ matrices were verified separately during this
  audit by direct diagonalization, but that check is not yet packaged as a
  repository utility.
\item The corrected DM reduction of App.~\ref{app:DMderivation} and the
  projected $\bar\chi^{2z}\sin\phi_{\rm qAFM}$ coupling of
  Eq.~\eqref{eq:n2qAFM} are therefore supported by a combination of code
  audit, bond-enumeration checks, and local symbolic reduction rather than
  by one all-encompassing script.
\item The CPTP Lindblad form of Sec.~\ref{sec:bloch} is now specified
  analytically in a complete projector-based form, but the repository does
  not yet contain a dedicated executable Choi-positivity check for that
  exact manuscript parameterization.
\end{enumerate}
The corrected manuscript should therefore be read as supported by a set of
partial executable checks plus local symbolic derivations, not by a single
monolithic verification script.

\end{document}
